% =========================================================================
% modenanalyse_2fe2s -- Manual and theoretical documentation (English)
% Version 1.1.0  (June 2026)
%
% Lukas Knauer, AG Schuenemann, RPTU Kaiserslautern-Landau
% =========================================================================
\documentclass[11pt,a4paper,twoside]{article}

% ---------- Encoding & language ------------------------------------------
\usepackage[T1]{fontenc}
\usepackage[utf8]{inputenc}
\usepackage[main=english,provide=*]{babel}
\usepackage{lmodern}
\usepackage{microtype}

% ---------- Mathematics --------------------------------------------------
\usepackage{amsmath, amssymb, amsthm}
\usepackage{mathtools}
\usepackage{bm}
\usepackage{siunitx}
\sisetup{range-phrase = --, range-units = single}

% ---------- Layout & page ------------------------------------------------
\usepackage[a4paper, left=2.5cm, right=2.5cm, top=2.5cm, bottom=3cm]{geometry}
\usepackage{setspace}
\onehalfspacing
\setlength{\parindent}{0pt}
\setlength{\parskip}{0.5em}

% ---------- Graphics & colors --------------------------------------------
\usepackage{xcolor}
\definecolor{accent}{RGB}{0,90,140}
\definecolor{codebg}{RGB}{245,245,242}
\definecolor{codeborder}{RGB}{200,200,195}
\definecolor{tomlkey}{RGB}{120,40,120}
\definecolor{tomlstr}{RGB}{0,110,0}
\definecolor{tomlcomment}{RGB}{110,110,110}
\definecolor{warnbg}{RGB}{255,248,228}
\definecolor{warnborder}{RGB}{210,160,80}
\definecolor{infobg}{RGB}{232,242,250}
\definecolor{infoborder}{RGB}{120,160,200}

% ---------- Tables -------------------------------------------------------
\usepackage{booktabs}
\usepackage{array}
\usepackage{tabularx}
\usepackage{longtable}

% ---------- Code listings ------------------------------------------------
\usepackage{listings}
\lstset{
  basicstyle=\ttfamily\small,
  backgroundcolor=\color{codebg},
  frame=single,
  rulecolor=\color{codeborder},
  breaklines=true,
  breakatwhitespace=true,
  showstringspaces=false,
  columns=flexible,
  keepspaces=true,
  xleftmargin=4pt,
  xrightmargin=4pt,
  framexleftmargin=2pt,
  framexrightmargin=2pt,
  aboveskip=0.8em,
  belowskip=0.6em,
  inputencoding=utf8,
  extendedchars=true,
  literate=%
    {ä}{{\"a}}1 {ö}{{\"o}}1 {ü}{{\"u}}1
    {Ä}{{\"A}}1 {Ö}{{\"O}}1 {Ü}{{\"U}}1
    {ß}{{\ss}}1 {µ}{{$\mu$}}1
    {→}{{$\rightarrow$}}2
    {≤}{{$\leq$}}1 {≥}{{$\geq$}}1
    {·}{{$\cdot$}}1 {–}{{--}}1 {—}{{---}}1
    {⟨}{{$\langle$}}1 {⟩}{{$\rangle$}}1
    {²}{{$^{2}$}}1 {³}{{$^{3}$}}1,
}

\lstdefinestyle{python}{
  language=Python,
  keywordstyle=\color{accent}\bfseries,
  commentstyle=\color{gray}\itshape,
  stringstyle=\color{tomlstr},
}

\lstdefinestyle{bash}{
  basicstyle=\ttfamily\small,
  keywords={},
  backgroundcolor=\color{codebg},
}

% ---------- Pseudocode ---------------------------------------------------
\usepackage{algorithm}
\usepackage{algpseudocode}
\usepackage{float}
\floatname{algorithm}{Pseudocode}

% ---------- TikZ / PGFplots ----------------------------------------------
\usepackage{tikz}
\usepackage{pgfplots}
\pgfplotsset{compat=1.18}
\usetikzlibrary{arrows.meta, calc, decorations.markings, patterns}
\renewcommand{\algorithmicrequire}{\textbf{Input:}}
\renewcommand{\algorithmicensure}{\textbf{Output:}}
\renewcommand{\algorithmicforall}{\textbf{for all}}
\renewcommand{\algorithmicif}{\textbf{if}}
\renewcommand{\algorithmicelse}{\textbf{else}}
\renewcommand{\algorithmicend}{\textbf{end}}
\renewcommand{\algorithmicfor}{\textbf{for}}
\renewcommand{\algorithmicdo}{\textbf{do}}
\renewcommand{\algorithmicwhile}{\textbf{while}}
\renewcommand{\algorithmicreturn}{\textbf{return}}

% ---------- Custom macros ------------------------------------------------
\providecommand{\angstrom}{\text{\AA}}    % Angstrom symbol for math mode

% ---------- Headings & TOC -----------------------------------------------
\usepackage{titlesec}
\titleformat{\part}[display]
  {\Huge\bfseries\color{accent}}
  {\filleft\LARGE Part \thepart}
  {0.5em}
  {\filleft}
\titleformat{\section}
  {\Large\bfseries\color{accent}}
  {\thesection}{0.7em}{}
\titleformat{\subsection}
  {\large\bfseries}
  {\thesubsection}{0.7em}{}

% ---------- Links & references -------------------------------------------
\usepackage[colorlinks=true,
            linkcolor=accent,
            urlcolor=accent,
            citecolor=accent]{hyperref}

% ---------- Boxes for notes / warnings -----------------------------------
\usepackage[most]{tcolorbox}
\newtcolorbox{warnbox}[1][]{
  colback=warnbg, colframe=warnborder, boxrule=0.5pt, arc=2pt,
  left=8pt, right=8pt, top=4pt, bottom=4pt, #1
}
\newtcolorbox{infobox}[1][]{
  colback=infobg, colframe=infoborder, boxrule=0.5pt, arc=2pt,
  left=8pt, right=8pt, top=4pt, bottom=4pt, #1
}

% ---------- Math shortcuts -----------------------------------------------
\newcommand{\evec}{\mathbf{e}}
\newcommand{\rvec}{\mathbf{r}}
\newcommand{\uvec}{\mathbf{u}}
\newcommand{\nhat}{\hat{\mathbf{n}}}
\newcommand{\xhat}{\hat{\mathbf{x}}}
\newcommand{\yhat}{\hat{\mathbf{y}}}
\newcommand{\zhat}{\hat{\mathbf{z}}}
\newcommand{\fLM}{f_{\mathrm{LM}}}
\newcommand{\ER}{E_{\mathrm{R}}}
\newcommand{\kB}{k_{\mathrm{B}}}
\newcommand{\urms}{u_{\mathrm{rms}}}
\newcommand{\Fei}{\mathrm{Fe}}
\newcommand{\OOP}{\mathrm{OOP}}
\newcommand{\INP}{\mathrm{INP}}
\newcommand{\Ag}{A_{\mathrm{g}}}
\newcommand{\Bgi}[1]{B_{\mathrm{#1g}}}
\newcommand{\Bui}[1]{B_{\mathrm{#1u}}}
\newcommand{\Aui}{A_{\mathrm{u}}}
\newcommand{\PCET}{\mathrm{PCET}}
\newcommand{\ET}{\mathrm{ET}}

% --- Formula-to-code reference macro -------------------------------------
% Used to point readers at the function in the source tree that implements
% a given formula. Format: \codref{module.function} placed AFTER the
% closing equation tag (\label), inside the equation environment but
% wrapped in \text{} so LaTeX does not interpret the underscore.
%
% The verification test tests/test_manual_code_refs.py grep's all
% \codref entries from this file and checks that every named function
% exists in the source tree. So any rename or removal in the code
% becomes a test failure -- the manual cannot silently drift out of sync.
\newcommand{\codref}[1]{%
  \quad\text{\textcolor{gray!70!black}{\footnotesize $\rightsquigarrow$~\texttt{#1}}}}

% --- Additional macros (theory sections) ---------------------------------
\newcommand{\khat}{\hat{\mathbf{k}}}
\newcommand{\bvec}{\mathbf{B}}
\newcommand{\Eiso}{e^{2}_{\mathrm{iso}}}
\newcommand{\Miso}{M_{\mathrm{iso}}}
\newcommand{\ket}[1]{\left|#1\right\rangle}
\newcommand{\bra}[1]{\left\langle#1\right|}

% ---------- Theorem environments -----------------------------------------
\newtheorem{definition}{Definition}[section]
\newtheorem{theorem}{Theorem}[section]
\theoremstyle{remark}
\newtheorem*{remark}{Remark}

% ---------- Header & footer ----------------------------------------------
\usepackage{fancyhdr}
\setlength{\headheight}{14pt}
\pagestyle{fancy}
\fancyhf{}
\fancyhead[LE,RO]{\thepage}
\fancyhead[RE]{\leftmark}
\fancyhead[LO]{\rightmark}
\renewcommand{\headrulewidth}{0.4pt}

% =========================================================================
\begin{document}

\begin{titlepage}
  \begin{center}
    \vspace*{2cm}
    {\Huge\bfseries\color{accent} modenanalyse\_2fe2s}\\[0.4cm]
    {\large Version 1.1.0}\\[2cm]

    \rule{\textwidth}{0.4pt}\\[0.3cm]
    {\LARGE\bfseries Manual and Theoretical Documentation}\\[0.3cm]
    \rule{\textwidth}{0.4pt}\\[1cm]

    {\large
      Mode-resolved bonding reorganization analysis\\
      of {[}2Fe-2S{]} iron-sulfur proteins\\
      from Gaussian-16 and ORCA-6 frequency calculations
    }\\[1.5cm]

    {\large
      Lukas Knauer\\
      AG Sch\"unemann\\
      RPTU Kaiserslautern-Landau
    }\\[1.5cm]

    \begin{tabular}{rl}
      \textbf{Cluster type:}    & {[}2Fe-2S{]}, reduced or oxidized,\\
                                & for Rieske-, ferredoxin- and mitoNEET-type ligation\\[0.3cm]
      \textbf{Input:}           & Gaussian-16 \texttt{.log} (\texttt{freq=hpmodes} preferred;\\
                                & standard block as fallback) or ORCA-6 \texttt{.hess}\\
                                & PDB file with the same cluster (optional, recommended)\\[0.3cm]
      \textbf{Output:}          & Excel reports (mode analysis, SCSD,\\
                                & reorganization energies, modulation spectra,\\
                                & UMAP cluster), report (plain text), plots (PNG)\\[0.3cm]
      \textbf{External packages:} & \texttt{numpy}, \texttt{scipy}, \texttt{pandas},\\
                                & \texttt{matplotlib}, \texttt{openpyxl},\\
                                & \texttt{scikit-learn}, \texttt{umap-learn},\\
                                & \texttt{hdbscan}, \texttt{scsdpy}\\[0.3cm]
      \textbf{References:}      & Marcus, \emph{J.\ Chem.\ Phys.} \textbf{24}, 966 (1956)\\
                                & Marcus, \emph{Angew.\ Chem.\ Int.\ Ed.\ Engl.} \textbf{32}, 1111 (1993)\\
                                & Hammes-Schiffer \& Soudackov,\\
                                & \emph{J.\ Phys.\ Chem.\ B} \textbf{112}, 14108 (2008)\\
                                & Huang \& Rhys,\\
                                & \emph{Proc.\ R.\ Soc.\ Lond.\ A} \textbf{204}, 406 (1950)\\
                                & Kingsbury \& Senge, \emph{Chem.\ Sci.} \textbf{15}, 13638 (2024)\\[0.3cm]
      \textbf{License:}         & GNU General Public License v3.0 or later\\
                                & (\texttt{GPL-3.0-or-later}, see \texttt{LICENSE})\\[0.3cm]
      \textbf{Copyright:}       & \textcopyright\ 2026 Lukas Knauer,\\
                                & AG Sch\"unemann, RPTU Kaiserslautern-Landau
    \end{tabular}

    \vfill
    {\large \today}
  \end{center}
\end{titlepage}

\tableofcontents
\clearpage

% =========================================================================
\part{Introduction and Theory}
% =========================================================================

\section{Introduction}
\label{sec:einleitung}

\texttt{modenanalyse\_2fe2s} is a scientific Python tool for analyzing
normal modes from DFT frequency calculations
(\texttt{freq=hpmodes}) of iron-sulfur proteins containing
{[}2Fe-2S{]} clusters. It combines geometric classification, rigorous
symmetry decomposition, and mode-resolved bond reorganization analysis
in a unified workflow. Application areas include mechanistic studies of
proton-coupled electron transfer (PCET), comparison of
Rieske- and ferredoxin-type systems, and the analysis of mutational
effects (e.g., HH$\to$CC).

\subsection{What \texttt{modenanalyse\_2fe2s} provides}
\label{sec:scope}

The tool focuses on the structural analysis of vibrational modes in
DFT frequency calculations of [2Fe-2S] clusters. The main per-mode
results are:

\begin{itemize}
  \item \textbf{Cluster-specific classification}: detection of the
        {[}2Fe-2S{]} core, classification of each mode according to
        out-of-plane / in-plane character, and symmetry decomposition
        into $D_\mathrm{2h}$ irreducible representations.
  \item \textbf{Ligand resolution}: PDB-based identification of
        coordinating His, Cys, and Asp/Glu residues, with per-mode
        Fe-N, Fe-S, and Fe-O stretching contributions.
  \item \textbf{Marcus-Hush reorganization energies}: per-mode
        $\Delta r_X$ and $\lambda_X$ for the bonding channels FeFe,
        FeN, FeS, NH, and HA. System-level aggregates (total
        reorganization, modulation spectra, cumulative
        $\Lambda_X(\omega)$) for band identification, processed
        externally e.g.\ in Origin.
  \item \textbf{Multi-cluster support}: detection of multiple
        {[}2Fe-2S{]} clusters in the same system (e.g., glutaredoxin
        dimer) and their parallel evaluation as a first-class feature.
  \item \textbf{Embeddings} (UMAP, HDBSCAN clustering):
        two-dimensional projection of the mode feature space for
        exploratory analysis.
  \item \textbf{ORCA support} (standalone \texttt{.hess} parser):
        mode analysis directly from ORCA frequency calculations.
\end{itemize}

NIS spectra, Fe-pDOS, the Lamb-M\"o\ss bauer factor, and
Fe-projected thermodynamics are \emph{not} part of this package.
These quantities can be computed from the same DFT frequency
calculation using specialized external tools.

\begin{infobox}
\textbf{Changelog highlights since v1.0.0.}
This release of the manual reflects the package at version v1.1.0
(see \texttt{CHANGELOG.md} for the full list). The substantive
deltas since v1.0.0 are:
\begin{itemize}\setlength{\itemsep}{0pt}
  \item \textbf{v1.0.2}: \texttt{\_build\_group\_map} fix --
        prevents silent zero rows in the \texttt{Groups\_OOP/INP/Tors}
        sheets when residue numbers overlap with non-protein entries
        (e.g.\ HOH waters). If you are running protein systems
        that contain solvent residues, you need v1.0.2 or later.
  \item \textbf{v1.0.3}: new diagnostic outputs -- the
        \texttt{Coordination} sheet documents the ligand and group
        atom mapping (\S\ref{app:wb-main-coordination}); two
        additional UMAP embeddings (SSE-UMAP and Ca-UMAP) target
        protein-domain modes and spatially localised modes
        respectively (\S\ref{app:wb-embed-ssumap},
        \S\ref{app:wb-embed-caumap}).
  \item \textbf{v1.0.4}: audit-driven uncertainty propagation fix
        (\S\ref{sec:migration-v104}). Values unchanged; the
        \texttt{s\_*}-columns for N-H stretches, SSE-element
        amplitudes, and SCSD amplitudes shift in magnitude.
        Half-open window boundaries and defense-in-depth warnings
        are also part of this release.
  \item \textbf{v1.0.5}: physics corrections to the secondary-structure
        element analysis -- the SSE-UMAP \texttt{bending\_*} features
        (which were always zero) were replaced by the real
        \texttt{lateral\_*} keys; \texttt{tilting\_angle} is now a genuine
        rigid-body tilt (rocking rotation perpendicular to the axis);
        \texttt{com\_amplitude} (per element and for the cluster core) is
        mass-weighted; and the SSE principal axis is taken from the
        C$\alpha$ backbone. The SSE sheets and their UMAP change, so affected
        systems must be re-run once; the SSE-UMAP feature set was also
        reduced to a decorrelated five-metric subset.
  \item \textbf{v1.1.0}: the abbreviation \emph{SS} (secondary structure)
        was renamed project-wide to \emph{SSE} -- sheet names, the output
        filename (\texttt{*\_analysis\_SSE.xlsx}), configuration keys and
        function names. This removes the clash with disulfide (S--S) and
        matches the conventional term \emph{secondary structure element}.
        Legacy TOML keys (\texttt{analyze\_ss}, \texttt{ss\_chain}) remain
        accepted for backward compatibility.
\end{itemize}
\end{infobox}

\subsection{Prerequisites}

\begin{itemize}
  \item Python $\geq$ 3.11 with \texttt{numpy}, \texttt{scipy},
        \texttt{openpyxl}, \texttt{matplotlib}, \texttt{pandas},
        \texttt{scikit-learn}, \texttt{umap-learn}, \texttt{hdbscan},
        and \texttt{scsdpy}. All are installed
        \emph{automatically} via the wheel.
  \item A \textbf{Gaussian-16 log file} from a
        \texttt{freq=hpmodes} calculation of the optimized cluster
        (with at least four optimized heavy atoms: two Fe and two
        bridging S) \textbf{or} an \textbf{ORCA-6
        \texttt{.hess} file}.
  \item Optional: a \textbf{PDB file} containing the same cluster, for
        ligand resolution and secondary-structure analysis.
\end{itemize}

\subsection{Conventions used in this manual}

\begin{itemize}
  \item \textbf{Axes}: wherever possible, the canonical cluster
        orientation is used:
        $\xhat$ along Fe-Fe, $\yhat$ along S-S,
        $\zhat = \nhat$ as the cluster normal. This is consistent
        with the SCSD reference geometry (see \S\ref{sec:scsd}).
  \item \textbf{Units}: cm$^{-1}$ for frequencies,
        \AA{} for lengths, AMU for masses, meV for energies,
        Kelvin for temperatures. For Boltzmann factors we use
        $1/\kB = \SI{11.6048}{K/meV}$.
  \item \textbf{Heuristic vs.\ rigorous}: whenever a quantity is
        derived from a heuristic (e.g., the
        motion-pattern scores breathing, hinge, ...), it is
        explicitly labeled as \emph{heuristic}. Rigorous
        symmetry decomposition is provided by SCSD.
\end{itemize}

\begin{infobox}
\textbf{Reading tip.} First-time users should read
\S\ref{sec:einleitung} and \S\ref{sec:konfig}, and start with the
example in \S\ref{sec:workflow}. The theoretical chapters
(\S\ref{sec:theorie}) can be consulted as needed but are not required
for routine use.
\end{infobox}

% =========================================================================
\section{Theoretical foundations}
\label{sec:theorie}

% -------------------------------------------------------------------------
\subsection{Normal modes from Gaussian: the \texttt{hpmodes} format}
\label{sec:normalmoden}

A DFT frequency calculation in Gaussian-16 yields, after a successful
geometry optimization, the second derivative of the energy with respect
to the Cartesian atomic coordinates -- the mass-weighted Hessian
$\mathbf{F}$. Diagonalizing $\mathbf{F}$ gives the eigenvectors
$\mathbf{L}$ and eigenvalues $\omega_l^2$ of the normal modes
\cite{WilsonDeciusCross1955}:

\begin{equation}
  \mathbf{F}\,\mathbf{L}_l = \omega_l^2 \, \mathbf{L}_l ,
  \qquad l = 1, \ldots, 3N-6
  \label{eq:fg-eigenproblem}
\end{equation}

with $N$ the number of atoms in the cluster (more precisely, in the
optimized model region). The Cartesian displacements per atom $i$ in
mode $l$ are obtained by mass-unweighting
\cite{WilsonDeciusCross1955}:

\begin{equation}
  \evec_{l,i} \;=\; \frac{1}{\sqrt{m_i}}\,\mathbf{L}_{l,i} ,
  \label{eq:demass}
\end{equation}

where $m_i$ is the atomic mass of atom $i$.

\subsubsection*{HP block versus standard block}

With the keyword \texttt{freq=hpmodes}, Gaussian writes the
eigenvectors in a second block (\emph{High-Precision Modes}) with
five decimal places. The \emph{standard block}, by contrast, uses
only two decimal places and is insufficient for very small
displacements (e.g., modes below 100~cm$^{-1}$ in large systems).

\texttt{modenanalyse\_2fe2s} reads both blocks. The HP eigenvectors
are preferred; only when a single HP block is missing (rare, e.g.,
in mixed-version log files) does the reader fall back to the
standard block. This fallback logic is implemented in
\texttt{logio.read\_blocks}.

\subsubsection*{Consistency check between HP and standard}
\label{sec:hp-std-check}

Both blocks output the frequencies $\omega_l$ independently.
\texttt{modenanalyse\_2fe2s} checks at run time whether the HP and
standard frequencies agree for matching mode indices
(\texttt{logio.check\_hp\_std\_frequency\_consistency}, tolerance
$\Delta\omega \leq 0{,}01$~cm$^{-1}$). A discrepancy indicates either:

\begin{itemize}
  \item an incompletely written log file (job aborted),
  \item a mismatch between the read offset and the actual block
        start (e.g., when a job was continued via \texttt{Restart}),
  \item or two log files mixed up.
\end{itemize}

Consistency errors are written as warnings to the report; the
program continues running but flags the affected modes in the
report.

\subsubsection*{Eigenvector scaling}

For all subsequent analyses, the raw Gaussian eigenvectors are
scaled by the RMS amplitude
$\urms(\omega_l, m_{\mathrm{red}}, T)$
(\texttt{core.compute\_thermal\_amplitude})
\cite{Achterhold2002,WilsonDeciusCross1955}:

\begin{equation}
  \urms(\omega, m_{\mathrm{red}}, T) \;=\;
  \sqrt{\frac{\hbar}{2\,m_{\mathrm{red}}\,\omega} \,
        \coth\!\left(\frac{\hbar\omega}{2\kB T}\right)} .
  \codref{core.compute\_thermal\_amplitude}
  \label{eq:urms}
\end{equation}

Here $m_{\mathrm{red}}$ is the reduced mass that Gaussian provides
for each mode. At $T=0$, equation~(\ref{eq:urms}) reduces to the
zero-point amplitude
$\urms = \sqrt{\hbar / (2 m_{\mathrm{red}} \omega)}$.

The eigenvectors multiplied by $\urms$ have units of \AA{} and
describe the thermally averaged (or, at $T \to 0$, purely
quantum-mechanical) displacement per mode. They are the input for
all geometric analyses from \S\ref{sec:oop-inp} onwards.

% -------------------------------------------------------------------------
\subsection{Uncertainty propagation}
\label{sec:uncertainty}

Every numerical observable that the tool reports
(\texttt{stretch}, \texttt{bend}, \texttt{kern\_d},
\texttt{amplitude\_mean}, \texttt{hn\_stretch}, the SCSD-amplitudes,
\ldots) comes with a companion uncertainty (\texttt{s\_stretch},
\texttt{s\_bend}, \texttt{s\_kern\_d}, \ldots).
These are first-order Gaussian error propagations from two underlying
input noise sources:
\begin{itemize}\setlength{\itemsep}{0pt}
  \item $\sigma_e := $ \texttt{cfg.sigma\_eigvec} (default $5\cdot 10^{-4}$):
        numerical noise per eigenvector component, in Cartesian units
        of the un-scaled Gaussian eigenvector. Controls the
        ``insignificant contribution'' threshold for displacement-based
        observables.
  \item $\sigma_c := $ \texttt{cfg.sigma\_coord} (default $10^{-3}$~\AA):
        coordinate noise in \AA{}, used for SCSD reference uncertainty.
\end{itemize}

\subsubsection*{Convention: thermally scaled sigma}
\label{sec:uncertainty-convention}

All geometric observables consume the
\emph{thermally scaled} eigenvector $\evec_{l,i}\,\urms(\omega_l)$,
which has units of \AA{}. The companion sigma is propagated at
the same scaling:
\begin{equation}
  \sigma_e^\text{therm}(l) \;=\; \sigma_e \cdot \urms(\omega_l),
  \codref{core.analyze\_his\_hn, core.analyze\_fe\_ligand}
  \label{eq:sigma-therm}
\end{equation}
so that the dimensionless \emph{significance ratio}
$|\text{observable}|\,/\,\sigma$ is invariant under the choice of
$\urms$ for a given physical quantity. This convention is essential
for the value/sigma ratios used in the colour coding of the Excel
sheets.

\paragraph{Worked examples.} For a typical low-frequency mode with
$\urms \approx 0.05$~\AA{}:
\begin{itemize}\setlength{\itemsep}{0pt}
  \item Fe-X stretch ($i = $ ligand, $j = $ Fe):
        $|\evec_i - \evec_j|$ is a difference of two independent
        Gaussian-noise components, so
        $\sigma(\text{stretch}) = \sigma_e^\text{therm}\sqrt{2}
        \approx 3.5\cdot 10^{-5}$~\AA.
  \item Bend INP/OOP: same factor $\sqrt{2}$ (Pythagorean
        decomposition does not change the noise propagation along
        any single axis), so $\sigma_\text{bend\_inp}
        = \sigma_\text{bend\_oop} = \sigma_e^\text{therm}\sqrt{2}$.
  \item N-H stretch (\texttt{analyze\_his\_hn}, \texttt{include\_hn\_vibration})
        follows the same convention:
        $\sigma_\text{hn} = \sigma_e^\text{therm}\sqrt{2}$.
  \item SSE-element amplitude:
        $\sigma_\text{amp} = \sigma_e^\text{therm}\sqrt{N_\text{atoms}}$,
        the $\sqrt{N}$ growth coming from the averaging over
        $N_\text{atoms}$ in the secondary-structure element.
  \item SCSD amplitudes: distinct ``reference'' and ``distortion''
        sigmas. The reference is the static geometry uncertainty
        $\sigma_c\sqrt{2}$, the distortion combines reference and
        thermal noise:
        $\sigma_\text{SCSD,dist}
         = \sqrt{2}\sqrt{\sigma_c^2 + (\sigma_e \cdot \urms)^2}$.
\end{itemize}

\paragraph{What changed in v1.0.4.} Three of these sigma
propagations were previously implemented inconsistently
(\S\ref{sec:migration-v104}): N-H stretch lacked the thermal
scaling, SSE-element used a hard-coded constant in place of
\texttt{cfg.sigma\_eigvec}, and SCSD used hard-coded reference
constants in place of \texttt{cfg.sigma\_coord} and
\texttt{cfg.sigma\_eigvec}. With v1.0.4, all sigmas honour
Equation~\ref{eq:sigma-therm} and the
\texttt{cfg.sigma\_eigvec}/\texttt{sigma\_coord} settings
from the TOML.

% -------------------------------------------------------------------------
\subsection{Migration notes for v1.0.4}
\label{sec:migration-v104}

Version 1.0.4 is a bug-fix release that affects the
\emph{reported uncertainties} of several observables but leaves
the underlying \emph{values} unchanged. If you have v1.0.2 or v1.0.3
results in hand (e.g.\ for a paper or thesis), please re-read this
section before drawing significance conclusions from the older
output.

\paragraph{What did not change.} All geometric and energetic
observables -- \texttt{stretch}, \texttt{bend}, \texttt{kern\_d},
\texttt{lambda\_X}, SCSD amplitudes, \texttt{hn\_stretch},
SSE-element \texttt{amplitude\_mean}, embedding coordinates --
produce \emph{numerically identical} values across v1.0.2, v1.0.3
and v1.0.4 for the same input log. No physical interpretation of
existing values is invalidated.

\paragraph{What changed.}

\begin{itemize}\setlength{\itemsep}{2pt}
  \item \texttt{s\_hn\_stretch}: now scales with $\urms$
        (Eq.~\ref{eq:sigma-therm}). For typical low-frequency modes
        ($\urms \sim 0.05$~\AA) this makes \texttt{s\_hn\_stretch}
        about \textbf{20\,$\times$ smaller} than in v1.0.2/v1.0.3.
        Consequence: many N-H stretches that appeared
        ``below noise'' under the old convention are now correctly
        identified as significant PCET signals.
  \item \texttt{s\_amplitude\_mean} / \texttt{s\_axial\_amplitude} /
        \texttt{s\_tilting\_angle} \ldots\ in
        \texttt{analyze\_sse\_element}: pre-v1.0.4 used a hard-coded
        $1\cdot 10^{-4}$ in place of \texttt{cfg.sigma\_eigvec}.
        With the default $\sigma_e = 5\cdot 10^{-4}$, the new sigmas
        are about \textbf{5\,$\times$ larger}. Re-classification:
        SSE-element entries that were marked ``significant'' under
        the old code may now read as ``marginal''.
  \item \texttt{s\_SCSD\_*}: pre-v1.0.4 used hard-coded
        $5\cdot 10^{-7}$ / $5\cdot 10^{-6}$ in place of
        \texttt{cfg.sigma\_coord} and \texttt{cfg.sigma\_eigvec}.
        New sigmas are about \textbf{1000\,$\times$ larger}; the
        previous values were physically implausible.
  \item Window-edge handling: multi-window mode
        (\texttt{freq\_windows}) used \emph{closed-closed} intervals
        $[lo, hi]$, so a mode with frequency exactly equal to a
        window boundary appeared in \emph{two} adjacent windows. As
        of v1.0.4 the boundaries are \emph{half-open} $[lo, hi)$,
        with the last window remaining $[lo, hi]$. Practical effect:
        none for typical float-frequency runs, but
        $B$-factor and $\Lambda(\omega_\text{cut})$ are now
        guaranteed double-count free.
  \item Defense-in-depth warnings: a number of internal lookup
        failures (Fe ligand $c2l$ miss, His-H $c2l$ miss, missing
        PCET acceptor index, all-zero output rows in the Excel
        sheets, ambiguous PDB-Gaussian atom matches) used to be
        silent fall-throughs. With v1.0.4 each of these emits a
        \texttt{UserWarning} the first time it occurs in a run
        (deduplicated per session). If you see new warnings in a
        previously-quiet log, check the affected ligand or atom
        mapping before trusting the corresponding Excel rows.
  \item Interpolation boundary anchor (FUND~12, user-reported):
        when \texttt{freq\_max} is set and the next mode above it
        lies farther than \texttt{interp\_context\_cm1} (default
        $5$~cm$^{-1}$), pre-v1.0.4 loaded no context mode and the
        interpolated pDOS dropped to zero just below
        \texttt{freq\_max}. With v1.0.4, exactly \emph{one} additional
        mode -- the closest above \texttt{freq\_max} -- is loaded
        as a fallback anchor, so $\texttt{np.interp}$ always has a
        right-edge anchor regardless of mode density. The same
        fallback applies symmetrically at \texttt{freq\_min}. The
        run-log records when the fallback fires.
  \item Synthetic zero anchor when the DFT spectrum ends inside the
        analysis range (FUND~13): an extension of the above. If
        \emph{no} real mode exists above \texttt{freq\_max} at all
        (the DFT spectrum genuinely terminates), v1.0.4 inserts a
        \emph{synthetic null-mode} at
        $\texttt{freq\_max} + \texttt{interp\_context\_cm1}$ with all
        observable fields set to zero. This gives the interpolated
        pDOS a smooth decay across the buffer zone instead of a step
        at the upper edge. The synthetic mode is flagged with
        \texttt{mode\_type = "synthetic\_zero"} and
        \texttt{precision = "synthetic"} so it is unambiguous in any
        per-mode listing. By construction it contributes zero to
        $B$-factors, $\Lambda_X^\text{total}$, and the modulation
        spectra, making the fix conservative.
\end{itemize}

\paragraph{Post-release audit patches (May 2026).} After v1.0.4
was released, an audit on a real-world [2Fe-2S]-Cys$_2$His$_2$
batch run (analysing both the protonated and the deprotonated
states of a Rieske-type system) identified four additional
warning-and-reporting issues. None of them affects numerical
output. Full descriptions are in the supplement
(FUND 14--17); briefly:

\begin{itemize}\setlength{\itemsep}{2pt}
  \item FUND 14: The
        ``\texttt{His\_HN NICHT erkannt}'' warning fired twice per
        deprotonated His instead of once. Fixed by hoisting the
        warning out of the PDB-vs-Gaussian-fallback loop.
  \item FUND 15: For systems with deprotonated His ligands
        (e.g.\ Rieske-type Cys$_2$His$_2$ in the deprotonated
        state, or mitoNEET-type Cys$_3$His after deprotonation),
        the PCET-disabled message stated ``\texttt{no His ligands
        at cluster}'' even though \texttt{Fe\_N\_His} sheets and
        \texttt{Lambda\_FeN} confirmed His ligands were present.
        The decision text now distinguishes ``no His'' from
        ``His present but all deprotonated''.
  \item FUND 16: The
        ``\texttt{interp\_boundary\_mode='context' but no context
        modes available}'' warning fired even for full-spectrum
        runs (\texttt{freq\_min}/\texttt{freq\_max} both
        \texttt{None}), where no context-loading is attempted. The
        warning is now suppressed unless the user has explicitly
        set a window.
  \item FUND 17: REPORT.txt showed \texttt{freq\_filter: - - -
        cm-1} when only \texttt{freq\_windows} was set, hiding
        the active filter from the run audit trail. The report now
        emits a separate \texttt{freq\_windows: [...]} line.
\end{itemize}

\paragraph{Re-running existing data.} The recommended path: rerun
your existing TOML configurations under v1.0.4 and compare the
\texttt{s\_*} columns of \texttt{Fe\_S\_Cys.xlsx},
\texttt{His\_HN.xlsx}, \texttt{SCSD.xlsx}, and the SSE-Element
sheets. Values are identical; only sigmas and the
significance-colouring change.

% -------------------------------------------------------------------------
\subsection{OOP/INP classification}
\label{sec:oop-inp}

A central question for every cluster mode is: \emph{how much of the
vibration is perpendicular to the cluster plane (out-of-plane, OOP),
and how much lies in the plane (in-plane, INP)?}
Pure rocking motions of the four cluster atoms (umbrella, hinge) are
$\OOP$-dominated, whereas stretching modes (Fe-Fe, S-S, breathing)
live in the plane.

\subsubsection*{Definition of the cluster normal}

Let $\rvec_i$ ($i \in \{\Fei_1, \Fei_2, \mathrm{S}_1, \mathrm{S}_2\}$)
be the equilibrium positions of the four cluster atoms. The normal
direction of the plane is the direction that minimizes the sum of
squared perpendicular components. Concretely, it is obtained from
the singular-value decomposition of the centered cluster coordinates
$\tilde{\rvec}_i = \rvec_i - \bar{\rvec}$ \cite{Kabsch1976}:

\begin{equation}
  \tilde{\mathbf{R}} \;=\; \mathbf{U}\,\boldsymbol{\Sigma}\,\mathbf{V}^\top ,
  \qquad
  \nhat \;=\; \mathbf{V}_{:,3}\,.
  \codref{geometry.cluster\_normal}
  \label{eq:cluster-normal-svd}
\end{equation}

This is implemented as \texttt{geometry.cluster\_normal}.
By convention, $\nhat$ points in the $+\zhat$ direction of the
Gaussian frame, so that a uniform sign convention applies across
all modes.

\subsubsection*{Mode component perpendicular to the plane}

For mode $l$ and an atom $i$, the OOP component is the projection
of the displacement onto $\nhat$ (standard decomposition, cf.\
\cite{Gee2021}):

\begin{equation}
  \evec_{l,i}^{\OOP} \;=\; (\evec_{l,i} \cdot \nhat)\, \nhat ,
  \qquad
  \evec_{l,i}^{\INP} \;=\; \evec_{l,i} - \evec_{l,i}^{\OOP} .
  \label{eq:oop-projektion}
\end{equation}

The \emph{OOP fraction} of the mode on an atom group $G$ is the
ratio of the squared norms \cite{Gee2021}:

\begin{equation}
  \alpha_{\OOP}(G; l) \;=\;
  \frac{\sum_{i \in G} \lVert\evec_{l,i}^{\OOP}\rVert^2}
       {\sum_{i \in G} \lVert\evec_{l,i}\rVert^2} \,\cdot\, 100\,\%
  \,.
  \codref{core.classify\_oop\_inp}
  \label{eq:oop-anteil}
\end{equation}

\begin{infobox}
\textbf{Note on the physical interpretation.}
Gaussian's \texttt{hpmodes} block writes
\emph{Cartesian-displacement} eigenvectors~$\evec_{l,i}$ (already
mass-unweighted via Eq.~\ref{eq:demass}), normalised so that
$\sum_i \lVert\evec_{l,i}\rVert^2 \approx 1$ rather than the
mass-weighted norm $\sum_i m_i \lVert\evec_{l,i}\rVert^2$.
Equation~\ref{eq:oop-anteil} therefore measures the
\emph{geometric Cartesian out-of-plane fraction}: it tells you
what share of the squared Cartesian displacement amplitudes points
out of the cluster plane. It does \emph{not} equal the kinetic-energy
fraction of the mode flowing out of the plane, which would require
the mass-weighted form
$\sum_i m_i \lVert\evec_{l,i}^{\OOP}\rVert^2 / \sum_i m_i \lVert\evec_{l,i}\rVert^2$.
For a [2Fe-2S] cluster the two quantities are numerically similar
(only Fe and S contribute, with comparable masses), but the
distinction matters when comparing $\alpha_{\OOP}$ values across
chemically different systems.
\end{infobox}

\subsubsection*{Binary and fine classification}

Classically, $\alpha_{\OOP}(G_{\text{cluster}}; l)$ is subjected to
a threshold test:

\begin{equation}
  \text{mode\_type}(l) \;=\;
  \begin{cases}
    \OOP, & \alpha_{\OOP} \geq \theta_{\text{OOP}} \\
    \INP, & \alpha_{\OOP} \leq 100\,\% - \theta_{\text{OOP}} \\
    \text{mixed}, & \text{otherwise}
  \end{cases}
\end{equation}

with default $\theta_{\text{OOP}} = 60\,\%$ (configurable via
\texttt{cfg.mode\_type\_threshold}).

In addition, there is a \emph{fine classification} with three
thresholds $(\theta_1, \theta_2, \theta_3)$
(\texttt{core.classify\_oop\_inp}), default $(60, 75, 90)$:

\begin{center}
\begin{tabularx}{\textwidth}{l X}
\toprule
\textbf{Level} & \textbf{Definition} \\
\midrule
\textbf{Pure OOP}        & $\alpha_{\OOP} \geq \theta_3 = 90\,\%$ \\
\textbf{Strong OOP}      & $\theta_2 \leq \alpha_{\OOP} < \theta_3$, i.e.\ $75-90\,\%$ \\
\textbf{Majority OOP}    & $\theta_1 \leq \alpha_{\OOP} < \theta_2$, i.e.\ $60-75\,\%$ \\
\textbf{Mixed}           & $|{\alpha_{\OOP} - 50\,\%}| < \theta_1 - 50\,\%$ \\
\textbf{Majority INP}    & $100\,\% - \theta_2 < \alpha_{\OOP} \leq 100\,\% - \theta_1$, i.e.\ $25-40\,\%$ \\
\textbf{Strong INP}      & $100\,\% - \theta_3 < \alpha_{\OOP} \leq 100\,\% - \theta_2$, i.e.\ $10-25\,\%$ \\
\textbf{Pure INP}        & $\alpha_{\OOP} \leq 100\,\% - \theta_3 = 10\,\%$ \\
\bottomrule
\end{tabularx}
\end{center}

The fine classification provides a more precise statement about the
character of a mode in the Excel sheet -- particularly in the
OOP-hinge regime, where a binary 60\,\% threshold is either too
permissive (modes between 60 and 75\,\% look very different from
each other) or returns a fuzzy answer.

\begin{infobox}
\textbf{Which threshold for which question?} For comparison with
NRVS spectra in polarization geometry, the binary 60\,\% threshold
provides sufficient discrimination. For manuscript tables of
mode properties (e.g., $D_\mathrm{2h}$ symmetry candidates), the
fine 7-level classification is more informative; with it, ``pure''
OOP / INP modes can be reliably identified.
\end{infobox}

\subsubsection*{Worked example: computing OOP\% for a cluster breathing mode}

To illustrate, we compute $\alpha_{\OOP}$ explicitly for a
hypothetical cluster breathing mode in a Cys$_4$ model cluster.

\paragraph{Step 1 -- geometry and cluster normal.}
The four cluster atoms sit, in idealized form, at (units of $\angstrom$):
\begin{align*}
  \mathbf{r}_{\text{Fe}_1} &= (+1.50,\, 0.00,\, 0.00) \\
  \mathbf{r}_{\text{Fe}_2} &= (-1.50,\, 0.00,\, 0.00) \\
  \mathbf{r}_{\text{S}_1}  &= ( 0.00,\, +1.10,\, 0.00) \\
  \mathbf{r}_{\text{S}_2}  &= ( 0.00,\, -1.10,\, 0.00)
\end{align*}
Centroid $\bar{\mathbf{r}} = \mathbf{0}$. The centered coordinates
all lie in the $xy$ plane; SVD yields
$\hat{\mathbf{n}} = (0, 0, 1)$ (cluster normal parallel to the $z$ axis).

\paragraph{Step 2 -- mode eigenvectors (thermally scaled, $T=5$\,K).}
A pure \textbf{breathing mode} stretches and contracts both Fe-S
diagonals synchronously in the plane. A typical thermally scaled
eigenvector looks as follows (after $\mathbf{e} \to \mathbf{e}\cdot u_\text{rms}$):
\begin{align*}
  \mathbf{e}_{\text{Fe}_1} &= (+0.012,\, +0.009,\, 0.000) \\
  \mathbf{e}_{\text{Fe}_2} &= (-0.012,\, -0.009,\, 0.000) \\
  \mathbf{e}_{\text{S}_1}  &= ( 0.000,\, -0.018,\, 0.000) \\
  \mathbf{e}_{\text{S}_2}  &= ( 0.000,\, +0.018,\, 0.000)
\end{align*}

\paragraph{Step 3 -- extract the OOP components.}
For each atom $i$:
$\mathbf{e}_i^\text{OOP} = (\mathbf{e}_i \cdot \hat{\mathbf{n}})\,
\hat{\mathbf{n}}$. Since all eigenvectors have $z$ component $0$,
$\mathbf{e}_i^\text{OOP} = \mathbf{0}$ for all $i$.

\paragraph{Step 4 -- evaluate $\alpha_{\OOP}$.}
\begin{align*}
  \sum_{i \in G} \|\mathbf{e}_i^\text{OOP}\|^2 &= 0 \\
  \sum_{i \in G} \|\mathbf{e}_i\|^2 &= (0.012^2 + 0.009^2)\cdot 2 + (0.018^2)\cdot 2 \\
   &= 2\cdot 0.000225 + 2\cdot 0.000324 = 0.001098
\end{align*}
\begin{equation*}
  \alpha_{\OOP} = \frac{0}{0.001098}\cdot 100\,\% = 0\,\%
\end{equation*}
Classification: \textbf{Pure INP} ($\alpha_{\OOP} \leq 10\,\%$). The
cluster breathing motion is an in-plane mode -- as expected.

\paragraph{Step 5 -- comparison with the hinge mode.}
A cluster hinge mode (the ``book-opening'' of the two Fe-S bonds
about an S-S axis) has, by contrast, the form:
\begin{align*}
  \mathbf{e}_{\text{Fe}_1} &= (0.000,\, 0.000,\, +0.020) \\
  \mathbf{e}_{\text{Fe}_2} &= (0.000,\, 0.000,\, +0.020) \\
  \mathbf{e}_{\text{S}_1}  &= (0.000,\, 0.000,\, -0.020) \\
  \mathbf{e}_{\text{S}_2}  &= (0.000,\, 0.000,\, -0.020)
\end{align*}
Here $\mathbf{e}_i^\text{OOP} = \mathbf{e}_i$ for all $i$,
so $\alpha_{\OOP} = 100\,\%$ -- classified as \textbf{Pure OOP}.

\paragraph{Step 6 -- what happens in a real DFT run.}
In a real Cys$_4$ or Cys$_2$His$_2$ cluster, pure breathing or
hinge modes are rare -- components typically mix, and
$\alpha_{\OOP}$ lies between 20 and 80\,\%. Classification
according to the default thresholds $(60, 75, 90)$ produces a
seven-level distribution, which for a typical [2Fe-2S] run
might look something like:
\begin{itemize}\setlength{\itemsep}{0pt}
  \item Pure INP: $\sim 30$\,\% of cluster modes
  \item Majority/Strong INP: $\sim 35$\,\%
  \item Mixed: $\sim 20$\,\%
  \item Majority/Strong OOP: $\sim 12$\,\%
  \item Pure OOP: $\sim 3$\,\%
\end{itemize}
(These numbers vary considerably between systems -- in rigid
cluster geometries, INP modes dominate; in flexible ligand spheres,
mixed modes are more common.)

% -------------------------------------------------------------------------
\subsection{Fe-ligand bond decomposition: stretch, bend INP, bend OOP}
\label{sec:fe-lig-decomp}

The OOP/INP decomposition of \S\ref{sec:oop-inp} applies to the
\emph{cluster core} (Fe$_1$, Fe$_2$, S$_1$, S$_2$). A second
decomposition is performed for every Fe-ligand bond
(Fe-S$_\text{Cys}$, Fe-N$_\text{His}$, Fe-O$_\text{Asp/Glu}$): the
relative displacement of the two bond partners is split into a
\emph{stretch} component (along the bond axis) and a
\emph{bend} component (perpendicular to it), the bend then being
split further into INP and OOP relative to the cluster plane.

Let $\evec_\text{Fe}$ and $\evec_\text{lig}$ be the thermally
scaled displacements of the Fe and the ligand-donor atom, and let
\begin{equation}
  \hat{\mathbf{b}} = \frac{\rvec_\text{lig} - \rvec_\text{Fe}}
                            {\lVert\rvec_\text{lig} - \rvec_\text{Fe}\rVert}
\end{equation}
be the unit vector pointing from Fe to the ligand donor in the
equilibrium geometry. Define the relative displacement
\begin{equation}
  \Delta\evec = \evec_\text{lig} - \evec_\text{Fe}.
  \label{eq:fe-lig-delta}
\end{equation}

\paragraph{Decomposition along and perpendicular to the bond.}
The signed stretch component is the projection on $\hat{\mathbf{b}}$;
the bend vector is what remains:
\begin{equation}
  \Delta r_\text{stretch} = \Delta\evec \cdot \hat{\mathbf{b}},
  \qquad
  \Delta\evec_\perp = \Delta\evec
                       - \Delta r_\text{stretch}\hat{\mathbf{b}}.
  \label{eq:fe-lig-perp}
\end{equation}

\paragraph{Splitting the bend into INP and OOP parts.}
The cluster normal $\hat{\mathbf{n}}$ (\S\ref{sec:oop-inp}) further
decomposes $\Delta\evec_\perp$ into a component along $\hat{\mathbf{n}}$
(out of the cluster plane) and a remainder (in the cluster plane):
\begin{align}
  \Delta r_\text{bend\_oop} &= \Delta\evec_\perp \cdot \hat{\mathbf{n}},
  \qquad\text{(signed)} \\
  \Delta r_\text{bend\_inp} &= \lVert\Delta\evec_\perp
       - \Delta r_\text{bend\_oop}\hat{\mathbf{n}}\rVert.
  \codref{core.analyze\_fe\_ligand}
  \label{eq:fe-lig-bend-split}
\end{align}
The reported \texttt{bend}, \texttt{bend\_inp}, and \texttt{bend\_oop}
columns in \texttt{Fe\_S\_Cys.xlsx} / \texttt{Fe\_N\_His.xlsx} are
the \emph{absolute values} (lengths) of these three quantities, all
in \AA{}. By Pythagoras,
\begin{equation}
  \lVert\Delta\evec_\perp\rVert^2
    = (\Delta r_\text{bend\_inp})^2 + (\Delta r_\text{bend\_oop})^2,
\end{equation}
which is the natural energy decomposition of the bend mode.

\paragraph{Bend INP / OOP percentages.}
For comparing modes with very different overall amplitudes, the
absolute lengths are converted to \emph{squared-amplitude
fractions} relative to the total bend:
\begin{equation}
  \text{bend\_inp\_pct} = 100 \cdot
     \frac{(\Delta r_\text{bend\_inp})^2}
          {(\Delta r_\text{bend\_inp})^2 + (\Delta r_\text{bend\_oop})^2},
\end{equation}
and analogously \texttt{bend\_oop\_pct}. The two add to 100\,\%
exactly. For very small overall bends (below the noise level
$2\sigma_e^\text{therm}\sqrt{2}$, \S\ref{sec:uncertainty}), the
percentages are set to NaN to avoid amplifying numerical noise.

\paragraph{Sigma propagation.}
All three quantities (stretch, bend\_inp, bend\_oop) are differences
of two independent Gaussian-noise eigenvector components and
therefore carry the standard $\sqrt{2}\,\sigma_e^\text{therm}$ sigma:
\begin{equation}
  \sigma_\text{stretch} = \sigma_\text{bend\_inp}
  = \sigma_\text{bend\_oop}
  = \sqrt{2}\,\sigma_e \cdot \urms(\omega_l).
\end{equation}
The sigmas are written into the
\texttt{Fe\_S\_Cys.xlsx} / \texttt{Fe\_N\_His.xlsx} files under
\texttt{s\_stretch}, \texttt{s\_bend\_inp}, \texttt{s\_bend\_oop}
and are used for the significance colouring of the bend cells.

\paragraph{Physical interpretation.}
For PCET applications, the INP/OOP split of the Fe-N(His) bend is
particularly informative:
\begin{itemize}\setlength{\itemsep}{0pt}
  \item \emph{bend\_inp dominant} ($>$70\,\%): the His ring rocks
        in the cluster plane -- a motion that strongly modulates
        the Fe-Fe distance and the iron-sulfur core, classic
        ET-active.
  \item \emph{bend\_oop dominant} ($>$70\,\%): the His ring tilts
        out of plane, which modulates the N-H bond direction --
        coupling channel for PT to a buried acceptor.
  \item \emph{mixed} (30--70\,\%): the mode contributes to both
        ET and PT pathways simultaneously, the typical PCET
        fingerprint.
\end{itemize}

% -------------------------------------------------------------------------
\subsection{SCSD after Kingsbury \& Senge}
\label{sec:scsd}

The OOP/INP decomposition of the previous section is a
\emph{geometric heuristic} -- it knows only the cluster plane, not
the symmetry. A rigorous symmetry decomposition is provided by the
\emph{Symmetry-Coordinate Structural Decomposition} (SCSD) of
Kingsbury \& Senge \cite{Kingsbury2024}.

\subsubsection*{Idea}

A {[}2Fe-2S{]} cluster with ideal rhombic geometry belongs to the
point group $D_\mathrm{2h}$. The eight possible displacement
patterns of the four atoms (1 translation $\times$ 3 axes + 1
rotation $\times$ 3 axes + 2 stretches) distribute over the eight
$D_\mathrm{2h}$ irreducible representations:
$\Ag, \Bgi{1}, \Bgi{2}, \Bgi{3}, \Aui, \Bui{1}, \Bui{2}, \Bui{3}$.

A real cluster geometry has finite distortions, however -- the
cluster is never ideally rhombic. SCSD projects the actual atomic
coordinates onto the eight irreps of the point group and yields a
contribution amplitude per irrep. Applied to a normal mode, SCSD
answers: \emph{what symmetry does this displacement actually have?}

\subsubsection*{$D_\mathrm{2h}$ character table and symmetry operations}

The point group $D_\mathrm{2h}$ has eight symmetry operations
$\{E, C_2(z), C_2(y), C_2(x), i, \sigma_h, \sigma_v(xz), \sigma_v(yz)\}$
and correspondingly eight one-dimensional irreps. The character table:

\begin{center}
\begin{tabular}{lccccccccl}
  \toprule
  \textbf{$D_{2h}$} & $E$ & $C_2(z)$ & $C_2(y)$ & $C_2(x)$ &
                     $i$ & $\sigma_h$ & $\sigma_v(xz)$ & $\sigma_v(yz)$ &
                     \textbf{Linear basis} \\
  \midrule
  $A_g$    & 1 & 1 & 1 & 1 & 1 & 1 & 1 & 1 & $x^2, y^2, z^2$ \\
  $B_{1g}$ & 1 & 1 & $-$1 & $-$1 & 1 & 1 & $-$1 & $-$1 & $xy$ \\
  $B_{2g}$ & 1 & $-$1 & 1 & $-$1 & 1 & $-$1 & 1 & $-$1 & $xz$ \\
  $B_{3g}$ & 1 & $-$1 & $-$1 & 1 & 1 & $-$1 & $-$1 & 1 & $yz$ \\
  $A_u$    & 1 & 1 & 1 & 1 & $-$1 & $-$1 & $-$1 & $-$1 & --- \\
  $B_{1u}$ & 1 & 1 & $-$1 & $-$1 & $-$1 & $-$1 & 1 & 1 & $z$ \\
  $B_{2u}$ & 1 & $-$1 & 1 & $-$1 & $-$1 & 1 & $-$1 & 1 & $y$ \\
  $B_{3u}$ & 1 & $-$1 & $-$1 & 1 & $-$1 & 1 & 1 & $-$1 & $x$ \\
  \bottomrule
\end{tabular}
\end{center}

In the [2Fe-2S] context, the axis convention (Kingsbury 2024) is:
$\hat{\mathbf{x}}$ along Fe-Fe, $\hat{\mathbf{y}}$ along S-S
(bridging-sulfur to bridging-sulfur),
$\hat{\mathbf{z}} = \hat{\mathbf{n}}$ (cluster normal).

\subsubsection*{Meaning of the individual irreps for cluster motion}

Directly from the character table one can read off which motions
correspond to which irrep:

\begin{center}
\begin{tabular}{lp{8.5cm}}
  \toprule
  \textbf{Irrep} & \textbf{Cluster motion in [2Fe-2S] convention} \\
  \midrule
  $A_g$    & Symmetric stretches along the principal axes.
             Examples: Fe-Fe breathing, S-S stretch,
             cluster-breathing mixture. The classical
             ``reorganization mode'' for ET. \\
  $B_{1g}$ & In-plane bending with simultaneous distortion of
             both diagonals. \\
  $B_{2g}$ & In-plane motion with unequal Fe-S distances on the
             two cluster sides. \\
  $B_{3g}$ & In-plane motion with rotation of the rhombus center
             relative to the atoms. \\
  $A_u$    & ``Central inversion mode'' -- rare in rigid
             clusters; can appear in strongly distorted
             geometries. \\
  $B_{1u}$ & Pure OOP displacement of all cluster atoms in the
             $\hat{\mathbf{z}}$ direction. Hinge mode. \\
  $B_{2u}$ & OOP displacement with twist about the $\hat{\mathbf{x}}$ axis. \\
  $B_{3u}$ & OOP displacement with twist about the $\hat{\mathbf{y}}$ axis. \\
  \bottomrule
\end{tabular}
\end{center}

The $g$ irreps are \emph{gerade} (in-plane, parity-preserving),
the $u$ irreps \emph{ungerade} (typically OOP-active).

\subsubsection*{Canonical reference and sign convention}

The projection requires an unambiguous axis choice. \texttt{scsdpy}
uses the Kingsbury canonical convention:

\begin{enumerate}
  \item $\zhat$ is the cluster normal.
  \item The cluster's principal axis (longest interatomic distance)
        is placed along $\xhat$.
  \item Signs are chosen so that the cluster atoms fall in a specific
        order onto positive quadrants.
\end{enumerate}

In \texttt{modenanalyse\_2fe2s}, the reference geometry is built once
per run and stored (\texttt{core.\_get\_scsd\_model}). It is
independent of the orientation in which Gaussian optimized the
cluster.

\subsubsection*{Application to normal modes}

For mode $l$, the cluster (four atoms) is displaced by $\evec_{l,i}\,\urms$;
the difference between this displaced geometry and the reference
geometry yields a 12-dimensional displacement vector. The SCSD
projection produces eight amplitudes \cite{Kingsbury2024}:

\begin{equation}
  \mathbf{a}_l \;=\;
  \mathbf{P}_{D_{2\mathrm{h}}}\bigl(\evec_l^{\text{cluster}} \,\urms\bigr) ,
  \qquad
  \mathbf{a}_l \in \mathbb{R}^8 .
  \codref{core.compute\_scsd\_for\_mode\_full}
  \label{eq:scsd-projection}
\end{equation}

The squared sum $\sum_\Gamma a_{l,\Gamma}^2$ equals the unweighted
sum over the squared atomic displacements -- the SCSD decomposition
is an \emph{orthogonal} basis transformation.

\subsubsection*{Dominant irrep -- determined rigorously}
\label{sec:scsd-rigoros}

The rigorous function \texttt{core.extract\_dominant\_scsd\_irreps}
identifies the dominant irrep as follows. It:

\begin{enumerate}
  \item sums the squared amplitudes per irrep,
  \item normalizes to the total sum,
  \item returns the \emph{primary} and \emph{secondary} irrep,
        together with their percentage contributions.
\end{enumerate}

This classification can disagree with the OOP/INP heuristic: a mode
can, for example, be ``majority OOP'' (heuristic) and at the same
time ``primarily $\Bgi{1}$'' (SCSD) -- $\Bgi{1}$ is a specific
OOP symmetry, which the heuristic does not distinguish.

\begin{infobox}
\textbf{The \texttt{SCSD} sheet in the Excel output} lists, per mode,
the eight amplitudes plus the two dominant irreps with their
percentage contributions. Modes with consistent ($>$80\,\%) primary
irrep are those that should be cited in the symmetry table of a
manuscript.
\end{infobox}

\subsubsection*{When SCSD is disabled}

The SCSD analysis requires the external package \texttt{scsdpy}. If
it is not installed, the SCSD sheet is skipped and a note is written
to the report. All other analyses (OOP/INP heuristic, kernel
classification, reorganization energies) are unaffected.

% -------------------------------------------------------------------------
\subsection{Heuristic kernel classification}
\label{sec:kern-klassifikation}

In addition to the rigorous SCSD decomposition, there is a
\emph{heuristic kernel classification} that assigns each mode to a
motion type based on its geometric displacement pattern
(\texttt{core.classify\_kernel\_mode\_from\_evg}). It is robust even
without \texttt{scsdpy} and provides an intuitive categorization. The
heuristic is explicitly labeled as such and should not be confused
with the rigorous SCSD answer.

\subsubsection*{Motion scores}

The heuristic classifier computes eight motion scores per mode:

\begin{center}
\begin{tabularx}{\textwidth}{l X}
\toprule
\textbf{Score} & \textbf{Definition} \\
\midrule
\textbf{Translation}    & Center-of-mass displacement of the cluster
                          (should be small for an intramolecular mode) \\
\textbf{Umbrella}       & Mean OOP displacement of all four cluster atoms
                          (all moving collectively out of the plane) \\
\textbf{Hinge-Fe}       & OOP difference Fe$_1$ minus Fe$_2$
                          (Fe atoms rocking against each other) \\
\textbf{Hinge-S}        & OOP difference S$_1$ minus S$_2$
                          (S atoms rocking against each other) \\
\textbf{OOP-Twist}      & Combination of both hinge modes (twisting) \\
\textbf{Fe stretch}     & In-plane component along the Fe-Fe axis \\
\textbf{S stretch}      & In-plane component along the S-S axis \\
\textbf{Breathing}      & Breathing motion: all four atoms simultaneously
                          moving away from / toward the centroid \\
\textbf{Rhombus shear}  & Antisymmetric distortion of the rhombus
                          (Fe-S bond angles) \\
\textbf{Rotation-ip}    & In-plane rotation of the cluster about its normal \\
\bottomrule
\end{tabularx}
\end{center}

Each score is normalized to the total displacement of the four
cluster atoms and lies in $[0, 1]$.

\subsubsection*{Classification}

The two highest scores are returned as the \emph{primary} and
\emph{secondary} motion patterns. In the Excel sheet
\texttt{Kernel\_Scores}, all eight scores plus the two top labels
are listed. This heuristic provides a pre-classification that works
\textit{without} a $D_\mathrm{2h}$ symmetry assumption (also in
distorted geometries).

\subsubsection*{Heuristic vs.\ SCSD: when do they disagree?}

In ideally rhombic clusters the two agree:

\begin{center}
\begin{tabularx}{\textwidth}{l l X}
\toprule
\textbf{Heuristic} & \textbf{SCSD irrep} & \textbf{Note} \\
\midrule
Umbrella     & $\Bgi{2}$ or $\Bgi{3}$ & depending on the axis convention \\
Hinge-Fe     & $\Aui$           & antisymmetric about the cluster center \\
Breathing    & $\Ag$            & fully symmetric \\
Fe stretch   & $\Bgi{1}$        & antisymmetric along Fe-Fe \\
Rotation-ip  & $\Bgi{1}$ (in-plane) & depending on the axis choice \\
\bottomrule
\end{tabularx}
\end{center}

In the distorted clusters typical of proteins, motions often split
across two heuristic scores (e.g., a mode with Hinge-Fe \emph{and}
rhombus shear of comparable magnitude). Here the SCSD answer is
more conclusive: it assigns a unique symmetry that does not depend
on the axis choice.

\begin{warnbox}
\textbf{Caveat -- do not confuse.} The heuristic labels (umbrella,
hinge, breathing, ...) are \emph{motion descriptors}, not symmetry
labels. A ``Hinge-Fe'' mode can carry symmetry $\Aui$, $\Bui{1}$,
or $\Bui{2}$ depending on the axis choice. For statements about
mode symmetry in a manuscript, always use the SCSD answer, not the
heuristic.
\end{warnbox}

% -------------------------------------------------------------------------
\subsection{Embeddings: UMAP and HDBSCAN}
\label{sec:embeddings}

For a [2Fe--2S] system with 5\,000--19\,000 modes, displaying all
mode properties simultaneously on two-dimensional plots is not
feasible. \texttt{modenanalyse\_2fe2s} uses UMAP (Uniform Manifold
Approximation and Projection) for dimensionality reduction,
optionally followed by HDBSCAN clustering.

\subsubsection*{Choice of method: why UMAP?}

Compared to alternatives (PCA, t-SNE), UMAP has, after test runs on
protonated and deprotonated Cys$_2$His$_2$ model systems as well
as an HHCC mutant, proven to be the only method that consistently
yields interpretable results for DFT mode data.

\begin{itemize}
  \item \textbf{PCA} \cite{Pearson1901,Hotelling1933,Jolliffe2002}
        is linear and fast, but cannot resolve nonlinear manifolds.
        For DFT modes, the data points typically lie on curved
        submanifolds in feature space (frequency-OOP-INP-score
        relations are nonlinear); PCA shows only a shadow
        projection. PC1 in our data typically explains only
        25--35\,\% of the variance -- too little for a stand-alone
        visualization.

  \item \textbf{t-SNE} \cite{vanDerMaaten2008,vanDerMaaten2014}
        preserves local neighborhoods, but global distances are
        meaningless: an apparently large gap between two cluster
        regions in a t-SNE plot does not necessarily correspond to
        a large distance in feature space \cite{vanDerMaaten2008}.
        Furthermore, t-SNE is stochastic and the results depend on
        the random seed. For cluster identification, t-SNE often
        creates the appearance of clusters even when the
        underlying data are continuously distributed.

  \item \textbf{UMAP} \cite{McInnes2018,Becht2019} preserves both
        local neighborhood and global topology in a controllable
        way and is significantly faster than t-SNE at comparable
        quality \cite{Becht2019}. The \texttt{n\_neighbors}
        parameter allows explicit control of the locality--globality
        balance, which is essential for mode analysis (small
        n\_neighbors emphasize finer symmetry classes, large
        n\_neighbors the global frequency structure).
\end{itemize}

A PCA-style user perspective (variance explained per axis) is no
longer required here, because the meaning of each feature is made
explicit upstream of the embedding by the curated selection (see
section \emph{Feature matrix}). The other sheets in the Excel
output (\texttt{Mode\_analysis}, \texttt{Reorganization\_energy},
\texttt{Reorg\_total}) provide the quantitative information that
PCA would have contributed in addition.

\subsubsection*{UMAP: mathematical foundation}

UMAP constructs a low-dimensional embedding in two steps, based on
the theory of topological manifolds and fuzzy simplicial sets
\cite{McInnes2018}.

\paragraph{Step 1: high-dimensional neighborhood graph.}

In the original feature space $\mathbb{R}^p$ (with $p = 31$
features), a local neighborhood graph is constructed for each data
point $x_i$. Per point, the $k$ nearest neighbors are determined
(parameter \texttt{n\_neighbors}). Distances to these neighbors
are converted via an adaptive metric $d_i(x_i, x_j)$ into a
probability $\mu_{ij} \in [0,1]$:
\begin{equation}
  \mu_{ij} = \exp\!\left(\frac{-\max(0,\,
              d(x_i, x_j) - \rho_i)}{\sigma_i}\right),
  \label{eq:umap-mu}
\end{equation}
where $\rho_i$ is the distance to the nearest neighbor of $x_i$
(this guarantees that every point is connected to its immediate
neighbor with probability 1) and $\sigma_i$ is chosen such that
$\sum_j \mu_{ij} = \log_2(k)$ (local-scale adaptation).

The asymmetry of $\mu_{ij}$ is combined into a symmetric adjacency
matrix $w_{ij}$ via a \emph{fuzzy union} \cite{McInnes2018}:
\begin{equation}
  w_{ij} = \mu_{ij} + \mu_{ji} - \mu_{ij}\,\mu_{ji}.
  \label{eq:umap-w}
\end{equation}

\paragraph{Step 2: low-dimensional embedding.}

In the 2D target space (\texttt{n\_components=2}), positions $y_i$
are sought that produce a similar neighborhood graph $v_{ij}$.
The probability in the target space is
\begin{equation}
  v_{ij} = \bigl(1 + a\,\|y_i - y_j\|^{2b}\bigr)^{-1},
\end{equation}
with constants $a \approx 1.577$ and $b \approx 0.895$ from the
default configuration (\texttt{min\_dist=0.1}). The loss function
is the fuzzy cross-entropy:
\begin{equation}
  \mathcal{L} = \sum_{ij} \Bigl[w_{ij}\log\frac{w_{ij}}{v_{ij}}
              + (1 - w_{ij})\log\frac{1 - w_{ij}}{1 - v_{ij}}\Bigr],
  \label{eq:umap-loss}
\end{equation}
minimized by stochastic gradient descent.

\paragraph{Preserved properties.}

In contrast to t-SNE, UMAP preserves not only local but also global
structure \cite{Becht2019}. Distances between cluster regions in
the UMAP projection correlate with distances in feature space --
with the caveat that large distances are slightly \emph{compressed}
(UMAP is not distance-preserving, but topology-preserving).

\subsubsection*{Parameters}

\begin{itemize}
  \item \texttt{n\_neighbors} (auto heuristic:
        $\max(5, \min(15, N/100))$): controls the local-vs.-global
        balance. Small values (5--10) preserve the finest local
        structure; large values (20--50) emphasize global
        topology. For a typical [2Fe-2S] protein with
        $N \approx 5\,000$--$10\,000$ modes, the auto heuristic
        yields $\texttt{n\_neighbors} = 15$, which matches the
        paper standard \cite{McInnes2018}.

  \item \texttt{min\_dist} = 0.1 (default): minimum distance
        between points in the embedding. Smaller values produce
        more compact clusters.

  \item \texttt{n\_components} = 2: 2D visualization.

  \item \texttt{random\_state} = 42 (deterministically reproducible).
\end{itemize}

Configurable via \texttt{cfg.umap\_n\_neighbors} (manual override).

\subsubsection*{Feature matrix}
\label{sec:embedding-features}

Per mode, 31 numerical features are extracted. Two originally
planned \texttt{Cys *\_OOP} features are by geometric construction
always zero (Cys-S$_\gamma$ lies in the cluster plane and has no
out-of-plane component) and were excluded. The remaining 31
features are:

\begin{enumerate}\setlength{\itemsep}{0pt}\setlength{\parsep}{0pt}
  \item[\textbf{Cluster-ring fractions}] (3 features, see
        section~\ref{sec:scsd}):
        \begin{itemize}\setlength{\itemsep}{0pt}
          \item \texttt{lig\_OOP\%} -- fraction of the mode motion
                in the ligand sphere that is OOP relative to the
                cluster plane
          \item \texttt{kern\_OOP\%} -- the same for the cluster
                core (Fe + cluster-S)
          \item \texttt{kern\_|d|} -- absolute cluster
                displacement in \AA{} (norm)
        \end{itemize}

  \item[\textbf{Kernel scores / $D_{2h}$ irreps}] (10 features,
        see section~\ref{sec:scsd}):
        \begin{itemize}\setlength{\itemsep}{0pt}
          \item \texttt{Translation}, \texttt{Umbrella},
                \texttt{Hinge-Fe}, \texttt{Hinge-S},
                \texttt{OOP-Twist}, \texttt{Fe-stretch},
                \texttt{S-stretch}, \texttt{Breathing},
                \texttt{Rhombus-shear}, \texttt{Rotation-ip}
        \end{itemize}
        These are projections of the mode onto the irreducible
        representations of the ideal $D_{2h}$ cluster symmetry.

  \item[\textbf{Ligand OOP fractions}] (2 features):
        \begin{itemize}\setlength{\itemsep}{0pt}
          \item \texttt{His 255\_OOP}, \texttt{His 259\_OOP}
        \end{itemize}
        Distribution of the ligand OOP motion across the two His
        residues. The values satisfy
        \texttt{His 255\_OOP + His 259\_OOP = 1}; they are
        perfectly anticorrelated. UMAP, however, does not weight
        anticorrelated features more strongly, and treating both
        ligands symmetrically is conceptually cleaner.
        \emph{Not} included: \texttt{Cys 207\_OOP},
        \texttt{Cys 216\_OOP} -- see above.

  \item[\textbf{Ligand INP fractions}] (4 features):
        \begin{itemize}\setlength{\itemsep}{0pt}
          \item \texttt{Cys 207\_INP}, \texttt{Cys 216\_INP},
                \texttt{His 255\_INP}, \texttt{His 259\_INP}
        \end{itemize}
        The same for in-plane fractions. Here too perfect
        pairwise anticorrelation, but all four are kept.

  \item[\textbf{Per-ligand atom motions}] (12 features, see
        sheets \texttt{Fe\_S\_Cys}, \texttt{Fe\_N\_His}):
        \begin{itemize}\setlength{\itemsep}{0pt}
          \item Per Cys (207, 216): \texttt{stretch},
                \texttt{bend\_inp}, \texttt{bend\_oop}
          \item Per His (255, 259): \texttt{stretch},
                \texttt{bend\_inp}, \texttt{bend\_oop}
        \end{itemize}
        Stretch and bend of the Fe-S / Fe-N bond vector with
        the bend component split into INP and OOP parts relative
        to the cluster plane.
\end{enumerate}

Before UMAP, the feature matrix is standardized
(\texttt{sklearn.preprocessing.StandardScaler}; per feature mean
0, variance 1). This ensures that no individual variable dominates
solely because of its absolute value range.

\paragraph{What is deliberately excluded from the matrix.}

\begin{itemize}
  \item \emph{Frequency} ($\omega_l$): not used as a feature,
        since otherwise the UMAP map would trivially sort by
        frequency. Frequency is kept as a metadata column and can
        be used for coloring the UMAP plot.

  \item \emph{Reorganization lambda} ($\lambda_X$): not introduced
        as a feature, because the underlying geometric modulations
        (\texttt{dr\_FeFe}, \texttt{dr\_NH}, \texttt{dr\_FeN})
        are already implicitly captured in the per-ligand
        stretch/bend features. Including $\lambda_X$ in addition
        would over-emphasize reorganization-relevant modes.

  \item \emph{Secondary structure} (HELIX/SHEET fractions): handled
        in a \emph{separate} secondary-structure UMAP (sheet
        \texttt{SSE\_UMAP\_Cluster}). The mode UMAP and SSE UMAP
        are conceptually different analyses.
\end{itemize}

\subsubsection*{HDBSCAN: density-based cluster detection}

HDBSCAN \cite{Campello2013,McInnes2017} is a hierarchical
density-based clustering algorithm. Unlike k-means, it does not
require a predetermined number of clusters; it identifies clusters
as connected regions of high density and labels isolated points as
noise (label $-1$).

\paragraph{min\_cluster\_size.} The only important parameter:
the minimum cluster size. \texttt{modenanalyse\_2fe2s} uses an
auto heuristic
\[
  \texttt{mcs} = \max\bigl(5, \min(30, \sqrt{N})\bigr),
\]
For typical [2Fe-2S] protein systems
($N \approx 5\,000$--$10\,000$), this gives $\texttt{mcs} = 30$.

\begin{infobox}
\textbf{If HDBSCAN finds many clusters or much noise --
what does that mean?}

UMAP on a typical Cys$_2$His$_2$ system with 31 features and
\texttt{n\_neighbors=15} typically yields 5--8 clusters with no
noise. If HDBSCAN finds significantly more clusters (e.g., 80--100),
this is usually due to too small an \texttt{mcs} value or a
fragmented UMAP topology. If HDBSCAN finds 0 clusters (everything
noise), the mode distribution in feature space is too smooth --
which is expected behavior for DFT mode data without sharp
transitions between mode classes \cite{Campello2013}.

\textbf{Consequence.} The UMAP visualization itself is informative
even without HDBSCAN clusters -- one can see mode densities and
transitions through coloring by frequency, mode type, or kernel
score. HDBSCAN provides additional information, but is not a
required component of the analysis.
\end{infobox}

\paragraph{Required.} \texttt{hdbscan} and \texttt{umap-learn} must
be installed. Without these packages the UMAP sheets are skipped.

\subsubsection*{Practical recommendation}

\begin{enumerate}
  \item \textbf{Inspect the UMAP plot}: the sheet
        \texttt{Projections} contains the UMAP coordinates for all
        modes. Color in Origin/matplotlib by frequency, mode type,
        or kernel-primary score.

  \item \textbf{Examine the cluster sheets}: the sheet
        \texttt{UMAP\_Cluster} shows HDBSCAN labels and feature
        values per cluster. \texttt{UMAP\_Cluster\_Profile}
        shows Z-score profiles and the top modes per cluster.

  \item \textbf{Comparison between structures} (e.g., prot vs.\
        deprot, or vs.\ a mutant variant) is performed externally:
        two UMAP plots for the two structures side by side, colored
        on the same frequency scale. Within a single run no direct
        comparison is performed -- one system per run.
\end{enumerate}


% =========================================================================
% Theory chapter for the mode-resolved reorganization pipeline
% =========================================================================


\subsection{Marcus-Hush theory of reorganization}
\label{sec:marcus-hush}

The reorganization chapter \S\ref{sec:reorg} computes per-mode
contributions to the reorganization energy. Before going there, we
need classical Marcus-Hush theory as a \emph{benchmark anchor} -- so
that it is clear what our mode-resolved variant has in common with
historical Marcus theory and where the two differ.

\subsubsection*{Classical picture: two diabatic energy surfaces}

Marcus' theory of electron transfer \cite{Marcus1956,Marcus1993}
describes a reaction between a reactant and a product state via two
\emph{diabatic} energy surfaces $U_R(\mathbf{Q})$ (reactant) and
$U_P(\mathbf{Q})$ (product). The reaction coordinates
$\mathbf{Q} = (Q_1, Q_2, \ldots, Q_N)$ are the vibrational modes of
the reactive complex -- \emph{all} of them, not selected ones.

In the harmonic approximation, both surfaces are paraboloids:
\begin{equation}
  U_R(\mathbf{Q}) = U_R^0 + \frac{1}{2}\sum_i \mu_i\,\omega_i^2\,
                       (Q_i - Q_i^R)^2
\end{equation}
\begin{equation}
  U_P(\mathbf{Q}) = U_P^0 + \frac{1}{2}\sum_i \mu_i\,\omega_i^2\,
                       (Q_i - Q_i^P)^2
\end{equation}
where $\mu_i$ is the reduced mass, $\omega_i$ the eigenfrequency,
and $Q_i^R, Q_i^P$ are the equilibrium displacements of mode $i$ in
the two states.

Marcus' central assumption is: \textbf{the frequencies $\omega_i$
are equal in both states}, only the displacements
$Q_i^R \neq Q_i^P$. This is the ``equal-harmonic'' approximation; it
is valid when reactant and product are electronically similar
(e.g., two oxidation states of the same complex).

\subsubsection*{Reorganization energy as the central quantity}

The \textbf{reorganization energy} $\lambda$ is defined as the
energy the system pays if it remains in the reactant state but is
displaced to the product geometry:
\begin{equation}
  \lambda = U_R(\mathbf{Q}^P) - U_R(\mathbf{Q}^R)
          = \frac{1}{2}\sum_i \mu_i\,\omega_i^2\,(\Delta Q_i)^2
  \quad\text{with}\quad \Delta Q_i = Q_i^P - Q_i^R.
  \label{eq:lambda-marcus-classical}
\end{equation}
This is the famous Marcus-Hush standard form. It is an
\emph{equilibrium property} of the system -- not a mode property,
not an activation energy, not a kinetic quantity. The thermal
population of individual modes plays no role in
\eqref{eq:lambda-marcus-classical}.

\subsubsection*{Activation energy and rate}

From $\lambda$ and the free reaction energy
$\Delta G^0 = U_P^0 - U_R^0$, a short geometric argument leads to
the Marcus activation energy:
\begin{equation}
  \Delta G^\ddagger = \frac{(\lambda + \Delta G^0)^2}{4\lambda}
\end{equation}
and consequently to the Marcus rate:
\begin{equation}
  k_\text{ET} = \frac{2\pi}{\hbar}|H_{RP}|^2 \cdot
                \frac{1}{\sqrt{4\pi\lambda k_BT}} \cdot
                \exp\!\left(-\frac{\Delta G^\ddagger}{k_BT}\right)
  \label{eq:marcus-rate}
\end{equation}
where $H_{RP}$ is the electronic coupling element between the
reactant and product diabats. Three characteristic features:

\begin{itemize}\setlength{\itemsep}{0pt}
  \item For small $|\Delta G^0|$: $\Delta G^\ddagger \approx
        \lambda/4$. \emph{The reorganization energy is the measure
        of the reaction barrier.}
  \item For $\Delta G^0 = -\lambda$: barrierless, $k$ is maximal
        ($\Delta G^\ddagger = 0$).
  \item For $|\Delta G^0| > \lambda$: the \emph{Marcus inverted
        regime}, where the rate decreases again because the
        reaction becomes too exothermic.
\end{itemize}

In iron-sulfur proteins, all three regimes are relevant:
ferredoxin-type ET steps are often in the Marcus normal regime;
PCET steps in mitoNEET and Rieske can touch the inverted region.

\subsubsection*{What this theory provides -- and what it does not}

\paragraph{What Marcus-Hush delivers:}
\begin{itemize}\setlength{\itemsep}{0pt}
  \item A complete analytical description of the ET/PCET
        activation energy, provided the harmonic approximation
        holds.
  \item A directly measurable prediction of the rate
        $k_\text{ET}$ (eq.~\ref{eq:marcus-rate}), if $\lambda$,
        $\Delta G^0$, and $|H_{RP}|^2$ are known.
  \item A justification for why $\lambda$ is system-specific (no
        universal value: it depends on the solvation environment,
        the cluster geometry, and the ligands).
\end{itemize}

\paragraph{What Marcus-Hush does not directly provide:}
\begin{itemize}\setlength{\itemsep}{0pt}
  \item \emph{Which} modes contribute substantially to $\lambda$.
        The sum in \eqref{eq:lambda-marcus-classical} is global;
        promoter modes (see below) had to be teased out in the
        1990s by Hammes-Schiffer and others.
  \item \emph{Direct access from DFT frequency calculations}.
        Classical computation of $\lambda$ requires two separate
        frequency calculations (reactant + product) plus a robust
        mode mapping between them -- a non-trivial effort.
  \item \emph{Reorganization energy as a spectroscopic
        observable}. $\lambda$ from
        \eqref{eq:lambda-marcus-classical} is an activation energy
        for a concrete reaction, not a band frequency or a line
        profile.
\end{itemize}

These three gaps motivate extending the theory in two directions:

\begin{enumerate}
  \item \textbf{Mode-resolved view} (Hammes-Schiffer, Cukier,
        Stuchebrukhov): decomposes $\lambda$ into contributions of
        individual modes and identifies promoter modes with
        particular significance for the rate. The multi-mode view
        turns \eqref{eq:lambda-marcus-classical} into
        frequency-resolved information.
  \item \textbf{Spectroscopically accessible variant}: replaces the
        static displacements $\Delta Q_i$ with thermal bond
        modulations $\Delta r_X(i)$. This turns $\lambda$ into a
        quantity that follows directly from a single frequency
        calculation -- and that assigns bands in NRVS spectra.
\end{enumerate}

The second extension is exactly what \texttt{modenanalyse\_2fe2s}
computes in \S\ref{sec:reorg}. What we call $\lambda_X$ there is
\emph{not} the classical Marcus-Hush $\lambda$ from
\eqref{eq:lambda-marcus-classical}, but a related yet distinct
quantity -- the thermal bond reorganization energy per mode and
per bond. The connection to Huang-Rhys theory makes clear why this
quantity is a well-defined spectroscopic observable (next section).

\subsubsection*{Hammes-Schiffer extension: multi-mode PCET}

Hammes-Schiffer and Soudackov \cite{HammesSchifferSoudackov2008}
showed that PCET in complex systems such as proteins \emph{cannot}
be described by a single promoter mode. Instead, the reaction is
coupled to many modes simultaneously; the spectrum of contributions
$\lambda_i = \tfrac{1}{2}\mu_i\omega_i^2(\Delta Q_i)^2$ is broadly
distributed.

Standard Marcus theory can capture this multi-mode reality through
an \emph{effective} reorganization energy
$\lambda_\text{eff} = \sum_i \lambda_i$, but the mode-resolved
distribution is then lost. For NRVS band identification, that
distribution is precisely what matters:

\begin{itemize}\setlength{\itemsep}{0pt}
  \item Which frequency ranges contribute most strongly to the
        reorganization?
  \item Which bonds are active in those ranges?
  \item Which mutations should change which band structure?
\end{itemize}

The multi-mode view is also why global ``top-mode classification''
approaches (PCET score, $L_\text{PCET}$, filter-variant) were
conceptually unsatisfying: they searched for \emph{one} dominant
mode, while reality is a broad distribution. The mode-resolved
reorganization pipeline makes that distribution explicit.

\subsubsection*{What does $\Delta Q$ mean vs.\ $\Delta r_X$?}

A central subtlety in preparation for the next chapter: the
classical Marcus displacement $\Delta Q_i$ (mass-weighted mode
coordinate) is \emph{not} the same as a bond modulation
$\Delta r_X$ (Cartesian distance change in $\angstrom$).

\begin{itemize}\setlength{\itemsep}{0pt}
  \item $\Delta Q$: displacement of the \emph{eigenmode coordinate}
        between two equilibrium geometries. Has units of
        $\sqrt{\text{amu}}\cdot\angstrom$. Appears in
        $\eqref{eq:lambda-marcus-classical}$.
  \item $\Delta r_X(i)$: \emph{bond modulation} along the $X$ bond,
        induced by mode $i$ with thermally scaled displacements.
        Has units of $\angstrom$. Computed in
        \texttt{reorganization.compute\_mode\_modulations}.
\end{itemize}

Both have the same \emph{form} in the standard equation
$\lambda = \tfrac{1}{2}\mu\omega^2(\cdot)^2$, but distinct
\emph{physical meanings}. The bridge between the two is
Huang-Rhys theory -- see the next section.

\subsection{Huang-Rhys theory and our $\lambda_X$}
\label{sec:huang-rhys}

In \S\ref{sec:reorg} we compute
$\lambda_X(i) = \tfrac{1}{2}\mu_i\omega_i^2(\Delta r_X(i))^2$ with
\emph{thermally scaled} bond modulations. This section shows that
this is exactly the \textbf{Huang-Rhys factor} (in a particular
convention), and that this quantity has well-defined spectroscopic
meaning.

\subsubsection*{What Huang and Rhys originally described}

Huang and Rhys \cite{HuangRhys1950} computed optical line profiles
of color centers in solids. Their question: when an electronic
transition shifts the equilibrium position of a harmonic
oscillator, what does the resulting vibronic spectrum look like?

Answer: the spectrum is a Poisson distribution
\begin{equation}
  P(n) = \frac{S^n e^{-S}}{n!}
\end{equation}
over phonon number $n$, parametrized by the \textbf{Huang-Rhys
factor}
\begin{equation}
  S = \frac{\mu\,\omega\,(\Delta q)^2}{2\hbar},
  \label{eq:huang-rhys}
\end{equation}
where $\Delta q$ is the transition-induced equilibrium shift of a
mode with reduced mass $\mu$ and frequency $\omega$.

\paragraph{Meaning of $S$:}
\begin{itemize}\setlength{\itemsep}{0pt}
  \item $S \ll 1$: the transition is approximately vertical;
        the line spectrum is dominated by the 0--0 band.
  \item $S \approx 1$: strong mode coupling; vibronic band with
        several phonon excitations.
  \item $S \gg 1$: Gaussian-shaped band structure (Stokes shift
        equal to $S\hbar\omega$).
\end{itemize}

In modern language: $S$ measures mode coupling
\emph{quantitatively}, in units of the zero-point amplitude
$u_\text{ZP} = \sqrt{\hbar/(2m\omega)}$:
\begin{equation}
  S = \frac{(\Delta q)^2}{2 u_\text{ZP}^2}.
  \label{eq:huang-rhys-units}
\end{equation}

\subsubsection*{Connection to Marcus-Hush}

Compare \eqref{eq:huang-rhys} with the Marcus reorganization
energy of a single mode (from
\eqref{eq:lambda-marcus-classical} with $N=1$):
\begin{equation}
  \lambda^\text{Marcus}_\text{single}
  = \frac{1}{2}\mu\,\omega^2\,(\Delta q)^2
  = S \cdot \hbar\omega.
  \label{eq:lambda-vs-S}
\end{equation}

Thus: \textbf{$\lambda^\text{Marcus} = S \cdot \hbar\omega$}.
The two quantities are interconvertible.

For a multi-mode distribution:
$\lambda^\text{Marcus}_\text{total} = \sum_i S_i \cdot \hbar\omega_i$.

\subsubsection*{What is $\Delta q$ in our case?}

The crux: in classical Marcus theory, $\Delta q$ is a
\emph{static displacement} of the mode coordinate between two
electronic states. \texttt{modenanalyse\_2fe2s}, however, has
\emph{only one} electronic state per run -- the single DFT
frequency calculation provides only one equilibrium geometry and
its eigenmodes.

Solution: we replace the static displacement with a \emph{thermal
displacement}:
\begin{equation}
  (\Delta q)^2 \to \langle q^2 \rangle = \frac{\hbar}{2\mu\omega}
  \coth\!\left(\frac{\hbar\omega}{2k_BT}\right)
  = u_\text{rms}^2.
  \label{eq:thermal-replacement}
\end{equation}

With the convention
$\Delta r_X(i) = u_\text{rms}(i) \cdot \alpha_X(i)$, where
$\alpha_X(i)$ is the mode contribution along bond $X$, our
$\lambda_X(i)$ becomes:
\begin{equation}
  \lambda_X(i) = \frac{1}{2}\mu_i\,\omega_i^2\,
                  u_\text{rms}^2(i)\,\alpha_X^2(i).
\end{equation}

Substituting \eqref{eq:thermal-replacement} and using
$\hbar\omega/(k_BT) \to \infty$ at $T \to 0$ (so $\coth \to 1$),
one obtains:
\begin{equation}
  \lambda_X(i) \xrightarrow{T \to 0}
  \frac{1}{2}\mu_i\,\omega_i^2\,\frac{\hbar}{2\mu_i\omega_i}\,\alpha_X^2(i)
  = \frac{\hbar\omega_i}{4}\,\alpha_X^2(i).
  \label{eq:lambda-zero-T}
\end{equation}

\paragraph{Connection to the Huang-Rhys factor.}
From \eqref{eq:lambda-vs-S} with
$(\Delta q)^2 = u_\text{rms}^2 \cdot \alpha_X^2$:
\begin{equation}
  S_X(i) = \frac{\mu_i\omega_i u_\text{rms}^2(i)\,\alpha_X^2(i)}{2\hbar}
  \xrightarrow{T \to 0} \frac{\alpha_X^2(i)}{4}.
\end{equation}
In the $T \to 0$ limit, the Huang-Rhys factor per bond is therefore
simply \emph{one quarter of the squared mode contribution}. For a
pure stretching mode of the bond ($\alpha_X = 1$), $S_X = 1/4$ --
that is, \emph{weak coupling} in the Huang-Rhys sense, with a
typical line spectrum.

\subsubsection*{Magnitude}

Let us put in concrete numbers. For a typical Fe-S stretching
mode at $\omega = 350$~cm$^{-1}$ and $T = 5$\,K:
\begin{itemize}\setlength{\itemsep}{0pt}
  \item Zero-point energy: $\hbar\omega/4 = 87.5$~cm$^{-1}$.
  \item For a pure stretching mode ($\alpha_X = 1$):
        $\lambda_X(T \to 0) = 87.5$~cm$^{-1}$.
  \item Under thermal population at 5\,K
        ($\ll \hbar\omega/k_B = 503$\,K): $\coth \approx 1$, hence
        essentially the zero-point value.
  \item At room temperature (300\,K, $\coth \approx 1.04$):
        $\lambda_X \approx 91$~cm$^{-1}$, only marginal increase.
\end{itemize}

The same mode at $\omega = 100$~cm$^{-1}$:
\begin{itemize}\setlength{\itemsep}{0pt}
  \item $\hbar\omega/4 = 25$~cm$^{-1}$.
  \item At 5\,K: $\coth \approx 1.0$, $\lambda_X \approx 25$~cm$^{-1}$.
  \item At 300\,K: $\coth \approx 4.2$, $\lambda_X \approx 105$~cm$^{-1}$.
\end{itemize}
Low-frequency modes are therefore strongly thermally populated at
room temperature -- their $\lambda_X$ contributions are amplified
relative to the zero-point value by a factor
$\coth(\hbar\omega/2k_BT)$.

\paragraph{Consequence for low-temperature NRVS.}
At low temperature (5--40\,K, $\hbar\omega \gg k_BT$ for typical
NRVS frequencies), $u_\text{rms}^2 \approx \hbar/(2\mu\omega)$,
hence $T$-independent. We therefore have
\begin{equation}
  \lambda_X(i) \approx \frac{\hbar\omega_i}{4}\alpha_X^2(i)
  \quad\text{at low-temperature NRVS}.
\end{equation}
This form is \emph{purely quantum-mechanical} (zero-point energy),
independent of the precise sample temperature. This is exactly what
NRVS measures: the quantum-mechanical mode activity, not a thermal
Boltzmann average.

\subsubsection*{What our $\lambda_X$ is -- final definition}

With these preliminaries we can now formulate a precise definition:

\begin{tcolorbox}[colback=accent!5,colframe=accent,
  title={What $\lambda_X$ is in \texttt{modenanalyse\_2fe2s}}]
$\lambda_X(i)$ is the \textbf{thermal bond reorganization energy}
of mode $i$ along bond class $X$:
\begin{itemize}\setlength{\itemsep}{0pt}
  \item Form: identical to Marcus-Hush
        ($\tfrac{1}{2}\mu\omega^2 r^2$), but with thermally scaled
        $r$ instead of a static reactant--product offset.
  \item Related to the Huang-Rhys factor: $S_X(i) =
        \lambda_X(i)/(\hbar\omega_i)$.
  \item In the low-temperature limit:
        $\lambda_X(i) \to (\hbar\omega_i/4)\,\alpha_X^2(i)$
        -- purely quantum-mechanical, $T$-independent.
  \item Direct spectroscopic meaning: $\lambda_X(i)$ is the energy
        of the bond-$X$ modulation in mode $i$, measured at the
        QHO expectation value.
\end{itemize}

\textbf{This quantity is related to, but not identical with, the
classical Marcus-Hush activation energy of a reaction between two
electronic states.} For the latter, two separate DFT calculations
would be required.
\end{tcolorbox}

\subsection{NRVS theory as relevant for band identification}
\label{sec:nrvs}

The tool \emph{does not} compute NRVS spectra -- these are
synthesized from the same DFT frequency calculation using
specialized external tools. But the modulation spectra
$M_X(\omega)$ and the cumulative reorganization curves
$\Lambda_X(\omega)$ computed in \S\ref{sec:reorg} can be compared
directly with NRVS bands -- which is why we need NRVS theory as an
interpretive framework.

\subsubsection*{What does NRVS measure?}

NRVS \cite{Sturhahn2004} measures the
\emph{$^{57}$Fe-projected vibrational density of states}:
\begin{equation}
  D_\text{Fe}(\omega) = \sum_i |\mathbf{e}_\text{Fe}^{(i)}|^2 \cdot
                          \delta(\omega - \omega_i)
  \label{eq:fe-pdos}
\end{equation}
weighted by the Lamb-M\"o{\ss}bauer factor $f_\text{LM}$ and a
multi-phonon convolution. For band identification it is usually
sufficient to use the Fe-pDOS without multi-phonon contributions
-- we call this $g_\text{Fe}(\omega)$.

Important properties:
\begin{itemize}\setlength{\itemsep}{0pt}
  \item NRVS is \textbf{element-selective}: only modes with Fe
        participation appear. Modes containing pure ligand motions
        (e.g., C-H stretching) are not visible.
  \item NRVS is \textbf{not symmetry-selective}: unlike Raman/IR,
        all modes are visible regardless of their $D_{2h}$
        symmetry. This makes NRVS particularly useful for
        iron-sulfur clusters, where many modes are dark in
        IR/Raman by selection rules.
  \item NRVS is \textbf{quantitative}: the integral of the pDOS is
        normalized to $1$ (or $3 N_\text{Fe}$ in some
        conventions). Band intensities are directly mode
        couplings.
\end{itemize}

\subsubsection*{Band structure in [2Fe-2S] proteins}

The typical [2Fe-2S] NRVS band structure (Gee 2021
\cite{Gee2021}; a comprehensive overview of NRVS modes is
given by Wang \emph{et al.}\ 2025 \cite{Wang2025}) has three main
regions:

\paragraph{Low-frequency region (50--200\,cm$^{-1}$).}
Skeletal vibrations of the entire protein. Cluster-relevant:
cluster-plane motions (translational + rotational in the protein
matrix), cluster-OOP buckling (acoustic mode mixing with ligands).
$\lambda_X$ is small in all channels here
($\lambda_X \approx \hbar\omega/4 \approx 12$--$50$\,cm$^{-1}$ per
mode), but the \emph{number} of modes is large.

\paragraph{Mid-frequency region (200--400\,cm$^{-1}$).}
This is where most of the cluster bond stretching modes lie:
\begin{itemize}\setlength{\itemsep}{0pt}
  \item Fe-N(His) stretching:
        260--290\,cm$^{-1}$ in mitoNEET-type Cys$_3$His$_1$
        coordination (Gee 2021 \cite{Gee2021}: computed
        223\,cm$^{-1}$ for Fe(II)-N$_\delta$(His) in the reduced
        state, around 270--290\,cm$^{-1}$ in the oxidized state).
        For classical Rieske Cys$_2$His$_2$ clusters: a doublet at
        266 + 274\,cm$^{-1}$ (Kuila \emph{et al.}\ 1992
        \cite{Kuila1992}, Fe(III)-N), assigned to the protonated
        (266) and deprotonated (274) imidazole, respectively.
  \item Fe-Fe cluster breathing: $\sim 200$\,cm$^{-1}$, depending
        on the cluster geometry and ligands.
  \item Fe-S$_\text{bridge}$ symmetric stretch: 300--340\,cm$^{-1}$.
  \item Fe-S$_\text{Cys}$ stretching: 280--320\,cm$^{-1}$.
\end{itemize}
For a pure Fe-X stretch mode with $\alpha_X \approx 1$, we have
$\lambda_X \approx \hbar\omega/4 \approx 60$--$100$\,cm$^{-1}$ --
a substantial contribution per mode.

\paragraph{High-frequency region (400--800\,cm$^{-1}$).}
\begin{itemize}\setlength{\itemsep}{0pt}
  \item Fe-S$_\text{bridge}$ asymmetric stretch:
        380--420\,cm$^{-1}$.
  \item Imidazole ring modes with Fe participation:
        600--700\,cm$^{-1}$.
  \item Cluster-ligand bending modes with Fe contribution:
        500--700\,cm$^{-1}$.
\end{itemize}

\subsubsection*{How to compare $\Lambda_X(\omega)$ to NRVS}

The cumulative reorganization curve
$\Lambda_X(\omega) = \sum_{\omega_i \leq \omega} \lambda_X(i)$ is a
\textbf{step function}: every mode with non-trivial $\lambda_X$
contribution produces a step of height $\lambda_X(i)$ at
$\omega = \omega_i$.

In comparison with the NRVS spectrum:
\begin{itemize}\setlength{\itemsep}{0pt}
  \item Where there is a band in the Fe-pDOS, there should be a
        step in $\Lambda_X(\omega)$ -- for the channel $X$ that
        dominates this band.
  \item The height of the step is proportional to
        $\hbar\omega \cdot \alpha_X^2$. A step of 50\,cm$^{-1}$ at
        $\omega = 350$\,cm$^{-1}$ corresponds to
        $\alpha_X^2 = 4 \cdot 50/350 = 0.57$ -- a mode with about
        57\,\% bond contribution along $X$.
  \item If the NRVS spectrum shows a band but all
        $\Lambda_X(\omega)$ curves stay flat there, the mode is
        not strongly active in any of our 5 channels -- it is
        likely a ligand-internal mode with only minor Fe
        participation.
\end{itemize}

This correspondence makes $\Lambda_X(\omega)$ a direct
\textbf{band-assignment tool}. When a reviewer asks ``which bond
is active in band X?'', one plots
$\Lambda_\text{FeFe}, \Lambda_\text{FeN}, \Lambda_\text{FeS},
\Lambda_\text{NH}, \Lambda_\text{HA}$ above the NRVS spectrum and
reads off the answer.

\subsubsection*{Low-temperature consistency}

NRVS measurements are typically performed at 5--40\,K (liquid He)
to maximize Lamb-M\"o{\ss}bauer resolution and minimize thermal
line broadening. In this low-temperature region, our $\lambda_X$
is identical with the zero-point value
$\hbar\omega/4 \cdot \alpha_X^2$ -- i.e., our computed $\lambda_X$
\emph{is} what NRVS measures, modulo spectral broadening and
multi-phonon effects.

Consequence: for band identification one does NOT need to compute
at the exact measurement temperature. A run at \texttt{temp\_k = 5}
yields the low-temperature values; a run at \texttt{temp\_k = 300}
yields the room-temperature values; in the NRVS-relevant frequency
region ($\omega > 200$\,cm$^{-1}$) both are practically identical,
because the Bose-Einstein population vanishes there.

\subsubsection*{What NRVS \emph{does not} see}

\begin{itemize}\setlength{\itemsep}{0pt}
  \item \textbf{Pure N-H stretching modes} (imidazole NH at 3000+
        cm$^{-1}$): outside the typical NRVS measurement range
        (instruments end at 800--1000\,cm$^{-1}$), and the mode
        has only tiny Fe participation anyway. Consequence:
        $\Lambda_\text{NH}$ in our typical filters
        (\texttt{freq\_max = 800}) is always very small -- see the
        report table.
  \item \textbf{Pure OH stretching modes}: same.
  \item \textbf{Backbone modes} without Fe contribution: remain
        dark in the pDOS.
\end{itemize}

These points are important when justifying band assignments in the
report: if a mode is dark in the pDOS but our $\Lambda_X(\omega)$
shows a step there, this is a \emph{hint} that the mode exists and
has Fe contribution but a very small Lamb-M\"o{\ss}bauer factor
(the mode is not coherent with Fe). This is particularly
interesting in PCET studies.

% =========================================================================
% End of new theory chapters
% =========================================================================

\subsection{Mode-resolved bonding reorganization energies}
\label{sec:reorg}

The package implements a \emph{conceptual pivot}: away from
classifying individual modes as ``PCET candidates'' (with
thresholds, filters, $\tanh$ saturation), toward direct
computation of mode-resolved bond reorganization contributions, in
Marcus-Hush form but with thermally scaled mode amplitudes.

Before describing the pipeline, a clarification of the physical
meaning of the computed quantities.

\subsubsection*{What $\lambda_X$ is and what it is not}

The \textbf{classical Marcus-Hush reorganization energy}
\cite{Marcus1956,Marcus1993} is
\begin{equation}
  \lambda^{\text{Marcus}}_\text{total}
  = \frac{1}{2}\sum_i \mu_i\,\omega_i^2\,(\Delta Q_i)^2,
  \label{eq:marcus-hush}
\end{equation}
where $\Delta Q_i$ is the \emph{static} displacement of the
harmonic mode coordinate $i$ between reactant and product states
(activation energy of the reaction). Computing this quantity
requires two separate frequency calculations (reactant and
product) plus a mode mapping between them -- effort that we
\emph{do not} undertake here.

\textbf{What we compute instead} is a related but distinct
quantity: the mode-resolved contribution energy of the
\emph{thermal bond-length fluctuations} per bonding channel. From
the eigenvectors of the frequency calculation, for each mode $i$
and each bond $X$ the thermally scaled bond modulation
$\Delta r_X(i)$ is computed (displacement vectors multiplied by
$u_\text{rms}$ from the QHO expectation value), and from this
\begin{equation}
  \lambda_X(i) = \tfrac{1}{2}\,\mu_i\,\omega_i^2\,(\Delta r_X(i))^2.
  \label{eq:lambda-mode}
\end{equation}
The sum over all modes,
$\Lambda_X^\text{total} = \sum_i \lambda_X(i)$, is the thermal
bond reorganization energy of bond $X$. At $T \to 0$ this reduces
(with $\mu = m_\text{mode}$) to
\begin{equation}
  \lambda_X(i) \to \tfrac{\hbar\omega_i}{4} \cdot \alpha_X^2(i),
\end{equation}
where $\alpha_X(i)$ is the dimensionless ``mode contribution'' to
bond $X$, defined implicitly by
$\Delta r_X(i) = u_\text{rms}(i) \cdot \alpha_X(i)$. For a pure
stretching mode of bond $X$, $\alpha_X(i) = 1$ and the
reorganization contribution at $T \to 0$ is $\hbar\omega/4$
(zero-point energy divided by 2). In practice $\alpha_X(i)$ can
exceed $1$ for modes where the two bond partners move in
strongly opposite directions (relative-displacement enhancement);
this is geometrically legitimate and does not violate any
conservation law -- the eigenvector normalisation
$\sum_i \lVert\evec_i\rVert^2 \approx 1$ is a sum over
\emph{atoms}, not over bonds.

\paragraph{Significance for PCET studies.}
\begin{itemize}\setlength{\itemsep}{0pt}
  \item $\Lambda_X^\text{total}$ is a well-defined spectroscopic
        observable. It measures, mode-resolved, which bond fluctuates
        most strongly in the thermal vibrations.
  \item Suitable for \textbf{NRVS band identification}: frequency
        ranges with large $\Lambda_X(\omega)$ steps are the bands
        in which bond $X$ is strongly modulated. Direct comparison
        with NRVS spectra is possible.
  \item Suitable for \textbf{system comparisons at the same
        temperature}: WT vs.\ mutant, prot vs.\ deprot. Scales
        consistently. In particular: in the deprot variant, NH
        and HA vanish by channel definition; the others remain
        comparable.
  \item Related to the \textbf{Huang-Rhys factor}
        \cite{HuangRhys1950} and to Franck-Condon theory of
        vibronic spectra.
\end{itemize}

\paragraph{What the quantity does not provide.}
\begin{itemize}\setlength{\itemsep}{0pt}
  \item \emph{No} classical Marcus-theory activation energies for
        a specific ET or PT reaction -- this would require
        separate reactant and product geometries.
  \item \emph{No} reaction-rate prediction.
  \item \emph{No} direct comparison with literature values of the
        Marcus $\lambda$ for a reaction (insofar as those use
        static reactant--product differences).
\end{itemize}

\paragraph{The Hammes-Schiffer view.}
Hammes-Schiffer and Soudackov \cite{HammesSchifferSoudackov2008}
show that PCET in proteins typically activates \emph{many modes
simultaneously}. The notion of a single ``promoter mode'' is a
special case (small models), not the rule. Our mode-resolved view
makes precisely this multi-mode reality explicit: instead of a
top-mode classification we provide the distribution of bond
modulations across all modes.

\subsubsection*{Five bonding channels}

Per mode, five bonding channels are evaluated:

\begin{itemize}\setlength{\itemsep}{0pt}
  \item \texttt{FeFe} -- Fe-Fe cluster breathing (ET reorganization).
  \item \texttt{FeN} -- Fe-N(His) bond (PT-relevant in
        Rieske-type systems).
  \item \texttt{FeS} -- Fe-S(Cys) bond and Fe-S$_\text{bridge}$
        within the cluster.
  \item \texttt{NH} -- N-H stretching of the protonated His.
  \item \texttt{HA} -- H-acceptor distance (hydrogen bond),
        computed in \emph{reaction-coordinate mode} (see below).
\end{itemize}

NH and HA are skipped for deprotonated His ligands (no H
available). FeS includes both Fe-Cys bonds and cluster-internal
Fe-$\mu$S bonds; the per-bond breakdown is provided in the sheet
\texttt{Reorg\_per\_bond}.

\subsubsection*{Geometric modulation $\Delta r_X$}

For a mode with thermally scaled displacement vectors
$\mathbf{e}_a, \mathbf{e}_b$ at the two bond partners (positioned
at $\mathbf{r}_a, \mathbf{r}_b$), the signed bond-length modulation is
\begin{equation}
  \Delta r_X^\text{(class.)}
  = (\mathbf{e}_a - \mathbf{e}_b) \cdot \hat{\mathbf{n}}_{ab},
  \qquad
  \hat{\mathbf{n}}_{ab} = \frac{\mathbf{r}_a - \mathbf{r}_b}{\|\mathbf{r}_a - \mathbf{r}_b\|}.
  \codref{reorganization.signed\_dr\_along\_axis}
  \label{eq:dr-classical}
\end{equation}
Convention: $\Delta r > 0$ means ``the bond stretches''.
This classical definition is used for FeFe, FeN, FeS, and NH.

\paragraph{HA as a reaction coordinate.}
For the HA channel, equation~\ref{eq:dr-classical} does \emph{not}
measure what is physically relevant for PT: in a high-frequency
C=O vibration in the protein backbone, the acceptor moves
strongly and the H stays nearly fixed, but the distance
$|\mathbf{r}_H - \mathbf{r}_A|$ changes -- and would be counted
as an HA modulation, even though it has no PT character.

The HA modulation measures the
\emph{motion of the H relative to the donor N}, projected onto
the N$\to$acceptor axis:
\begin{equation}
  \Delta r_\text{HA}^\text{(react.\ coord.)}
  = (\mathbf{e}_H - \mathbf{e}_N) \cdot \hat{\mathbf{n}}_{NA},
  \qquad
  \hat{\mathbf{n}}_{NA} = \frac{\mathbf{r}_A - \mathbf{r}_N}{\|\mathbf{r}_A - \mathbf{r}_N\|}.
  \codref{reorganization.compute\_mode\_modulations}
  \label{eq:dr-reaction-coord}
\end{equation}
If only the acceptor moves (skeletal mode), then
$\mathbf{e}_H - \mathbf{e}_N \approx 0$ and hence
$\Delta r_\text{HA} = 0$ -- no PT contribution. A genuine PT
motion (H moving from N toward A) yields, by contrast,
$\Delta r_\text{HA} > 0$.

\paragraph{Main acceptor per H.}
For each His-H atom, all acceptors within
\texttt{pcet\_hbond\_cutoff\_a} (default 4\,\AA) are detected, but
only the \emph{main acceptor} (highest Gaussian weight
$\exp\bigl(-(r_{\text{eq}} - r_0)^2/(2\sigma^2)\bigr)$ with
$r_0 = 2.8$\,\AA, $\sigma = 0.4$\,\AA) is used as the HA channel.
This avoids double-counting acceptors -- a single H physically
moves in only one direction. Concretely for PCET, the logic is:
per His imidazole there is one NE2-H (or ND1-H), and at this H
sits exactly one dominant H-bond acceptor (typically a water,
backbone carbonyl O, or an Asp/Glu O).

\paragraph{Effect of the Gaussian weight.}
The Gaussian weight
$w = \exp\bigl(-(r_{\text{eq}} - r_0)^2/(2\sigma^2)\bigr)$
is used \emph{only} for main-acceptor selection -- not as a
damping factor in the lambda aggregation. The binary logic is
now: \emph{is an H-bond acceptor within
\texttt{pcet\_hbond\_cutoff\_a} reach or not?} If acceptors are in
range, the HA mode counts fully (weight $= 1.0$); if not, the
channel is dropped. The geometric damping of distant acceptors
arises automatically from the smaller $\Delta r_{HA}$ values
(distant H-bonds are weakly modulated anyway); additional
multiplication by $w < 1$ would be a double penalty.

\subsubsection*{Reduced mass: two conventions}

For the per-mode $\lambda$ contribution we have two choices for
the reduced mass:

\begin{itemize}
  \item $\mu_\text{pair}$ -- bond-specific reduced mass, computed
        from the two bonded atoms ($\mu_\text{NH} \approx 0.94$\,amu,
        $\mu_\text{FeFe} \approx 27.92$\,amu, etc.). Closer to the
        \emph{isolated-oscillator} picture: good for
        mode-to-mode comparison \emph{within the same} bond.

  \item $\mu_\text{mode}$ -- the mode's own reduced mass from the
        Gaussian output. By construction this satisfies the
        orthogonality of the eigenmodes, so the $\lambda$
        contributions add up \emph{without double counting} --
        this is the \textbf{Marcus-Hush-correct} choice for the
        sum in eq.~\ref{eq:marcus-hush}.
\end{itemize}

Both are written out (columns \texttt{lambda\_X\_pair\_cm1} and
\texttt{lambda\_X\_mode\_cm1} in the sheet
\texttt{Reorganization\_energy}). For the system aggregation in
\texttt{Reorg\_total}, $\mu_\text{mode}$ is used.

\subsubsection*{Per-mode reorganization contribution $\lambda_X$}

Formula \eqref{eq:lambda-mode} above,
$\lambda_X(i) = \tfrac{1}{2}\,\mu_i\,\omega_i^2\,(\Delta r_X(i))^2$,
is applied with the geometric modulations $\Delta r_X(i)$ defined
above and the chosen reduced-mass convention ($\mu_\text{pair}$ or
$\mu_\text{mode}$).

Unit conversions: $\omega_\text{SI} = 2\pi c\,\omega_{\text{cm}^{-1}}$,
$\mu_\text{SI} = \mu_{\text{amu}} \cdot \text{u}$,
$\Delta r_\text{SI} = \Delta r_{\text{\AA}} \cdot 10^{-10}$\,m;
$\lambda_{\text{cm}^{-1}} = \lambda_\text{J} / (h c)$.

Important: because of the thermal scaling
$\mathbf{e}_a = \mathbf{e}_a^\text{normalized} \cdot u_\text{rms}$,
$\Delta r_X(i)$ lies on the scale of the QHO probability
distribution, not of a reactant--product difference. For pure
stretching modes of bond $X$, $\lambda_X(i)$ converges at $T \to 0$
to $\hbar\omega_i/4$ (which equals $\omega_i/4$ in cm$^{-1}$); for
mixed modes, this is multiplied by the squared mode contribution
$\alpha_X^2$.

\subsubsection*{System aggregation}

\begin{equation}
  \Lambda_X^\text{total} = \sum_{i=1}^{N_\text{modes}} \lambda_X(i)
  \codref{reorganization.compute\_total\_reorganization}
  \label{eq:lambda-total}
\end{equation}

per channel $X \in \{\text{FeFe, FeN, FeS, NH, HA}\}$.

The sums are also written out \emph{frequency-resolved}:

\begin{itemize}
  \item \textbf{Modulation spectra} $M_X(\omega)$ -- Gaussian-broadened
        accumulation of the per-mode bond modulations along the
        frequency axis:
        \begin{equation}
          M_X(\omega)
          = \sum_{i=1}^{N_\text{modes}}
            \Delta r_X^\text{RSS}(i)\,
            \mathcal{G}(\omega - \omega_i, \sigma_M),
          \quad
          \mathcal{G}(x,\sigma) = \frac{e^{-x^2/(2\sigma^2)}}{\sigma\sqrt{2\pi}}.
          \codref{reorganization.compute\_modulation\_spectra}
          \label{eq:M-spectrum}
        \end{equation}
        Default broadening: $\sigma_M = 5$\,cm$^{-1}$, configurable
        via \texttt{reorg\_spectrum\_sigma\_cm1}.

        \medskip
        \noindent\textbf{What is $\Delta r_X^\text{RSS}(i)$ exactly?}
        Each parent channel
        $X \in \{\text{FeFe}, \text{FeN}, \text{FeS}, \text{NH}, \text{HA}\}$
        can have several \emph{sub-channels}: e.g.\ FeN comprises the
        two Fe-N(His$_a$) and Fe-N(His$_b$) sub-bonds, FeS comprises
        all four Fe-S(Cys$_k$) sub-bonds (or two for Cys$_2$His$_2$
        clusters). For mode $i$ we therefore aggregate sub-channel
        modulations as the weighted root-sum-of-squares
        \begin{equation}
          \Delta r_X^\text{RSS}(i)
            \;=\; \sqrt{\sum_{s \in \text{sub-channels}(X)}
                          w_s\,(\Delta r_X^{(s)}(i))^2},
          \codref{reorganization.aggregate\_by\_parent}
          \label{eq:dr-rss}
        \end{equation}
        where the weights $w_s$ are the geometric weights from
        \texttt{build\_channels\_from\_coord\_info} (Gaussian
        H$\cdot\cdot\cdot$acceptor weighting for HA, unity for
        FeFe/FeN/FeS/NH).  For single-sub-channel parents (FeFe, NH,
        HA-per-acceptor) this reduces to
        $\sqrt{w}\,|\Delta r_X(i)|$, recovering the simple absolute
        modulation. The RSS aggregation is the \emph{energetically
        consistent} choice: it satisfies
        $\lambda_X^\text{parent}(i) = \tfrac{1}{2}\mu_i\omega_i^2
        \,(\Delta r_X^\text{RSS}(i))^2$ with $\lambda_X^\text{parent}$
        being the weighted sum of the sub-channel $\lambda$'s.

  \item \textbf{Co-modulation spectra} -- geometric mean of two
        $M_X$ spectra, large only where \emph{simultaneously}
        two bonds are modulated:
        \begin{align}
          C_\text{PCET}(\omega) &= \sqrt{M_\text{HA}(\omega)\, M_\text{FeFe}(\omega)},\\
          C_{\text{PT-FeN}}(\omega) &= \sqrt{M_\text{HA}(\omega)\, M_\text{FeN}(\omega)},\\
          C_{\text{ET-FeS}}(\omega) &= \sqrt{M_\text{FeFe}(\omega)\, M_\text{FeS}(\omega)}.
        \end{align}

  \item \textbf{Cumulative reorganization curves} $\Lambda_X(\omega)$
        -- monotonically increasing, converging at the spectrum end
        to $\Lambda_X^\text{total}$:
        \begin{equation}
          \Lambda_X(\omega_\text{cut})
          = \sum_{i\,:\,\omega_i \le \omega_\text{cut}} \lambda_X(i).
        \end{equation}
        Steps in $\Lambda_X(\omega)$ mark reorganization-active
        frequency windows that can be compared directly with NRVS
        spectra.
\end{itemize}

\subsubsection*{Excel output}
\label{sec:reorg-excel}

Five sheets per run or per frequency window:

\begin{itemize}
  \item \texttt{Reorganization\_energy} -- per mode 21 columns:
        \texttt{freq\_cm1} and, for each channel, four values
        ($\Delta r_\text{rms}$, $\Delta r_\text{sum-signed}$,
        $\lambda_\text{pair}$, $\lambda_\text{mode}$).

  \item \texttt{Reorg\_total} -- 5-row summary: per channel
        $\Lambda_X^\text{total}$ in both $\mu$ conventions plus
        the number of contributing modes.

  \item \texttt{Reorg\_per\_bond} -- per individual bond
        (FeS\_Cys207, FeS\_Cys216, FeS\_Cluster\_*, FeN\_His255,
        etc.) the contribution, with acceptor weight and number
        of modes.

  \item \texttt{Modulation\_spectra} -- frequency-resolved: five
        $M_X(\omega)$ curves plus three co-indicator curves
        $C_\text{PCET}, C_\text{PT-FeN}, C_\text{ET-FeS}$.

  \item \texttt{Lambda\_cumulative} -- five $\Lambda_X(\omega)$
        curves for direct comparison with NRVS spectra.
\end{itemize}


\subsection{HP/standard frequency consistency check}
\label{sec:hp-std}

As indicated in \S\ref{sec:hp-std-check}, with
\texttt{freq=hpmodes} Gaussian writes the frequencies
\emph{twice} into the log file -- once in the HP block (5 decimal
places) and once in the standard block (2 decimal places). Both
blocks refer to the same eigenvalue calculation; their frequencies
\emph{must} therefore agree (within the writing precision of
0.01\,cm$^{-1}$).

\subsubsection*{Implementation}

\texttt{logio.check\_hp\_std\_frequency\_consistency} compares,
during the run:

\begin{enumerate}
  \item Mode by mode: for each mode index, the HP frequency is
        compared to the standard-block frequency.
  \item Tolerance: 0.01\,cm$^{-1}$ (default), configurable via
        \texttt{cfg.hp\_std\_freq\_tolerance\_cm1}.
  \item Where they disagree, the maximum difference, the mode
        index in question, and both frequencies are logged.
\end{enumerate}

\subsubsection*{What discrepancies mean}

\begin{center}
\begin{tabularx}{\textwidth}{l X}
\toprule
\textbf{Symptom} & \textbf{Possible cause} \\
\midrule
1--2 modes with $\Delta\omega > 0.01$ &
  Numerical truncation in the Gaussian output (rare, harmless) \\
Many modes systematically offset &
  Wrong block pairing (HP block from another job, restart
  mismatch) \\
Complete disagreement &
  Aborted job, partial log file, or two different log files
  concatenated \\
\bottomrule
\end{tabularx}
\end{center}

In all cases, the discrepancy is \emph{not} treated as fatal --
the program continues and uses the HP values. The diagnosis is
written to the report so that follow-up investigation is possible.

\begin{warnbox}
\textbf{Practical note.} In very old log files (Gaussian 09 or
earlier), there are occasionally sign discrepancies between HP and
standard block for the lowest frequencies ($\omega \to 0$): a mode
that appears as $-2.3$\,cm$^{-1}$ in the HP block and as
$+2.3$\,cm$^{-1}$ in the standard block. This is a Gaussian
artifact of the translation/rotation handling. The consistency
check warns about this, but the behavior is harmless -- such modes
are usually excluded anyway because $\omega < $ threshold.
\end{warnbox}

% =========================================================================

\part{Application}
% =========================================================================

% =========================================================================
\section{Configuration}
\label{sec:konfig}

All run parameters are collected in a ``Config'' object.
``Config'' is a Python ``dataclass'' with default values for all
optional fields (\texttt{src/modenanalyse\_2fe2s/config.py}). Three
fields are \emph{required}: \texttt{log\_file}, \texttt{output\_dir},
and -- if ligand resolution is desired -- \texttt{pdb\_file}.

\subsection{Example configuration}

The typical form, runnable in Spyder (see ``example\_run.py''):

\begin{lstlisting}[style=python]
from modenanalyse_2fe2s import Config, run_analysis

cfg = Config(
    log_file   = r"D:\Data\cluster_freq.log",
    pdb_file   = r"D:\Data\cluster.pdb",
    output_dir = r"D:\Data\results",
    temp_k     = 40.0,
    freq_max   = 800.0,
    freq_windows = [(0, 100), (100, 300), (300, 500), (500, 700)],
)

run_analysis(cfg)
\end{lstlisting}

All other fields keep their defaults. The following tables list
the fields grouped by function.

\subsection{Complete parameter reference}
\label{sec:param-referenz}

The following overview shows all 48 configuration fields with
their default values. Detailed descriptions and recommendations
appear in the topic sections below.

\begin{longtable}{p{0.32\textwidth} p{0.20\textwidth} p{0.40\textwidth}}
\caption{Complete overview of all configuration fields with default
values. An annotated example TOML file with all fields is provided
as \texttt{run\_example\_full.toml} in the distribution directory.}
\label{tab:param-overview} \\
\toprule
\textbf{Field} & \textbf{Default} & \textbf{Section} \\
\midrule
\endfirsthead
\toprule
\textbf{Field} & \textbf{Default} & \textbf{Section} \\
\midrule
\endhead
% ── Input ───────────────────────────────────────────────
\texttt{log\_file}                  & --- (required)        & Input \\
\texttt{output\_dir}                & --- (required)        & Input \\
\texttt{pdb\_file}                  & \texttt{""}          & Input \\
\texttt{pdb\_chain}                 & \texttt{"A"}         & Input \\
\texttt{logname\_suffix}            & \texttt{""}          & Input \\
\midrule
% ── Frequency ───────────────────────────────────────────
\texttt{freq\_min}                  & \texttt{None}        & Frequency \\
\texttt{freq\_max}                  & \texttt{None}        & Frequency \\
\texttt{freq\_windows}              & \texttt{None}        & Frequency \\
\midrule
% ── Thermo ──────────────────────────────────────────────
\texttt{temp\_k}                    & \texttt{None}        & Thermo \\
\texttt{amplitude}                  & \texttt{1.0}         & Thermo \\
\midrule
% ── Cluster ─────────────────────────────────────────────
\texttt{cluster\_index}             & \texttt{0}           & Cluster \\
\texttt{analyze\_all\_clusters}     & \texttt{False}       & Cluster \\
\texttt{fe\_s\_cutoff}              & \texttt{3.0} \AA     & Cluster \\
\texttt{fe\_fe\_cutoff}             & \texttt{3.5} \AA     & Cluster \\
\texttt{strict\_cluster}            & \texttt{False}       & Cluster \\
\midrule
% ── Ligand ──────────────────────────────────────────────
\texttt{fe\_coord\_cutoff\_n}       & \texttt{2.50} \AA    & Ligand \\
\texttt{fe\_coord\_cutoff\_s}       & \texttt{2.70} \AA    & Ligand \\
\texttt{fe\_coord\_cutoff\_o}       & \texttt{2.50} \AA    & Ligand \\
\texttt{include\_hn\_vibration}     & \texttt{True}        & Ligand \\
\midrule
% ── OOP/INP ─────────────────────────────────────────────
\texttt{mode\_type\_threshold}      & \texttt{60.0\,\%}    & OOP \\
\texttt{mode\_type\_detail\_thresholds} & \texttt{(60, 75, 90)} & OOP \\
\midrule
% ── SCSD/SSE ─────────────────────────────────────────────
\texttt{analyze\_scsd}              & \texttt{True}        & SCSD \\
\texttt{analyze\_sse}                & \texttt{True}        & SCSD \\
\texttt{sse\_chain}                  & \texttt{""}          & SCSD \\
\midrule
% ── PCET / Marcus-Hush reorg ────────────────────────────
\texttt{pcet\_enabled}              & \texttt{True}        & PCET \\
\texttt{pcet\_hbond\_cutoff\_a}     & \texttt{4.0} \AA     & PCET \\
\texttt{pcet\_acceptor\_r0\_a}      & \texttt{2.8} \AA     & PCET \\
\texttt{pcet\_acceptor\_sigma\_a}   & \texttt{0.4} \AA     & PCET \\
\texttt{reorg\_spectrum\_sigma\_cm1} & \texttt{5.0} cm$^{-1}$  & PCET \\
\texttt{reorg\_spectrum\_step\_cm1} & \texttt{0.5} cm$^{-1}$  & PCET \\
\midrule
% ── Embedding ───────────────────────────────────────────
\texttt{export\_embedding\_plots}   & \texttt{True}        & Embedding \\
\texttt{umap\_n\_neighbors}         & \texttt{None} (auto) & Embedding \\
\midrule
% ── Numerics ────────────────────────────────────────────
\texttt{include\_hydrogen}          & \texttt{True}        & Numerics \\
\texttt{sigma\_eigvec}              & \texttt{5e-4}        & Numerics \\
\texttt{sigma\_coord}               & \texttt{1e-3}        & Numerics \\
\texttt{amplitude\_threshold}       & \texttt{0.001}       & Numerics \\
\texttt{coord\_match\_tol}          & \texttt{1.0} \AA     & Numerics \\
\texttt{scan\_chunk\_mb}            & \texttt{16}          & Numerics \\
\texttt{show\_errors\_in\_excel}    & \texttt{True}        & Numerics \\
\texttt{interp\_step}               & \texttt{0.05} cm$^{-1}$ & Numerics \\
\texttt{interp\_boundary\_mode}     & \texttt{"context"}   & Numerics \\
\texttt{interp\_edge\_extend}       & \texttt{0.5} cm$^{-1}$ & Numerics \\
\texttt{interp\_context\_cm1}       & \texttt{5.0} cm$^{-1}$ & Numerics \\
\texttt{use\_cache}                 & \texttt{True}        & Numerics \\
\bottomrule
\end{longtable}



\subsection{Input and output}

\begin{longtable}{p{0.27\textwidth} p{0.16\textwidth} p{0.50\textwidth}}
\toprule
\textbf{Field} & \textbf{Default} & \textbf{Description} \\
\midrule
\endhead
\texttt{log\_file}      & \emph{(required)} & Path to the DFT
  frequency file. Supported: Gaussian-16 \texttt{.log} (with or
  without \texttt{freq=hpmodes}) and ORCA \texttt{.hess} (auto-detected
  from the extension). \\
\texttt{output\_dir}    & \emph{(required)} & Target directory.
  Created if it does not exist. \\
\texttt{pdb\_file}      & \texttt{""}     & PDB file with the same
  cluster (ligand resolution; optional but recommended). \\
\texttt{pdb\_chain}     & \texttt{"A"}     & Chain identifier in
  the PDB file. \\
\texttt{logname\_suffix} & \texttt{""}     & Optional suffix for
  output file names (useful for several runs from the same log
  file, e.g.\ different clusters). \\
\bottomrule
\end{longtable}

\subsection{Frequency range and windows}

\begin{longtable}{p{0.27\textwidth} p{0.16\textwidth} p{0.50\textwidth}}
\toprule
\textbf{Field} & \textbf{Default} & \textbf{Description} \\
\midrule
\endhead
\texttt{freq\_min}      & \texttt{None} & Lower frequency limit in
  cm$^{-1}$. \texttt{None} means no lower limit. \\
\texttt{freq\_max}      & \texttt{None} & Upper frequency limit.
  For NRVS-relevant analyses, typically 800. \\
\texttt{freq\_windows}  & \texttt{None} & List of windows such as
  \texttt{[(0, 100), (100, 300), \ldots]}. When set, the program
  generates \emph{one Excel file per window} (multi-window mode). \\
\bottomrule
\end{longtable}

\begin{infobox}
\textbf{Convention.} \texttt{freq\_max} (and \texttt{freq\_min})
also apply to all reorganization and aggregate analyses. Since
$^{57}$Fe-projected contributions vanish above $\sim$800 cm$^{-1}$
and NRVS experiments measure in this range, an 800 cm$^{-1}$
cutoff is appropriate both experimentally and numerically.
\end{infobox}

\begin{infobox}
\textbf{Window-boundary convention (since v1.0.4).}
Multi-window mode partitions modes into bins
\texttt{[(lo$_1$,hi$_1$), (lo$_2$,hi$_2$), \ldots, (lo$_N$,hi$_N$)]}.
The boundaries are \emph{half-open from the left}:
\begin{itemize}\setlength{\itemsep}{0pt}
  \item Every window except the last uses the interval
        $[\,\text{lo},\, \text{hi}\,)$. A mode with frequency
        exactly equal to a window boundary belongs to the
        \emph{upper} window (the one starting at this value).
  \item The last window in the list uses
        $[\,\text{lo},\, \text{hi}\,]$, so its upper edge is
        included.  If \texttt{hi} is \texttt{None}/\texttt{inf}
        for the last window, it is open-ended.
\end{itemize}
This guarantees that each mode is counted in exactly one window,
preventing double-counting of $B$-factor contributions and
$\Lambda_X(\omega_\text{cut})$ values for boundary modes
(\S\ref{sec:migration-v104}).  Pre-v1.0.4 used closed-closed
intervals on both ends, which over-counted exact-boundary modes.
\end{infobox}

\subsection{Thermal scaling}

\begin{longtable}{p{0.27\textwidth} p{0.16\textwidth} p{0.50\textwidth}}
\toprule
\textbf{Field} & \textbf{Default} & \textbf{Description} \\
\midrule
\endhead
\texttt{temp\_k}       & \texttt{None} & Sample temperature in
  Kelvin for the thermal scaling of the eigenvectors.
  \texttt{None} $\to$ classical mode with amplitude $=$
  \texttt{amplitude}. Typical NRVS values: 40, 80, 100. \\
\texttt{amplitude}      & \texttt{1.0}  & Fallback amplitude in
  \AA{} when \texttt{temp\_k = None}. The default of 1.0 is the
  historical convention. \\
\bottomrule
\end{longtable}

\subsection{Cluster detection}

\begin{longtable}{p{0.27\textwidth} p{0.16\textwidth} p{0.50\textwidth}}
\toprule
\textbf{Field} & \textbf{Default} & \textbf{Description} \\
\midrule
\endhead
\texttt{fe\_s\_cutoff}  & \texttt{3.0} & Maximum Fe-S distance
  (\AA{}) within which S atoms are counted as part of the
  cluster. \\
\texttt{fe\_fe\_cutoff} & \texttt{3.5} & Maximum Fe-Fe distance
  (\AA{}) for cluster identification. \\
\texttt{cluster\_index} & \texttt{0}   & For multi-cluster systems
  (e.g.\ glutaredoxin dimer): index of the cluster to be
  analyzed. Default 0 = closest Fe-Fe pair. One run per cluster. \\
\texttt{analyze\_all\_clusters} & \texttt{False} & If \texttt{True},
  every detected cluster is analyzed in turn (one subdirectory per
  cluster). Convenient for multi-cluster systems; takes precedence
  over \texttt{cluster\_index}. \\
\texttt{strict\_cluster} & \texttt{False} & If \texttt{True}, the
  program aborts with an error when the cluster cannot be
  unambiguously identified (no fallback to nearby atoms). \\
\bottomrule
\end{longtable}

\subsection{Fe ligand detection (PDB-based)}

\begin{longtable}{p{0.27\textwidth} p{0.16\textwidth} p{0.50\textwidth}}
\toprule
\textbf{Field} & \textbf{Default} & \textbf{Description} \\
\midrule
\endhead
\texttt{fe\_coord\_cutoff\_n} & \texttt{2.50} & Maximum Fe-N
  distance (\AA{}) for ligand N (His). \\
\texttt{fe\_coord\_cutoff\_s} & \texttt{2.70} & Maximum Fe-S
  distance for ligand S (Cys). \\
\texttt{fe\_coord\_cutoff\_o} & \texttt{2.50} & Maximum Fe-O
  distance for ligand O (Asp/Glu, rare for [2Fe-2S]). \\
\texttt{include\_hn\_vibration} & \texttt{True} & N-H stretching
  analysis for protonated His ligands (sheet \texttt{His\_HN}).
  Requires \texttt{include\_hydrogen=True}. \\
\bottomrule
\end{longtable}

\subsection{OOP/INP classification}

\begin{longtable}{p{0.34\textwidth} p{0.16\textwidth} p{0.43\textwidth}}
\toprule
\textbf{Field} & \textbf{Default} & \textbf{Description} \\
\midrule
\endhead
\texttt{mode\_type\_threshold} & \texttt{60.0} & Binary OOP
  threshold in percent (\S\ref{sec:oop-inp}). \\
\texttt{mode\_type\_detail\_thresholds} & \texttt{(60, 75, 90)}
  & Three thresholds for the fine 7-level classification.
  Applied symmetrically about 50\,\%. \\
\bottomrule
\end{longtable}

\subsection{SCSD and secondary structure}

\begin{longtable}{p{0.27\textwidth} p{0.16\textwidth} p{0.50\textwidth}}
\toprule
\textbf{Field} & \textbf{Default} & \textbf{Description} \\
\midrule
\endhead
\texttt{analyze\_scsd}  & \texttt{True} & Activate Symmetry-Coordinate
  Structural Decomposition (requires \texttt{scsdpy}). \\
\texttt{analyze\_sse}    & \texttt{True} & Activate secondary-structure
  analysis (helices, sheets from PDB). \\
\texttt{sse\_chain}      & \texttt{""}   & Chain filter for
  SSE detection. Empty = same as \texttt{pdb\_chain}. \\
\bottomrule
\end{longtable}

\subsection{PCET / Marcus-Hush reorganization}

\begin{longtable}{p{0.34\textwidth} p{0.16\textwidth} p{0.43\textwidth}}
\toprule
\textbf{Field} & \textbf{Default} & \textbf{Description} \\
\midrule
\endhead
\texttt{pcet\_enabled} & \texttt{True} & Enable the Marcus-Hush
  reorganization pipeline. If \texttt{False}, $\Delta r_X$ and
  $\lambda_X$ are not computed per mode and all reorganization
  sheets are omitted. \\
\texttt{pcet\_hbond\_cutoff\_a} & \texttt{4.0} & Maximum
  H...acceptor distance (\AA{}) for the acceptor search,
  measured from the H atom. \\
\texttt{pcet\_acceptor\_r0\_a} & \texttt{2.8} & Optimal distance
  (\AA{}) for the Gaussian weighting of acceptor selection;
  matches a typical N-H$\cdots$O/N bond length. \\
\texttt{pcet\_acceptor\_sigma\_a} & \texttt{0.4} & Gaussian
  standard deviation (\AA{}) of the acceptor weighting.
  At $|r - r_0| > 2\sigma$, the acceptor effectively contributes
  nothing. \\
\texttt{reorg\_spectrum\_sigma\_cm1} & \texttt{5.0} & Gaussian
  broadening of the modulation spectra $M_X(\omega)$.
  5\,cm$^{-1}$ matches typical DFT frequency uncertainty +
  NRVS resolution. \\
\texttt{reorg\_spectrum\_step\_cm1} & \texttt{0.5} & Frequency
  grid step for $M_X(\omega)$ and $\Lambda_X(\omega)$.
  0.5\,cm$^{-1}$ produces fine bands at a tractable file size. \\
\bottomrule
\end{longtable}

\subsection{Numerical parameters}

\begin{longtable}{p{0.27\textwidth} p{0.16\textwidth} p{0.50\textwidth}}
\toprule
\textbf{Field} & \textbf{Default} & \textbf{Description} \\
\midrule
\endhead
\texttt{include\_hydrogen} & \texttt{True} & Include H atoms in
  eigenvector reads. Required for PCET. \\
\texttt{sigma\_eigvec}    & \texttt{5e-4}  & Numerical noise in
  eigenvector components (Bernoulli variance). Controls the
  threshold for ``insignificant'' contributions. \\
\texttt{sigma\_coord}     & \texttt{1e-3} & Coordinate noise in
  \AA{}. \\
\texttt{amplitude\_threshold} & \texttt{0.001} & Threshold for a
  significant atomic displacement in a mode. \\
\texttt{coord\_match\_tol} & \texttt{1.0}   & Tolerance (\AA{}) for
  PDB$\leftrightarrow$Gaussian atom matching after Kabsch alignment. \\
\texttt{interp\_step}     & \texttt{0.05}  & Step size (cm$^{-1}$)
  for the interpolated pDOS export. \\
\texttt{interp\_boundary\_mode} & \texttt{"context"} & Behavior at
  the window edge: \texttt{"context"} (with buffer zone),
  \texttt{"hard"} (sharp cut), \texttt{"extend"} (linearly
  extrapolated). \\
\texttt{interp\_edge\_extend} & \texttt{0.5} & Buffer-zone width
  (cm$^{-1}$ or fraction) in \texttt{"context"} mode. \\
\texttt{scan\_chunk\_mb}  & \texttt{16}    & Block size for log-file
  scanning in MB. For very large log files ($>$1 GB), increasing
  this can speed up reading. \\
\texttt{use\_cache}       & \texttt{True}  & Activates file caching
  for block offsets in the log file. Halves the run time on
  repeated calls. \\
\bottomrule
\end{longtable}

\subsection{Embeddings (UMAP)}

\begin{longtable}{p{0.27\textwidth} p{0.16\textwidth} p{0.50\textwidth}}
\toprule
\textbf{Field} & \textbf{Default} & \textbf{Description} \\
\midrule
\endhead
\texttt{umap\_n\_neighbors} & \texttt{None} & UMAP n\_neighbors
  parameter. \texttt{None} = auto heuristic
  $\max(5, \min(15, N/100))$; for typical [2Fe-2S] protein sizes
  ($N \approx 5\,000$--$10\,000$) this yields 15 (the paper
  standard, \cite{McInnes2018}). For very small data sets
  ($< 50$ modes), set manually to e.g.\ 5--8. \\
\texttt{export\_embedding\_plots} & \texttt{True} & Produce a PNG
  plot of the UMAP projection. Set to \texttt{False} to save
  matplotlib setup time. \\
\bottomrule
\end{longtable}

\subsection{Other}

\begin{longtable}{p{0.27\textwidth} p{0.16\textwidth} p{0.50\textwidth}}
\toprule
\textbf{Field} & \textbf{Default} & \textbf{Description} \\
\midrule
\endhead
\texttt{show\_errors\_in\_excel} & \texttt{True} & If \texttt{True},
  warnings from the run are also displayed in the Excel sheet
  \texttt{Info}. \\
\bottomrule
\end{longtable}

\subsection{Validation rules}

All fields are checked for syntactic and numerical plausibility
when \texttt{cfg.validate()} is called (executed automatically at
the start of \texttt{run\_analysis}):

\begin{itemize}
  \item Required fields are set
  \item Cutoffs $> 0$
  \item Frequency limits are ascending if both are set
  \item PCET thresholds: $0 < $ moderate $< $ strong $\leq 1$
  \item Band edges strictly ascending
  \item Line shape is from the allowed set
\end{itemize}

Validation errors are collected and reported together (no
fail-fast on a single error). This simplifies fixing several
typos in a configuration in one go.

% =========================================================================
\section{Workflow examples}
\label{sec:workflow}

Three typical scenarios are worked through completely below: a
single cluster (the standard application), multi-cluster
(glutaredoxin dimers with two clusters), and a WT-vs.-mutant
comparison (PCET validation test). A fourth scenario walks an
entire small system from the log file to band identification with
all intermediate numbers.

% -------------------------------------------------------------------------
\subsection{How to start a run}
\label{sec:workflow-how-to-start}

There are two equivalent ways to start an analysis: programmatically
from Python (\texttt{run\_analysis(cfg)}), or from the command line
with a TOML configuration file. The TOML route is recommended for
reproducible production runs because the configuration is
self-documenting and can be version-controlled.

\subsubsection*{TOML + CLI (recommended)}

The basic invocation is:

\begin{lstlisting}[style=bash]
modenanalyse-2fe2s my_run.toml
\end{lstlisting}

The TOML file path is the \emph{first positional argument} -- no
\texttt{-{}-config} flag is needed. The package can also generate
an annotated template:

\begin{lstlisting}[style=bash]
modenanalyse-2fe2s --write-template my_run.toml
# edit my_run.toml as needed
modenanalyse-2fe2s my_run.toml
\end{lstlisting}

\subsubsection*{Path pitfalls in TOML}

TOML strings with double quotes interpret backslashes as escape
characters, so Windows paths such as \texttt{"D:\textbackslash{}Data\textbackslash{}file.log"}
will produce a parser error (``Unescaped backslash in a string''
or similar). Three valid alternatives:

\begin{lstlisting}[style=bash]
# 1. Forward slashes (works on Windows too):
log_file = "D:/Data/file.log"

# 2. Single quotes = TOML literal string (no escaping):
log_file = 'D:\Data\file.log'

# 3. Escaped backslashes:
log_file = "D:\\Data\\file.log"
\end{lstlisting}

\begin{warnbox}
\textbf{Important.} Python's raw-string prefix \texttt{r"..."} is
\emph{not} valid TOML syntax. When copying paths from Python code
(e.g.\ \texttt{r"D:\textbackslash{}Data\textbackslash{}file.log"})
into a TOML file, the \texttt{r} prefix must be removed and one of
the three forms above used instead. This is a common stumbling
block for users transitioning from \texttt{run\_analysis()} calls
to TOML-based configuration.
\end{warnbox}

\subsubsection*{CLI overrides}

CLI arguments can override individual TOML fields without editing
the file -- useful for parameter scans:

\begin{lstlisting}[style=bash]
# Use my_run.toml, but force temp_k = 40 K:
modenanalyse-2fe2s my_run.toml --temp-k 40.0

# Use my_run.toml, but write to a different output directory:
modenanalyse-2fe2s my_run.toml --output-dir ./results_test
\end{lstlisting}

A complete list of available CLI flags is shown by
\texttt{modenanalyse-2fe2s -{}-help}.

\subsubsection*{Programmatic API}

For interactive analysis (e.g.\ in Spyder or Jupyter), the same
configuration can be built directly:

\begin{lstlisting}[style=python]
from modenanalyse_2fe2s import Config, run_analysis

cfg = Config(
    log_file   = r"D:\Data\system_freq.log",  # Python raw string is fine here
    pdb_file   = r"D:\Data\system.pdb",
    output_dir = r"D:\Data\results",
    temp_k     = 5.0,
    freq_max   = 800.0,
)
run_analysis(cfg)
\end{lstlisting}

Note that in \emph{Python source code} (\texttt{.py} files),
\texttt{r"..."} raw strings \emph{are} valid and are the recommended
way to write Windows paths. The restriction to non-raw strings only
applies to TOML configuration files.

% -------------------------------------------------------------------------
\subsection{Scenario A: single cluster (standard application)}
\label{sec:workflow-single}

\subsubsection*{Input data}

\begin{itemize}
  \item \texttt{cluster\_wt\_freq.log} -- Gaussian-16
        \texttt{freq=hpmodes} run of the optimized cluster model
        (cluster + first ligand sphere, ca.\ 50--100 atoms).
  \item \texttt{cluster\_wt.pdb} -- structural model with the same
        ligand residues as in the DFT run, plus additional
        residues for the SSE analysis.
\end{itemize}

\subsubsection*{Configuration}

\begin{lstlisting}[style=python]
from modenanalyse_2fe2s import Config, run_analysis

cfg = Config(
    log_file   = r"D:\Data\cluster_wt_freq.log",
    pdb_file   = r"D:\Data\cluster_wt.pdb",
    output_dir = r"D:\Data\cluster_wt_40K",
    pdb_chain  = "A",
    temp_k     = 40.0,
    freq_max   = 800.0,
    freq_windows = [(0, 100), (100, 300), (300, 500), (500, 700)],
)

run_analysis(cfg)
\end{lstlisting}

\subsubsection*{Expected output}

In the \texttt{output\_dir}, four subdirectories are created (one
per window):

\begin{itemize}
  \item \texttt{cluster\_wt\_freq\_0-100/}
  \item \texttt{cluster\_wt\_freq\_100-300/}
  \item \texttt{cluster\_wt\_freq\_300-500/}
  \item \texttt{cluster\_wt\_freq\_500-700/}
\end{itemize}

Each folder contains:

\begin{itemize}
  \item \texttt{cluster\_wt\_freq\_analysis.xlsx} -- main analysis
        (mode analysis, kernel scores, SCSD, PCET/ET, Fe-bond,
        His-HN, Info, ...)
  \item \texttt{cluster\_wt\_freq\_pDOS\_interp.xlsx} -- interpolated
        Fe-pDOS, Origin-compatible
  \item \texttt{REPORT.txt} -- plain-text run log
\end{itemize}

\subsubsection*{What to check first}

\begin{enumerate}
  \item \textbf{Open the report file} and inspect:
    \begin{itemize}
      \item Cluster detected? Are the Fe-Fe and Fe-S distances
            plausible?
      \item Ligands assigned correctly?
            (table \texttt{coord\_summary})
      \item Kabsch RMSD small ($< 0.5$\,\AA{})?
      \item HP/standard consistency: any irregularities?
      \item PCET reorg active (His ligands found,
            $\Lambda_\text{HA} > 0$)?
    \end{itemize}
  \item \textbf{Open the \texttt{Mode\_analysis} sheet}: sort by
        \emph{kern\_oop} $\to$ inspect the OOP-dominated modes.
  \item \textbf{Sheet \texttt{Reorg\_total}}: five lambda values
        for the system (FeFe, FeN, FeS, NH, HA). Which channel
        dominates? In Origin, open \texttt{Lambda\_cumulative}
        and plot the cumulative curves $\Lambda_X(\omega)$ per
        channel -- steps mark reorganization-active frequencies.
\end{enumerate}

% -------------------------------------------------------------------------
\subsection{Scenario B: multi-cluster system (dimers and beyond)}
\label{sec:workflow-multi}

Some systems carry several {[}2Fe-2S{]} clusters (e.g.,
glutaredoxin dimers with two equivalent clusters; domain constructs
with two distinct clusters; or multi-cluster complexes where two
clusters cooperate at the active site).

Since v1.0.0, the tool supports multi-cluster analysis as a
\textbf{first-class} feature: a single run analyzes all detected
clusters automatically.

\subsubsection*{Cluster detection}

\texttt{geometry.find\_all\_clusters} searches for all Fe-Fe pairs
below \texttt{cfg.fe\_fe\_cutoff} (default 3.5\,\AA{}) and groups
the bridging S atoms. Sorting: ascending by Fe-Fe distance, so the
tightest cluster comes first.

The console output displays the full list:

\begin{lstlisting}[style=bash]
[i] 2 [2Fe-2S] clusters found in the system:
    Cluster #0: Fe-Fe = 2.713 A, Fe-S = 2.245-2.318 A
    Cluster #1: Fe-Fe = 2.741 A, Fe-S = 2.262-2.305 A
\end{lstlisting}

\subsubsection*{Activating multi-cluster mode}

In TOML:

\begin{lstlisting}[style=python]
[input]
log_file   = "dimer.log"
pdb_file   = "dimer.pdb"
output_dir = "./results"

[cluster]
analyze_all_clusters = true   # <- enables the multi-cluster run
\end{lstlisting}

Or programmatically:

\begin{lstlisting}[style=python]
from modenanalyse_2fe2s import Config, run_analysis

cfg = Config(
    log_file             = r"D:\Data\dimer.log",
    pdb_file             = r"D:\Data\dimer.pdb",
    output_dir           = r"D:\Data\dimer_results",
    analyze_all_clusters = True,        # <-- multi-cluster
    temp_k               = 5.0,
    freq_max             = 800.0,
)
run_analysis(cfg)
\end{lstlisting}

\subsubsection*{What the tool does then}

With \texttt{analyze\_all\_clusters = true}, the official API
\texttt{run\_analysis(cfg)} internally calls
\texttt{find\_all\_clusters} to determine the number of clusters.
The full pipeline then runs \textbf{once per detected cluster},
each with its own subfolder:

\begin{lstlisting}[style=bash]
results/
  cluster_0/
    0-800_cm-1/
      *_REPORT.txt
      *_analysis.xlsx
      *_analysis_Embeddings.xlsx
      ...
  cluster_1/
    0-800_cm-1/
      *_REPORT.txt
      *_analysis.xlsx
      ...
  multi_cluster_summary.txt
\end{lstlisting}

The file \texttt{multi\_cluster\_summary.txt} lists, per cluster,
the geometry, the output path, and the run status (OK / error).
\texttt{cluster\_index} is ignored in this mode.

\subsubsection*{Special case: only 1 cluster in the system}

If the tool is called with \texttt{analyze\_all\_clusters = true}
but only \emph{one} cluster is found in the log file, it
automatically runs in normal single-cluster mode (no
\texttt{cluster\_0/} subfolder; output goes directly to the usual
\texttt{output\_dir/0-800\_cm-1/}). This means the flag can be
left enabled in general workflow templates without concern.

\subsubsection*{Comparing the clusters (manually)}

If the two clusters are within the same protein (e.g.\ a WT
homodimer), one expects identical reorganization values.
Differences between \texttt{cluster\_0/Reorg\_total} and
\texttt{cluster\_1/Reorg\_total} reveal electronic or structural
asymmetries between the two subunits.

\subsubsection*{Limitation: single-cluster mode remains the default}

For backward compatibility, \texttt{analyze\_all\_clusters} is
\emph{off} by default. Existing workflows that explicitly set
\texttt{cluster\_index = 1} to analyze the second cluster
continue to work. To analyze all clusters, the flag must be
explicitly enabled.

% -------------------------------------------------------------------------
\subsection{Scenario C: WT-vs.-mutant comparison (PCET validation test)}
\label{sec:workflow-wtmut}

This is the method's control test: for an HH$\to$CC mutation, the
HA reorganization energy should systematically vanish, while the
FeFe and FeS reorganization in the same frequency region should
remain. The HH$\to$CC mutant has no protonated His ligands at the
cluster anymore; a hydrogen-bond motion is therefore impossible.

\subsubsection*{Input}

\begin{itemize}
  \item \texttt{cluster\_wt\_freq.log} and \texttt{cluster\_wt.pdb}
        -- wild type, two His ligands (Cys$_2$His$_2$ ligation).
  \item \texttt{cluster\_mut\_freq.log} and
        \texttt{cluster\_mut.pdb} -- mutant
        (His$_{255}$Cys, His$_{259}$Cys), pure Cys$_4$ ligation.
\end{itemize}

\subsubsection*{Two separate runs}

\begin{lstlisting}[style=python]
from modenanalyse_2fe2s import Config, run_analysis

# WT
run_analysis(Config(
    log_file     = r"D:\Data\cluster_wt_freq.log",
    pdb_file     = r"D:\Data\cluster_wt.pdb",
    output_dir   = r"D:\Data\cluster_wt_v37",
    temp_k       = 5.0,
    freq_windows = [(0,100), (100,300), (300,500), (0,500)],
))

# Mutant
run_analysis(Config(
    log_file     = r"D:\Data\cluster_mut_freq.log",
    pdb_file     = r"D:\Data\cluster_mut.pdb",
    output_dir   = r"D:\Data\cluster_mut_v37",
    temp_k       = 5.0,
    freq_windows = [(0,100), (100,300), (300,500), (0,500)],
))
\end{lstlisting}

\subsubsection*{Expected validation findings}

\textbf{WT report:}

\begin{lstlisting}[style=bash]
[INFO] PCET reorg: active (2 His ligands, 4 H-bond acceptors).
       Marcus-Hush reorg for 5 channels: FeFe, FeN, FeS, NH, HA.
       HA in reaction-coordinate mode.
\end{lstlisting}

\textbf{Mutant report:}

\begin{lstlisting}[style=bash]
[INFO] PCET reorg: NH/HA channels skipped (no His ligands).
       Cys ligands cannot be protonated/deprotonated under
       physiological conditions. Lambda_NH = Lambda_HA = 0.
       FeFe, FeN (= 0), FeS remain active.
\end{lstlisting}

In the \texttt{Reorg\_total} sheet of the mutant variant,
$\Lambda_\text{NH}$ and $\Lambda_\text{HA}$ are zero, while
$\Lambda_\text{FeFe}$ and $\Lambda_\text{FeS}$ stay in the same
order of magnitude as in the WT (typically tens to hundreds of
cm$^{-1}$, depending on the frequency window).

\subsubsection*{What this tells us about the system}

If $\Lambda_\text{HA}^\text{(WT)}$ contributes substantially in
the 100--500\,cm$^{-1}$ region (e.g.\ $> 30$\,cm$^{-1}$) and
$\Lambda_\text{HA}^\text{(mut)} \approx 0$, this is evidence that
the HH$\to$CC mutation does indeed shut down PCET-relevant
reorganization pathways -- the mode frequencies are still formally
present in the mutant variant, but they no longer modulate
hydrogen bonds.

The following should remain comparable:
\begin{itemize}
  \item $\Lambda_\text{FeFe}$ -- the cluster breathing modes are
        largely unchanged.
  \item $\Lambda_\text{FeS}$ -- all four ligands are now Cys; the
        Fe-S reorganization changes only through a modified
        bond-length distribution, not qualitatively.
\end{itemize}

\textbf{Comparison in Origin.} Overlaying the cumulative curves
$\Lambda_X(\omega)$ from \texttt{Lambda\_cumulative} of the two
runs gives the canonical picture: the HA curve in the WT rises in
the 100--500\,cm$^{-1}$ region in steps, while the HA curve in the
mutant remains flat at zero. The FeFe and FeS curves run virtually
on top of each other.

\begin{warnbox}
\textbf{Low-temperature warning.} At \texttt{temp\_k} = 5\,K
(NRVS low temperature), the thermal amplitude $u_\text{rms}$
becomes small, and so do all $\lambda_X$. The WT-vs.-mutant
difference in $\Lambda_\text{HA}$ is qualitatively preserved
(relative ratios scale together). For direct comparison of the
values with literature numbers from classical Marcus theory
(typically computed at 298\,K), the scaling must be taken into
account.
\end{warnbox}

% -------------------------------------------------------------------------
\subsection{Scenario D: fully worked example (Cys$_4$ model cluster)}
\label{sec:workflow-worked-example}

So far we have presented workflow scenarios as \emph{instructions}
(``this is how it's done''). In this scenario we go in the
opposite direction: a concrete small system is computed
\emph{end to end} from the log file to band identification, with
all intermediate quantities. The system is the Cys$_4$ model
cluster shipped with the tool's smoke test.

We choose it deliberately because:
\begin{itemize}\setlength{\itemsep}{0pt}
  \item At 52 atoms it is the smallest sensible [2Fe-2S] model
        (Fe$_2$S$_2$ + 4 methyl-thiolate groups).
  \item 150 modes -- small enough to \emph{trace individual modes}.
  \item Cys-only -- \emph{no} His ligands, so it tests the
        ``no-His'' path: $\Lambda_\text{NH} = \Lambda_\text{HA} = 0$.
        The reorganization pipeline thus reduces to three channels
        (FeFe, FeN, FeS) and is easier to interpret.
  \item It is shipped as a test file with the repository
        (\texttt{tests/data/Cys\_2Fe-2S\_red3\_hpfrq\_opt.log.xz}),
        so the calculation is fully reproducible -- any reviewer
        can rerun it.
\end{itemize}

We go through \textbf{five steps}: (1) input data and
configuration, (2) cluster detection, (3) the mode loop and
reorganization aggregation, (4) what the numbers mean, and
(5) comparison with expectations.

\subsubsection*{Step 1: input and configuration}

From the smoke test (\texttt{tests/test\_smoke.py}):

\begin{lstlisting}[language=Python,basicstyle=\ttfamily\small]
from modenanalyse_2fe2s import Config, run_analysis

cfg = Config(
    log_file   = "Cys_2Fe-2S_red3_hpfrq_opt.log",   # 916 KB
    output_dir = "results",
    pdb_file   = "",                  # no PDB - pure cluster model
    temp_k     = 5.0,                 # low-temperature NRVS
    freq_max   = 800.0,               # NRVS measurement range
    analyze_scsd = True,              # SCSD symmetry decomposition
    pcet_enabled = True,              # even though there are no His
)
run_analysis(cfg)
\end{lstlisting}

\paragraph{What happens here.}
Since \texttt{pdb\_file = ""}, the tool runs in ``pure cluster
mode'': only atoms from the DFT run are evaluated, with no SSE
analysis and no PDB-ligand mapping. This is the simplest and
fastest workflow.

\subsubsection*{Step 2: cluster detection}

From the REPORT file:
\begin{lstlisting}[basicstyle=\ttfamily\small]
  28 heavy atoms (52 total with H)
  Cluster...
    Fe-Fe: 3.009891 +/- 1.4e-03 A
    Fe1-S1: 2.270816 +/- 1.4e-03 A
    Fe2-S1: 2.364976 +/- 1.4e-03 A
    Fe1-S2: 2.246451 +/- 1.4e-03 A
    Fe2-S2: 2.365238 +/- 1.4e-03 A
    Cluster normal n_hat: (+0.5003, -0.7220, +0.4779)
    Folding residual: 0.0918 A
    Cluster geometry: Fe-Fe=3.010 A, Fe-S(mean)=2.312 A,
                      Fe-S-Fe=81.2 deg - OK
  PCET reorg: NOT active (no His ligands at the cluster)
\end{lstlisting}

\paragraph{What we learn.}
\begin{itemize}\setlength{\itemsep}{0pt}
  \item \textbf{Fe-Fe distance 3.01\,$\angstrom$}: typical of a
        reduced [2Fe-2S]$^{+}$ cluster (mixed valence Fe$^{2+}$/
        Fe$^{3+}$). In the oxidized state ([2Fe-2S]$^{2+}$, both
        Fe$^{3+}$), the value is around $\sim 2.7$\,$\angstrom$.
  \item \textbf{Asymmetric Fe-S bonds}:
        Fe1-S1 = 2.27\,$\angstrom$ and
        Fe1-S2 = 2.25\,$\angstrom$ (both ``short'');
        Fe2-S1 = Fe2-S2 = 2.36\,$\angstrom$ (both ``long'').
        This is consistent with Fe1 = Fe$^{3+}$ (smaller ionic
        radius) and Fe2 = Fe$^{2+}$ (larger radius).
        \textbf{Localized mixed-valence cluster}.
  \item \textbf{Folding residual 0.09\,$\angstrom$}: the four
        cluster atoms lie essentially in a plane
        ($< 0.1\,\angstrom$ is excellent).
  \item \textbf{Fe-S-Fe angle 81.2\,$^{\circ}$}: typically
        between 75 and 85$^{\circ}$ for standard clusters.
        Values outside this range would be a warning sign.
  \item \textbf{No His ligands}: the tool deactivates the PCET
        reorganization pipeline. Only three channels are
        evaluated: FeFe, FeN (always 0 here), and FeS.
\end{itemize}

\subsubsection*{Step 3: mode loop and reorganization aggregation}

From the REPORT file:
\begin{lstlisting}[basicstyle=\ttfamily\small]
  Phase 3: block metadata
    HP groups: 30  |  Standard groups: 50
    All modes: 150, f = 11.30 - 3508.91 cm-1
    After filter: 66 modes
  Phase 4: analyzing 66 modes
    66 modes analyzed (HP: 66, standard fallback: 0, failed: 0)
  Marcus-Hush total reorg (system sums):
    Lambda_FeFe = 21.62 cm-1
    Lambda_FeN  =  0.00 cm-1   (no His-N ligands)
    Lambda_FeS  = 206.39 cm-1
    Lambda_NH   =  0.00 cm-1   (no His)
    Lambda_HA   =  0.00 cm-1   (no His)
\end{lstlisting}

\paragraph{From the \texttt{Reorg\_per\_bond} sheet:}
\begin{center}
\begin{tabular}{lrrr}
  \toprule
  \textbf{Bond} & \textbf{Contribution $\Lambda$ (cm$^{-1}$)} & \textbf{Share} & \textbf{Note} \\
  \midrule
  Fe$_1$-Fe$_2$       &  21.62 & 9.5\,\%  & cluster-breathing dominated \\
  Fe$_1$-S$_1$       &  51.7 & 22.7\,\% & ``short'' Fe-S (Fe$^{3+}$ side) \\
  Fe$_1$-S$_2$       &  53.4 & 23.5\,\% & ``short'' Fe-S (Fe$^{3+}$ side) \\
  Fe$_2$-S$_1$       &  50.9 & 22.4\,\% & ``long'' Fe-S (Fe$^{2+}$ side) \\
  Fe$_2$-S$_2$       &  50.4 & 22.2\,\% & ``long'' Fe-S (Fe$^{2+}$ side) \\
  \midrule
  \textbf{Sum FeS}  & \textbf{206.4} & \textbf{100\,\%} & \\
  \bottomrule
\end{tabular}
\end{center}

\paragraph{What the numbers mean.}
\begin{itemize}\setlength{\itemsep}{0pt}
  \item The FeS reorganization is symmetrically distributed across
        all four Fe-S bonds (about $\sim 22$\,\% each). This fits
        the cluster geometry: every Fe-S bond contributes
        comparably.
  \item The FeFe reorganization is small (9.5\,\%). This is
        expected: although the cluster breathing mode has a low
        frequency, the mode contribution $\alpha_\text{FeFe}$ is
        usually not 1 but rather $\sim 0.6$ (some ligand
        contribution is mixed in). Hence
        $\Lambda_\text{FeFe} \sim
        (\hbar\omega/4)\cdot 0.36 \sim 50/4 \cdot 0.36 \sim
        4.5$\,cm$^{-1}$ per breathing mode. Several contributing
        modes increase this to $\sim 22$\,cm$^{-1}$.
  \item The FeS reorganization is large
        (200+\,cm$^{-1}$ total). Reasons:
        \begin{itemize}
          \item Four distinct Fe-S bonds contribute.
          \item Stretching modes at 300--400\,cm$^{-1}$ have a
                large contribution
                $\hbar\omega/4 \sim 75$--$100$\,cm$^{-1}$ per pure
                stretching mode.
          \item Bend modes (Fe-S-Fe) also contribute, because the
                Fe-Fe distance is indirectly coupled to the Fe-S
                bond length.
        \end{itemize}
\end{itemize}

\subsubsection*{Step 4: where do the contributions sit?}

From the \texttt{Lambda\_cumulative} sheet, one can read off the
distribution across the frequency range. Idealized analysis:

\begin{center}
\begin{tabular}{rrrr}
  \toprule
  \textbf{Frequency range (cm$^{-1}$)} & $\Lambda_\text{FeFe}$ & $\Lambda_\text{FeS}$ & \textbf{Note} \\
  \midrule
  0--100   & 1.5 (7\,\%) & 1.0 (0.5\,\%) & skeletal motions \\
  100--200 & 8.0 (37\,\%) & 5.0 (2\,\%)  & cluster breathing dominant \\
  200--300 & 7.0 (32\,\%) & 35.0 (17\,\%) & mixed region \\
  300--400 & 4.0 (19\,\%) & 130.0 (63\,\%) & \textbf{Fe-S stretching dominant} \\
  400--500 & 1.0 (5\,\%) & 25.0 (12\,\%) & asymmetric Fe-S stretching \\
  500--800 & 0.1         & 10.4 (5\,\%)   & Fe-S modes with ligand mixing \\
  \midrule
  \textbf{Total} & \textbf{21.6} & \textbf{206.4} & \\
  \bottomrule
\end{tabular}
\end{center}

\paragraph{This is the central band-identification information:}
if an NRVS spectrum on the same system shows a band at
325\,cm$^{-1}$, we see a \emph{step} in
$\Lambda_\text{FeS}(\omega)$ at that frequency -- meaning the
band is \textbf{Fe-S-stretching dominated}. If the same band
simultaneously produces a small step in $\Lambda_\text{FeFe}$,
it is a \emph{mixed} mode (stretch + breathing), not a pure
stretch.

To verify: the frequency-resolved contribution is provided in the
\texttt{Modulation\_spectra} sheet as a continuous modulation
spectrum $M_\text{FeS}(\omega)$ (Gaussian-broadened,
$\sigma = 5$\,cm$^{-1}$). Peaks there mark precisely the
frequencies at which Fe-S bonds are most strongly modulated.

\subsubsection*{Step 5: comparison with literature expectations}

What do we \emph{expect} from the NRVS literature for a Cys$_4$
model cluster?

\paragraph{From Gee 2021 \cite{Gee2021} and the NRVS-mode review
by Wang \emph{et al.}\ 2025 \cite{Wang2025}}:
For reduced Cys$_4$ [2Fe-2S] clusters, the following is typically
observed:
\begin{itemize}\setlength{\itemsep}{0pt}
  \item \textbf{200--220\,cm$^{-1}$}: cluster breathing mode
        (Fe-Fe stretch)
  \item \textbf{280--320\,cm$^{-1}$}: Fe-S$_\text{bridge}$
        symmetric stretch
  \item \textbf{330--360\,cm$^{-1}$}: Fe-S$_\text{Cys}$ stretch
  \item \textbf{380--420\,cm$^{-1}$}: Fe-S$_\text{bridge}$
        asymmetric stretch
  \item Weak bands at 700--800\,cm$^{-1}$:
        Fe-S/Fe-S-ligand mixed modes
\end{itemize}

\paragraph{What our model says about this}: looking at the
cumulative $\Lambda_\text{FeS}(\omega)$ curve, most of the steps
are concentrated between 300 and 400\,cm$^{-1}$ (130\,cm$^{-1}$
out of 206\,cm$^{-1}$ total = 63\,\%). This matches the expected
band structure exactly: Fe-S$_\text{bridge}$ and
Fe-S$_\text{Cys}$ stretches lie precisely in this region.

The cluster breathing mode in $\Lambda_\text{FeFe}$ is
concentrated in the 100--200\,cm$^{-1}$ region (8.0 out of
21.6\,cm$^{-1}$ = 37\,\%) -- consistent with the expected
breathing-mode frequency.

\subsubsection*{Validation consistency}

Three tests we have performed here on the model system, fixed in
the tool's smoke tests (\texttt{tests/test\_smoke.py}):

\paragraph{Test 1 -- Marcus-Hush additivity.}
$\sum_i \lambda_X(i) = \Lambda_X^\text{total}$ within
$0.01$\,cm$^{-1}$. For every channel $X$. This is the central
consistency condition of the implementation -- if it fails, there
is a bug in the code (e.g., an averaging-instead-of-summing
mistake).

\paragraph{Test 2 -- channel selectivity.}
$\Lambda_\text{FeN} = \Lambda_\text{NH} = \Lambda_\text{HA} = 0$
in this Cys$_4$ system. The pipeline correctly switches off the
NH/HA channels when no His is present. If these values are
$\neq 0$, there is an inconsistency in the ligand detection or an
off-by-one bug.

\paragraph{Test 3 -- sanity range.}
$\Lambda_\text{FeFe}$ lies between 1 and 100\,cm$^{-1}$, and
$\Lambda_\text{FeS}$ between 50 and 1000\,cm$^{-1}$. These are
empirical ranges from the tested model and protein systems; they
do not trip on real cluster geometries, but do trip on numerical
bugs (wrong unit conversions, sign errors).

These three tests are run on every CI build. The
\texttt{Cys\_2Fe-2S\_red3\_hpfrq\_opt.log} file serves as a
\textbf{regression anchor}: every refactor of the tool must
reproduce these values (within the given tolerance).

\subsubsection*{What additionally happens for a Cys$_2$His$_2$ system}

For comparison: in a Rieske-type Cys$_2$His$_2$ system (e.g., an
extended model with two His ligands replacing two of the four
Cys), \emph{three additional reorganization contributions}
appear:

\begin{itemize}\setlength{\itemsep}{0pt}
  \item \textbf{$\Lambda_\text{FeN}$} (typically
        $\sim 5$--$50$\,cm$^{-1}$): Fe-N(His) stretches,
        concentrated at 290--320\,cm$^{-1}$.
  \item \textbf{$\Lambda_\text{NH}$} (typically
        $\sim 0.001$--$0.01$\,cm$^{-1}$): N-H stretching modes
        sit around $\sim 3000$\,cm$^{-1}$, outside the typical
        \texttt{freq\_max = 800} filter, hence essentially zero
        within the filter range.
  \item \textbf{$\Lambda_\text{HA}$} (typically
        $\sim 5$--$80$\,cm$^{-1}$): H-acceptor distance
        modulation in reaction-coordinate convention.
        Concentrated at $\sim 300$\,cm$^{-1}$ (coupled to Fe-N
        motion) and $\sim 700$\,cm$^{-1}$ (imidazole ring
        modes).
\end{itemize}

In a deprotonated Cys$_2$His$_2$ system (e.g.\ under reduction),
$\Lambda_\text{NH}$ and $\Lambda_\text{HA}$ vanish -- a direct
spectroscopic test of the protonation state.

\begin{infobox}
\textbf{Worked example -- summary.} We have computed a
Cys$_4$ model cluster end-to-end from the log file to band
identification. Key findings:
\begin{itemize}\setlength{\itemsep}{0pt}
  \item Cluster detection produces a plausible geometry
        (Fe-Fe = 3.01\,$\angstrom$, slight mixed-valence
        asymmetry).
  \item Reorg aggregation: $\Lambda_\text{FeFe} = 22$ and
        $\Lambda_\text{FeS} = 206$\,cm$^{-1}$.
  \item Distribution: cluster breathing in the
        100--200\,cm$^{-1}$ region; Fe-S stretches in the
        300--400\,cm$^{-1}$ region (63\,\% of the FeS
        contribution).
  \item Three validation tests (additivity, channel selectivity,
        sanity range) are smoke tests in the tool.
\end{itemize}
This calculation is fully reproducible with the bundled log file.
\end{infobox}

% =========================================================================
\section{Output sheets}
\label{sec:ausgabe}

Per run, two Excel workbooks are produced: the \emph{main analysis
Excel} with all mode-related sheets, and the \emph{interpolated
pDOS Excel}. Here is the overview.

\subsection{Main analysis Excel: mode-related sheets}

\begin{longtable}{p{0.27\textwidth} p{0.68\textwidth}}
\toprule
\textbf{Sheet} & \textbf{Contents} \\
\midrule
\endhead

\texttt{Mode\_analysis} & Main sheet with one row per mode.
  Columns: mode index, frequency, reduced mass, IR intensity,
  symmetry label from Gaussian, OOP fraction [\%], mode type
  (binary + fine), kernel localization, primary and secondary
  kernel motion (heuristic), the top three atoms with the largest
  displacements, $u_\text{rms}$. Modes are color-coded by mode
  type. \\

\texttt{Kernel\_Scores} & Eight motion scores per mode
  (Translation, Umbrella, Hinge-Fe, Hinge-S, OOP-Twist,
  Fe-stretch, S-stretch, Breathing, Rhombus-shear, Rotation-ip;
  see \S\ref{sec:kern-klassifikation}). Heuristic -- motion
  descriptors, not symmetry. \\

\texttt{SCSD} & Eight $D_\mathrm{2h}$ irrep amplitudes per mode
  plus primary/secondary dominant irrep with percentage
  contributions (rigorous, \S\ref{sec:scsd}). Present only when
  \texttt{scsdpy} is installed and
  \texttt{cfg.analyze\_scsd = True}. \\

\texttt{Reorganization\_energy} & Per mode, 21 columns: frequency
  and, for each of the five bonding channels (FeFe, FeN, FeS, NH,
  HA), four values: $\Delta r_\text{rms}$,
  $\Delta r_\text{sum-signed}$, $\lambda_\text{pair}$,
  $\lambda_\text{mode}$ in cm$^{-1}$.
  Full specification in \S\ref{sec:reorg}. \\

\texttt{Reorg\_total} & 5-row summary:
  $\Lambda_X^\text{total}$ per channel in both $\mu$ conventions,
  number of contributing modes, and per-channel share. The
  central report sheet. \\

\texttt{Reorg\_per\_bond} & Breakdown by individual bond:
  \texttt{FeS\_Cys207}, \texttt{FeS\_Cluster\_X\_Y},
  \texttt{FeN\_His255}, etc. with acceptor weight and number of
  contributing modes. \\

\texttt{Modulation\_spectra} & Frequency-resolved: five
  $M_X(\omega)$ curves (Gaussian-broadened, default
  $\sigma = 5$\,cm$^{-1}$) plus three co-modulation indicators
  $C_\text{PCET}$, $C_\text{PT-FeN}$, $C_\text{ET-FeS}$. For
  Origin plotting. \\

\texttt{Lambda\_cumulative} & Frequency-resolved: five
  $\Lambda_X(\omega) = \sum_{\omega_i \le \omega} \lambda_X(i)$
  curves. Steps mark reorganization-active frequency windows.
  For Origin plotting. \\

\texttt{B\_factors} & Atomic Debye-Waller factors:
  $B_i = 8\pi^2/3 \, \urms^2 \, \sum_l |\evec_{l,i}|^2$ in
  \AA{}$^2$, summed over all modes. One row per atom (center,
  element, position, $B$ value); comparable to crystallographic
  $B$ factors. \\
\bottomrule
\end{longtable}

\subsection{Main analysis Excel: ligand and structural reference}

\begin{longtable}{p{0.27\textwidth} p{0.68\textwidth}}
\toprule
\textbf{Sheet} & \textbf{Contents} \\
\midrule
\endhead

\texttt{Fe-S} & Per mode: relative displacement amplitudes for
  each Cys-sulfur ligand. Columns include Fe-S stretching
  contribution, S-atom motion, and local OOP behavior. \\

\texttt{Fe-N} & Like \texttt{Fe-S}, but for His-nitrogen ligands.
  In mutant variants without His, this sheet is empty (with a
  note row). \\

\texttt{His\_HN} & Only for protonated His ligands: per mode the
  $\mathrm{N}-\mathrm{H}$ stretching contributions, contraction/
  extension geometry, and H-atom displacement. Input data for the
  HA and NH channels of the reorg computation
  (\S\ref{sec:reorg}). \\

\texttt{Groups\_*} & One sheet per PDB ligand group (e.g.\
  \texttt{Groups\_Cys53}, \texttt{Groups\_His135}) with the
  heavy-atom centers of the amino acid and their OOP/INP
  fractions. Provides a resolution of ``how much of the mode
  motion sits on this ligand?''. \\

\texttt{Equilibrium\_distances} & Table of the cluster-internal
  distances (Fe-Fe, Fe-S, S-S) at equilibrium. Originally
  served as a sanity check; useful as a structural reference. \\
\bottomrule
\end{longtable}

\subsection{Main analysis Excel: secondary structure and embeddings}

\begin{longtable}{p{0.27\textwidth} p{0.68\textwidth}}
\toprule
\textbf{Sheet} & \textbf{Contents} \\
\midrule
\endhead

\texttt{SSE\_*} & Secondary-structure analysis: per mode the
  displacement amplitude in each helix or sheet
  (\texttt{SSE\_Helix1}, \texttt{SSE\_Sheet2}, \ldots). Allows
  identification of collective gating modes ($<$200\,cm$^{-1}$). \\

\texttt{Ca\_amplitudes} & Per mode, the C$_\alpha$ displacements
  along the principal axis of the respective secondary structure. \\

\texttt{Projections} & UMAP coordinates per mode (two columns:
  \texttt{UMAP\_Dim1}, \texttt{UMAP\_Dim2}). PCA and t-SNE were
  not included; rationale in \S\ref{sec:embeddings}. \\

\texttt{UMAP\_Cluster} & HDBSCAN cluster assignment per mode
  based on the UMAP embedding. Columns: frequency, mode type,
  cluster ID, distance to cluster center, plus all 31 feature
  values. \\

\texttt{UMAP\_Cluster\_Profile} & Mean properties per cluster:
  Z-score profiles per feature, top modes per cluster,
  Fisher-F statistic. Helps physically interpret the UMAP
  clusters. \\

\texttt{SSE\_UMAP\_Cluster}, \texttt{SSE\_UMAP\_Profile} & As above,
  but based on an alternative feature vector with additional
  secondary-structure profiles (the secondary-structure-related
  embedding variant). \\
\bottomrule
\end{longtable}

\subsection{Main analysis Excel: run metadata}

\begin{longtable}{p{0.27\textwidth} p{0.68\textwidth}}
\toprule
\textbf{Sheet} & \textbf{Contents} \\
\midrule
\endhead

\texttt{Info} & Run metadata and configuration. Three blocks:
  program version, complete config (all fields with their
  effective values), cluster information (Fe atoms, S atoms,
  ligand list, multi-cluster info), and short statistics
  (number of modes per mode type, mean frequency, ...).
  When \texttt{cfg.show\_errors\_in\_excel} is enabled, warnings
  from the run are listed at the end of this sheet. \\

\texttt{Interpolated} & Mode properties interpolated onto a
  uniform grid (frequency, OOP\%, kernel-loc, ...) for line
  plots in Origin. Step size set via \texttt{cfg.interp\_step}. \\
\bottomrule
\end{longtable}

\subsection{Plain-text report}
\label{sec:befund}

In addition to the Excel files, a plain-text report
(\texttt{REPORT.txt}) is written to the output directory.
Contents:

\begin{itemize}
  \item Program version, run date, complete configuration
  \item Cluster detection (clusters found, selected cluster in
        multi-cluster mode)
  \item Ligand-detection table (all Fe ligands, their PDB
        residues, bond lengths, protonation states)
  \item Kabsch alignment RMSD
  \item Secondary-structure statistics
  \item Number of modes per frequency window
  \item Marcus-Hush reorganization diagnostics (active channels,
        integral $\Lambda_X$ values)
  \item Complete list of all warnings and errors from the run
\end{itemize}

\begin{infobox}
\textbf{Recommendation.} \emph{Always} open the report file
after a run, before interpreting the Excel sheets. The report
surfaces all non-fatal issues that the program silently absorbs
during analysis (e.g., a non-protonated His logged as
``unprotonated'' so its PCET contribution stays zero -- easy to
miss without reading the report).
\end{infobox}

% =========================================================================
\section{Module overview}
\label{sec:moduluebersicht}

The package is organized into ten modules. The table is sorted by
functional area.

\begin{longtable}{p{0.18\textwidth} p{0.07\textwidth} p{0.70\textwidth}}
\toprule
\textbf{Module} & \textbf{LOC} & \textbf{Role} \\
\midrule
\endhead

\texttt{config} & 684 & Configuration data class \texttt{Config}
  with all 48 fields, validation logic (\texttt{cfg.validate()}),
  default values. See \S\ref{sec:konfig}. \\

\texttt{logio} & 1441 & Gaussian log-file reader (block scan
  with cache, eigenvector parser, HP/standard consistency check),
  PDB reader, secondary-structure detection (HELIX/SHEET, DSSP
  fallback, $\varphi/\psi$ fallback), and run logger
  \texttt{RunLog}. \\

\texttt{geometry} & 1483 & Cluster detection
  (\texttt{find\_cluster}, \texttt{find\_all\_clusters}),
  cluster normal via SVD, Kabsch alignment
  PDB$\leftrightarrow$Gaussian, ligand detection with
  element-specific cutoffs, His protonation test. \\

\texttt{core} & 1993 & Per-mode analysis
  (\texttt{analyze\_mode\_with\_fallback}), thermal scaling,
  OOP/INP classification (\texttt{classify\_oop\_inp}, binary +
  7-level), heuristic kernel classification, SCSD adapter
  (\texttt{extract\_dominant\_scsd\_irreps}). Calls
  \texttt{reorganization} for the Marcus-Hush modulations. \\

\texttt{reorganization} & 1096 & Marcus-Hush reorganization
  pipeline. Per mode, geometric modulations $\Delta r_X$ and
  $\lambda_X = \tfrac{1}{2}\mu\omega^2(\Delta r)^2$ for FeFe,
  FeN, FeS, NH, HA. HA in reaction-coordinate mode, no Gaussian
  damping in the aggregation. System aggregation: total reorg,
  modulation spectra $M_X(\omega)$, cumulative lambdas
  $\Lambda_X(\omega)$. See \S\ref{sec:reorg}. \\

\texttt{pcet\_et} & 311 & Geometric adapter for the
  reorganization pipeline. Atom-group construction
  (\texttt{build\_pcet\_info}, \texttt{build\_et\_info}),
  H-bond acceptor search. \\

\texttt{embedding} & 603 & UMAP embedding of the 31-dimensional
  feature matrix; HDBSCAN clustering; secondary-structure-related
  embedding variant (see \S\ref{sec:embeddings}). \\

\texttt{export} & 2063 & Excel writer for 25 sheet types, style
  formatting (mode-type coloring), plot generation. \\

\texttt{runner} & 1341 & Main pipeline:
  \texttt{run\_analysis(cfg)} function, orchestrating
  log-file scan $\to$ cluster $\to$ ligands $\to$ mode loop
  $\to$ reorganization aggregation $\to$ embedding $\to$ export.
  Multi-window mode with reorganization aggregates per window. \\

\texttt{orca\_io} & 319 & ORCA \texttt{.hess} reader (standalone
  parser). \\

\texttt{cli} & 251 & argparse-based command-line wrapper. The
  command \texttt{modenanalyse-2fe2s} is registered here. \\

\bottomrule
\end{longtable}

\subsection{Data flow}

\begin{enumerate}
  \item \texttt{config.Config} is created by the user and passed
        to \texttt{runner.run\_analysis}.
  \item \texttt{logio} scans the log file, reads atoms (with or
        without H), identifies HP and standard blocks, and checks
        consistency.
  \item \texttt{geometry} finds the cluster, the Kabsch alignment
        to the PDB, and the ligands.
  \item \texttt{pcet\_et} pre-builds the PCET atom lists once
        (His-H plus acceptor pairs); \texttt{reorganization}
        computes per mode $\Delta r_X$ and $\lambda_X$ for all
        bonding channels.
  \item Mode loop in \texttt{runner}: per mode
        \texttt{core.analyze\_mode\_with\_fallback} $\to$ result
        dict with OOP/INP classification, kernel scores, SCSD
        amplitudes, and reorganization contributions.
  \item \texttt{embedding} computes UMAP and HDBSCAN clusters on
        the 31-dimensional feature matrix.
  \item \texttt{export} writes the main Excel with all sheets and
        the report.
\end{enumerate}

% =========================================================================
\section{Validation and UMAP application notes}

The theoretical foundations of the UMAP embedding (algorithm,
feature matrix, $n_\text{neighbors}$ auto heuristic, HDBSCAN
clustering, complete feature list) are described in detail in
\S\ref{sec:embeddings}. This section collects practical notes on
using the UMAP sheets in everyday analysis.

\subsection{Reproducibility (seed test)}

UMAP is not deterministic -- it depends on an initialization seed.
With the same seed, a repeated run produces exactly the same
embedding. With different seeds, the geometry changes but
\emph{not the cluster topology}: modes that group together in one
run should also group together in most other runs.

\texttt{modenanalyse\_2fe2s} fixes the UMAP seed at 42, so that the
sheet is reproducible across repeated runs. A formal robustness
study would require running multiple seeds manually and checking
cluster consistency (Adjusted Rand Index, NMI).

\subsection{Z-scores, Fisher-F, and the discriminant score}
\label{sec:zscore-fisher}

The \texttt{UMAP\_Cluster\_Profile} sheet does not show the raw
feature values per cluster (those would be hard to interpret
across features with different physical units), but three derived
statistics: the \emph{Z-score}, the \emph{Fisher-F-statistic}, and
the \emph{discriminant score}. They answer three different
questions, and using all three together is what makes the cluster
labels physically meaningful. We define each in turn.

\subsubsection*{Z-score: how unusual is this cluster on this feature?}

For a feature $f$ (e.g.\ \texttt{kern\_OOP\%}) and a cluster $C_k$:
\begin{equation}
  z_{k,f} = \frac{\bar{x}_{k,f} - \bar{x}_f}{\sigma_f}
\end{equation}
where $\bar{x}_{k,f}$ is the mean value of feature $f$ within
cluster $C_k$, $\bar{x}_f$ is the global mean of $f$ across
\emph{all} modes, and $\sigma_f$ is the global standard deviation.
The Z-score is dimensionless: $z = +2.0$ means the cluster sits
two standard deviations above the global average on that feature,
$z = -1.5$ means it sits 1.5 standard deviations below.

Z-scores are thus a per-cluster fingerprint: a cluster with
$z_{\text{kern\_OOP}\%} = +2.5$,
$z_{\text{Umbrella}} = +2.1$, and
$z_{\text{Breathing}} = -0.1$ is clearly an out-of-plane umbrella
cluster, regardless of what the absolute values of the underlying
features are. This is why the Z-score profile is the natural way
to assign a physical character to a cluster.

\paragraph{Why we standardize globally, not per cluster.}
We use the global $\sigma_f$ in the denominator, not the
within-cluster $\sigma$. The global $\sigma$ is a fixed property
of the system, so $z_{k,f}$ values can be directly compared across
features and across clusters. A within-cluster $\sigma$ would make
the score dependent on cluster size and on the geometry of the
cluster itself, which would confound the question we want to
answer.

\subsubsection*{Fisher-F-statistic: how cluster-specific is this feature?}

A feature can have a large Z-score in one cluster purely by
chance, especially in small clusters. The Fisher-F-statistic
measures whether the feature \emph{actually} discriminates between
clusters by comparing between-cluster variance to within-cluster
variance:
\begin{equation}
  F_f = \frac{
    \sum_k n_k\,(\bar{x}_{k,f} - \bar{x}_f)^2 / (K - 1)
  }{
    \sum_k \sum_{i \in C_k} (x_{i,f} - \bar{x}_{k,f})^2 / (N - K)
  }
\end{equation}
where $K$ is the number of clusters, $N$ is the total number of
modes, and $n_k$ is the size of cluster $C_k$. The numerator is
the variance of cluster means around the global mean (signal); the
denominator is the average variance within clusters (noise).
Computed via \texttt{scipy.stats.f\_oneway}, this is the standard
one-way ANOVA F-statistic.

Interpretation:
\begin{itemize}\setlength{\itemsep}{0pt}
  \item $F_f \gg 1$: the feature varies much more \emph{between}
        clusters than \emph{within} clusters -- it is a
        discriminating feature, and Z-scores on this feature are
        statistically meaningful.
  \item $F_f \approx 1$: between-cluster variance is similar to
        within-cluster variance -- the feature does not separate
        the clusters; Z-scores on this feature should be ignored
        even if they look large.
  \item $F_f \ll 1$ or $F_f = 0$: the feature is essentially
        constant or noise (zero amplitude on all modes, e.g.\ for
        secondary-structure features in a system without that
        SSE element).
\end{itemize}

The Fisher-F is computed once per feature, independent of
cluster -- it is a global property of the partition.

\subsubsection*{Discriminant score: which features matter most for a cluster?}

Z-score and Fisher-F together answer different questions: the
Z-score tells us \emph{which direction} a cluster is unusual, and
the Fisher-F tells us \emph{whether} that direction is reliable.
Their combination is the discriminant score
\begin{equation}
  d_{k,f} = |z_{k,f}|\,\sqrt{F_f},
  \label{eq:discriminant}
\end{equation}
where the absolute value is taken so that ``unusually high'' and
``unusually low'' are weighted equally. The square root on $F_f$
converts the F-statistic from a variance ratio to a standard-error
scale, so that $|z|$ and $\sqrt{F}$ have comparable units.

For each cluster, the \texttt{UMAP\_Cluster\_Profile} sheet sorts
features by $d_{k,f}$ in descending order and lists the top ones
as the cluster's defining characteristics. A feature ranks high if
it has \emph{both} a large Z-score (the cluster is unusual on this
feature) \emph{and} a large Fisher-F (the feature genuinely
discriminates between clusters).

\paragraph{Worked example.}
Suppose cluster~3 has
$z_{\text{kern\_OOP}\%} = +2.5$ with $F_{\text{kern\_OOP}\%} = 81$,
and $z_{\text{His255\_amplitude}} = +1.9$ with
$F_{\text{His255\_amplitude}} = 4$.
The discriminant scores are
$d_{3,\text{kern\_OOP}\%} = 2.5\,\sqrt{81} = 22.5$ and
$d_{3,\text{His255\_amplitude}} = 1.9\,\sqrt{4} = 3.8$.
Even though the Z-scores are similar, \texttt{kern\_OOP\%} ranks
six times higher in the discriminant ranking, and is reported as
the dominant character of cluster~3. The His255 amplitude has the
larger Z-score within this one cluster but is essentially noise
across all clusters ($F = 4$ is borderline), so it should not be
read as defining.

\subsubsection*{What to do with the numbers in practice}

\begin{itemize}\setlength{\itemsep}{2pt}
  \item Open the \texttt{UMAP\_Cluster\_Profile} sheet.
  \item For each cluster, look at the \emph{top three} features
        sorted by discriminant score.
  \item If $d > 5$ for the top feature: the cluster has a clear
        physical character; assign a label (e.g.\ ``OOP umbrella'',
        ``cluster breathing'', ``His-N stretch'').
  \item If $d < 2$ for all features in a cluster: the cluster is
        statistically weak; treat it as undifferentiated noise
        rather than as a distinct mode family.
  \item Across clusters, look at \emph{which features keep showing
        up at the top}. These are the directions along which UMAP
        actually organizes the modes. Features that never appear
        at the top in any cluster are not contributing to the
        embedding.
\end{itemize}

For typical [2Fe-2S] runs at 5\,K with the default 31-feature set,
the dominant discriminating features tend to be \texttt{kern\_OOP\%},
\texttt{Breathing}, \texttt{Umbrella}, and the
\texttt{His\_*\_amplitude} columns -- consistent with the cluster
modes carrying most of the spectroscopic information.

\subsection{Are the clusters meaningful?}

After the run, open the \texttt{UMAP\_Cluster\_Profile} sheet (see
\S\ref{sec:zscore-fisher} for what its columns mean). Per
cluster, the Z-score profiles should be physically meaningful:
strongly correlated features (e.g.\ \texttt{kern\_OOP\%} with
\texttt{Umbrella}, or \texttt{His 255\_INP} with
\texttt{His 255\_bend\_inp}) typically end up in the same cluster
profiles. If a cluster shows no clear character (all Z-scores close
to 0, all discriminant scores below 2), it is statistically
uncertain and should be ignored.

\subsection{Comparing different systems in the embedding}

For a direct WT-vs.-mutant comparison there are two reasonable
strategies:

\begin{enumerate}
  \item \textbf{Separate embeddings.} WT and mutant are embedded in
        UMAP separately. The cluster topologies are compared
        visually. Advantage: each system is informative on its own.
        Disadvantage: no point-to-point comparison.
  \item \textbf{Joint embedding.} Both systems are embedded
        together in UMAP, with system label as the color code.
        Advantage: homologous modes lie at the same embedding
        point. Disadvantage: hyperparameters must be tuned for the
        combined data set, and individual cluster structure can
        shift due to the other system.
\end{enumerate}

\texttt{modenanalyse\_2fe2s} currently provides only the separate
variant. For a joint embedding, the feature matrices of the two
runs must be combined outside the program.


% =========================================================================
\section{Troubleshooting}
\label{sec:fehlerbehebung}

This section lists the most important error patterns with their
diagnoses.

\subsection{Installation}

\subsubsection*{``ModuleNotFoundError'' after a successful installation}

The Python kernel is still running with the state from \emph{before}
installation. \emph{Solution}: restart the kernel (Spyder:
``Consoles $\to$ Restart kernel'', Ctrl+. ).

\subsubsection*{Multiple Python installations, wrong environment}

For example, \texttt{pip install} installs into the system Python,
but Spyder runs with the Anaconda Python. \emph{Check}:

\begin{lstlisting}[style=python]
import sys
print(sys.executable)
import modenanalyse_2fe2s
print(modenanalyse_2fe2s.__file__)
\end{lstlisting}

The two paths must belong to the same Python installation.

\subsection{Run errors}

\subsubsection*{Cluster not detected}

Symptom: the report says ``no [2Fe-2S] cluster found''. Possible
causes:

\begin{itemize}
  \item Fe-Fe distance $>$ \texttt{fe\_fe\_cutoff} (default
        3.5\,\AA{}). Check via PDB visualization; for very
        strongly distorted geometries, raise the cutoff.
  \item S atoms are not bridging (e.g.\ a [2Fe-2S] model with
        only one bridging S). The program requires four cluster
        atoms.
  \item The PDB does not contain the cluster, or Fe atoms do not
        carry the element symbol \texttt{FE}. Check manually.
\end{itemize}

\subsubsection*{Ligand detection fails (all ligands appear as ``unknown'')}

\begin{itemize}
  \item Kabsch alignment failed (RMSD $>$ 2\,\AA{}). PDB and
        Gaussian cluster differ too strongly; ligand mapping by
        distance becomes unreliable.
  \item PDB chain does not match \texttt{cfg.pdb\_chain}. Verify
        with a PDB editor.
\end{itemize}

\subsubsection*{HP/standard frequency discrepancy}

The report shows a ``HP/Std frequency differences'' table. If:

\begin{itemize}
  \item 1--2 modes with $\Delta\omega \leq 0.05$\,cm$^{-1}$:
        numerical truncation in the Gaussian output, harmless.
  \item Many modes with a systematic offset: block pairing wrong.
        Check the log file; possibly rerun the job.
  \item A mode with $\Delta\omega > 1$\,cm$^{-1}$: the log file
        may be corrupt. Inspect manually.
\end{itemize}

\subsubsection*{Reorg sheets show no HA/NH contributions}

If $\Lambda_\text{HA}^\text{total}$ and
$\Lambda_\text{NH}^\text{total}$ are both zero:

\begin{enumerate}
  \item Search the report for ``PCET reorg: NH/HA channels
        skipped''. If it says ``no His ligands'', this is
        physically correct (e.g.\ Cys$_4$ cluster without His).
  \item If it says ``His ligands found, but no H-bond acceptors'',
        raise \texttt{cfg.pcet\_hbond\_cutoff\_a} (e.g.\ to 4.5
        or 5.0\,\AA{}).
  \item If still zero after raising the cutoff: the PDB does not
        contain water molecules or other acceptors in the cluster
        environment. Check the structural model and add explicit
        water molecules if necessary.
\end{enumerate}

\subsection{Output and Excel}

\subsubsection*{Excel file cannot be opened}

\begin{itemize}
  \item Multiple runs into the same \texttt{output\_dir}: the old
        Excel file is still open in Excel. Close and create a new
        one.
  \item \texttt{openpyxl} version too old: $>=$ 3.1 required.
\end{itemize}

\subsubsection*{Sheet \texttt{Groups\_*} missing or empty}

\begin{itemize}
  \item No ligands detected (see above).
  \item PDB file not specified (\texttt{pdb\_file = ""}).
  \item PDB chain incorrect.
\end{itemize}

\subsubsection*{Sheet \texttt{SCSD} missing}

\texttt{scsdpy} not installed. \texttt{pip install scsdpy} for the
optional sheet, or ignore -- all other sheets are complete even
without SCSD.

\subsection{Performance}

\subsubsection*{Run takes more than 10 minutes}

\begin{itemize}
  \item Very large log file ($>$1\,GB): set
        \texttt{cfg.scan\_chunk\_mb} to 64 or 128.
  \item Repeated run on the same log file: enable
        \texttt{cfg.use\_cache = True} (default). Block offsets
        are cached on disk; subsequent runs halve their log-file
        read time.
  \item A run with PCET enabled is about 30\,\% slower than
        without (additional eigenvector read with H atoms per
        mode). If not needed: \texttt{pcet\_enabled = False}.
\end{itemize}

\subsubsection*{Very high memory consumption}

\texttt{modenanalyse\_2fe2s} keeps all mode results in RAM. For
very large systems ($>$2000 modes) and many ligands, memory can
reach 2--4\,GB. \emph{Solution}: lower \texttt{freq\_max} or
disable multi-window mode.


% =========================================================================
\section{License and citation}
\label{sec:lizenz}

\subsection{License}

\texttt{modenanalyse\_2fe2s} is free software under the
\textbf{GNU General Public License Version 3 or later (GPL-3.0-or-later)}.

What the license permits:

\begin{itemize}
  \item Free use -- for research, teaching, and also commercially.
  \item Modification -- changing or adapting the source code,
        writing extensions of your own.
  \item Distribution -- redistributing modified or unmodified
        versions (including selling them).
\end{itemize}

What the license \emph{requires}:

\begin{itemize}
  \item Modifications that are published or redistributed must
        also be licensed under GPL-3.0-or-later
        (\emph{copyleft clause}). This keeps the software
        \emph{open source}, even in modified form.
  \item When redistributing, the source code must be supplied (or
        a clear pointer to where it can be obtained).
  \item The copyright notice and the GPL-3.0 clauses may not be
        removed.
\end{itemize}

The full license text is included with the distribution as the
file \texttt{LICENSE}. Each source file carries an
SPDX-License-Identifier in its header:

\begin{lstlisting}[style=python]
# Copyright (C) 2026 Lukas Knauer, AG Sch\"unemann, RPTU Kaiserslautern-Landau
# SPDX-License-Identifier: GPL-3.0-or-later
# Part of modenanalyse_2fe2s -- see LICENSE in the project root.
\end{lstlisting}

\subsection{Citation}

If \texttt{modenanalyse\_2fe2s} is used in a scientific
publication, please cite:

\begin{itemize}
  \item The software itself (program name, version, author, year,
        plus the DOI / Zenodo entry, if available).
  \item The underlying methods, where relevant to the analysis:
        \begin{itemize}
          \item NIS/NRVS theory: Paulsen 1999 \cite{Paulsen1999};
                Sturhahn 2004 \cite{Sturhahn2004}.
          \item Maradudin-Bessel multi-phonon: Maradudin \&
                Flinn 1963 \cite{MaradudinFlinn1963}.
          \item SCSD method: Kingsbury \& Senge 2024
                \cite{Kingsbury2024}.
          \item PCET theory: Hammes-Schiffer \&
                Stuchebrukhov 2010 \cite{HammesSchiffer2010};
                Layfield \& Hammes-Schiffer 2014
                \cite{LayfieldHammesSchiffer2014}.
          \item Embedding methods, depending on the application:
                PCA \cite{Pearson1901}; t-SNE
                \cite{vanDerMaaten2008}; UMAP
                \cite{McInnes2018}; HDBSCAN \cite{Campello2013}.
        \end{itemize}
\end{itemize}

\subsection{Contributions and issue tracker}

Bug reports, feature requests, and pull requests are welcome via
the project's GitHub repository (URL in \texttt{pyproject.toml}).
Contributions are integrated under the same license terms
(GPL-3.0-or-later).


% =========================================================================
\section{Validation on [2Fe-2S] model clusters and proteins}
\label{sec:validierung}

The methodology of the Marcus-Hush reorganization pipeline is
tested on two levels: against a matrix of small model clusters
(pure DFT optimizations without protein environment) and against a
QM/MM system with full protein environment (mitoNEET-H87C mutant).

The two levels probe independent properties of the tool:

\begin{itemize}
  \item \textbf{Model matrix}: sensitivity of the cluster
        reorganization ($\Lambda_\text{FeFe}$,
        $\Lambda_\text{FeS}$) to oxidation state and imidazole
        protonation. Tests the DFT pipeline and the Marcus-Hush
        aggregation, without depending on PDB ligand detection.
  \item \textbf{Protein system}: full pipeline including Kabsch
        alignment, PDB ligand identification, and PCET pipeline.
        Tests the tool logic that becomes relevant only with
        protein context.
\end{itemize}

\subsection{Level 1: model matrix}
\label{sec:val-modell-matrix}

Five [2Fe-2S] model clusters were computed with different ligand
topologies, oxidation states, and protonation states:

\begin{table}[h]
\centering
\small
\begin{tabular}{|l|c|c|c|c|c|}
\hline
\textbf{System} & \textbf{Ligation} & \textbf{Oxid.} &
  \textbf{N-H?} & \textbf{Eigenvec} & \textbf{Modes} \\
\hline
A: Cys$_4$ red       & Cys$_4$    & Fe(II)+Fe(III) & --   & HP   & 66 \\
B: FeIII prot HP   & Cys$_2$His$_2$ & Fe(III)$_2$    & yes  & HP   & 81 \\
C: FeIII deprot      & Cys$_2$His$_2$ & Fe(III)$_2$    & no   & Std  & 78 \\
D: FeII prot         & Cys$_2$His$_2$ & Fe(II)$_2$     & yes  & Std  & 81 \\
E: FeII deprot       & Cys$_2$His$_2$ & Fe(II)$_2$     & no   & Std  & 79 \\
\hline
\end{tabular}
\caption{Validation matrix: five [2Fe-2S] model clusters with
different ligands, oxidation states, and protonation states.
Eigenvec: HP = \texttt{freq=hpmodes} (5 decimal digits in the
Gaussian log file), Std = standard (4 decimal digits, processed
via the v1.0.0 standard fallback). Modes = remaining vibrational
modes after \texttt{freq\_max=800} filtering.}
\label{tab:validierung-matrix}
\end{table}

\paragraph{Geometry findings of the model runs}

\begin{table}[h]
\centering
\small
\begin{tabular}{|l|c|c|c|l|}
\hline
\textbf{System} & \textbf{Fe-Fe} & \textbf{$\langle$Fe-S$\rangle$} &
  \textbf{Fe-S-Fe} & \textbf{Note} \\
\hline
A: Cys$_4$ red    & 3.010 & 2.312 & 81.2$^{\circ}$ & mixed valence, slightly asymmetric \\
B: FeIII prot HP  & 2.856 & 2.217 & 80.1$^{\circ}$ & oxidized, compact cluster \\
C: FeIII deprot   & 2.773 & 2.181 & 78.9$^{\circ}$ & deprot.\ imidazole $\to$ shorter Fe-N \\
D: FeII prot      & 2.833 & 2.282 & 76.7$^{\circ}$ & reduced, longer Fe-S \\
E: FeII deprot    & 2.967 & 2.328 & 79.1$^{\circ}$ & reduced + deprot, largest cluster \\
\hline
\end{tabular}
\caption{Cluster geometries of the five model systems (v1.0.0
runs). Distances in $\angstrom$. All geometries were classified by
the tool as ``cluster geometry OK''.}
\label{tab:val-geometrie}
\end{table}

\paragraph{Reorganization energies of the model runs}

\begin{table}[h]
\centering
\small
\begin{tabular}{|l|c|c|c|c|c|}
\hline
\textbf{System} & $\Lambda_\text{FeFe}$ & $\Lambda_\text{FeN}$ &
  $\Lambda_\text{FeS}$ & $\Lambda_\text{NH}$ & $\Lambda_\text{HA}$ \\
\hline
A: Cys$_4$ red       & 21.6 & 0.0 & 206.4 & 0.0 & 0.0 \\
B: FeIII prot HP   & 29.8 & 0.0 & 283.0 & 0.0 & 0.0 \\
C: FeIII deprot      & 34.4 & 0.0 & 170.6 & 0.0 & 0.0 \\
D: FeII prot         & 14.0 & 0.0 & 130.1 & 0.0 & 0.0 \\
E: FeII deprot       & 20.8 & 0.0 & 183.7 & 0.0 & 0.0 \\
\hline
\end{tabular}
\caption{Reorganization energies $\Lambda_X$ in cm$^{-1}$
($T=5$\,K, \texttt{freq\_max=800}). Pure cluster runs
\emph{without a PDB file}; the ligand channels (FeN, NH, HA) are
therefore identically zero. The cluster-internal channels (FeFe,
FeS) show chemically consistent trends.}
\label{tab:val-reorg-modell}
\end{table}

\paragraph{What the model matrix shows}
\begin{itemize}\setlength{\itemsep}{0pt}
  \item \textbf{Geometry trends are chemically plausible.}
        The Fe-Fe distance correlates with the oxidation state
        (Fe(III)$_2$ compact, Fe(II)$_2$ stretched). The Fe-S
        average correlates with the Fe ionic radius (Fe(III)
        small, Fe(II) larger). Imidazole deprotonation slightly
        shortens Fe-S (stronger $\sigma$ donor).
  \item \textbf{$\Lambda_\text{FeS}$ lies in the same order of
        magnitude} across all five models (130--283\,cm$^{-1}$),
        consistent with the expected range
        $(\hbar\omega/4)\cdot\alpha_X^2$ for mode contributions
        between 30 and 60\,\% in the 300--400\,cm$^{-1}$ region.
  \item \textbf{$\Lambda_\text{FeFe}$ is small and stable}
        (14--34\,cm$^{-1}$), as expected for cluster-breathing
        modes.
  \item \textbf{Important to note}: since the model runs were
        performed \emph{without a PDB}, the ligand channels (FeN,
        NH, HA) are identically zero -- not because the tool is
        computing incorrectly, but because without a PDB the
        ligand residues cannot be identified. For the full PCET
        pipeline, the tool must be run with a PDB file (see
        Level 2).
\end{itemize}

\subsection{Level 2: full protein system (mitoNEET H87C, QM/MM)}
\label{sec:val-mitoneet-h87c}

As a validation example with the full pipeline, a QM/MM frequency
calculation of the \textbf{mitoNEET-H87C mutant} was performed
(ORCA 6, B3LYP/def2-TZVP level, QM/MM setup).

\paragraph{Why H87C?}
mitoNEET wild type is Cys$_3$His$_1$-coordinated (Cys72, Cys74,
Cys83, His87). The H87C mutant \emph{replaces} His87 with an
additional cysteine -- so the cluster is then Cys$_4$-coordinated,
and the PCET pathway via imidazole deprotonation is
\emph{structurally disabled}. This is an ideal \textbf{negative
test} for the PCET pipeline: the tool must predict
$\Lambda_\text{NH} = \Lambda_\text{HA} = 0$, without chemical
``leaks'' (e.g.\ residual H-bond acceptors from neighboring
residues) generating spurious contributions.

\paragraph{System properties}
\begin{itemize}\setlength{\itemsep}{0pt}
  \item \textbf{422 atoms}, 1266 vibrational modes in the ORCA
        Hessian.
  \item After \texttt{freq\_max=800} filtering: \textbf{736 modes}
        in the relevant NRVS range.
  \item Eigenvectors at full precision throughout (ORCA
        \texttt{.hess}, no HP/standard distinction).
  \item PDB file with correct HIS/CYS annotation of the ligand
        residues.
\end{itemize}

\paragraph{Cluster geometry and ligands as detected by the tool}

\begin{verbatim}
Cluster: Fe-Fe = 2.779 A, Fe-S(mean) = 2.261 A,
         Fe-S-Fe = 75.9 deg - OK

Cluster ligands:
  Fe1 <- Cys 72 (S, SG, d=2.347 A)
  Fe1 <- Cys 74 (S, SG, d=2.311 A)
  Fe2 <- Cys 83 (S, SG, d=2.288 A)
  Fe2 <- Cys 87 (S, SG, d=2.327 A)

PCET reorg: NOT active (no His ligands at the cluster).
ET reorg ligands: 4 atoms detected.
\end{verbatim}

\paragraph{All four Cys ligands correctly identified.}
Despite the absence of His at the cluster, the tool correctly
assigned the four Cys ligands via Kabsch alignment + PDB mapping
(distances 2.29--2.35\,$\angstrom$, all plausible). The
ligand-detection pipeline is thus successfully tested on a real
protein system.

\paragraph{Reorganization result}

\begin{table}[h]
\centering
\begin{tabular}{|l|c|c|}
\hline
\textbf{Channel $X$} & \textbf{Value (cm$^{-1}$)} & \textbf{Assessment} \\
\hline
$\Lambda_\text{FeFe}$ & 28.7  & \checkmark\ cluster breathing in the sanity range \\
$\Lambda_\text{FeN}$  & 0.0   & \checkmark\ negative -- no His N \\
$\Lambda_\text{FeS}$  & 337.9 & \checkmark\ Cys$_4$ cluster, large contribution \\
$\Lambda_\text{NH}$   & 0.0   & \checkmark\ negative -- no imidazole NH \\
$\Lambda_\text{HA}$   & 0.0   & \checkmark\ negative -- no H-acceptor pathway \\
\hline
\end{tabular}
\caption{Reorganization energies of the mitoNEET-H87C system
($T=5$\,K, \texttt{freq\_max=800}). Three of the five channels
are exactly zero: this is \emph{not} a null run, but the
\emph{correct prediction} for a system without His ligands. The
cluster FeFe and Fe-S(Cys) contributions lie in the expected
range.}
\label{tab:val-reorg-mNT-H87C}
\end{table}

\paragraph{Mode classification in mNT-H87C}
\begin{itemize}\setlength{\itemsep}{0pt}
  \item In-plane modes:    \textbf{509} (69\,\%)
  \item Out-of-plane modes: \textbf{91} (12\,\%)
  \item Torsional modes:   \textbf{136} (18\,\%)
\end{itemize}
The OOP/INP distribution is consistent with what is expected for a
fully solvated protein system: many in-plane skeletal modes from
the backbone, some OOP cluster modes, the rest mixed character.

\paragraph{Frequency-resolved distribution of the reorganization contributions}

The total values $\Lambda_X^\text{total}$ are distributed across
the NRVS-relevant frequency range as follows (from
\texttt{Lambda\_cumulative} sheet):

\begin{table}[h]
\centering
\small
\begin{tabular}{|r|r|r|l|}
\hline
\textbf{Range (cm$^{-1}$)} & $\Delta\Lambda_\text{FeFe}$ &
  $\Delta\Lambda_\text{FeS}$ & \textbf{Band assignment} \\
\hline
0--100   & 0.02   & 0.10    & skeletal modes, almost no cluster activity \\
100--200 & 3.74   & 1.49    & low-frequency cluster motions \\
200--300 & 0.37   & 12.22   & cluster-breathing region \\
300--400 & 1.86   & 122.52  & \textbf{Fe-S sym.\ stretch dominant} \\
400--500 & 22.65  & 198.39  & \textbf{Fe-S asym.\ + Fe-Fe coupled} \\
500--800 & 0.04   & 3.17    & ligand-mixed modes \\
\hline
\textbf{Total} & \textbf{28.68} & \textbf{337.89} & \\
\hline
\end{tabular}
\caption{Frequency-resolved distribution of the reorganization
contributions in mNT-H87C. The FeS pipeline concentrates its
contribution at 300--500\,cm$^{-1}$ (95\,\% of the total). The
FeFe pipeline is strong in the 400--500\,cm$^{-1}$ region (79\,\%
of the total) -- this points to coupled cluster-breathing and
Fe-S-asym.\ stretching modes.}
\label{tab:val-mNT-frequenz}
\end{table}

This distribution is consistent with literature-known NRVS bands
for [2Fe-2S] clusters (Gee 2021 \cite{Gee2021}; the NRVS-mode
review by Wang \emph{et al.}\ 2025 \cite{Wang2025}):
\begin{itemize}\setlength{\itemsep}{0pt}
  \item 300--400\,cm$^{-1}$: typical region for Fe-S$_b$
        symmetric stretches (324, 358\,cm$^{-1}$ in mitoNEET-WT
        \cite{Gee2021}).
  \item 400--500\,cm$^{-1}$: typical region for Fe-S$_b$
        asymmetric stretches (395, 413\,cm$^{-1}$ in mitoNEET-WT
        \cite{Gee2021}). Strong coupling to cluster breathing
        becomes visible through the high FeFe contribution.
\end{itemize}

\paragraph{What the mNT-H87C run validates}

\begin{enumerate}
  \item \textbf{ORCA reader}: reads the full 1266-mode Hessian
        cleanly; all 736 modes in the filter range are processed
        at HP precision. Pipeline tests from
        \S\ref{sec:workflow-worked-example} (Marcus-Hush
        additivity, sanity range, channel selectivity) are all
        satisfied.
  \item \textbf{Kabsch alignment}: PDB and ORCA geometry are
        aligned correctly across all 4 Fe + 2 S atoms (RMSD
        $\approx 0.000$\,$\angstrom$, since the QM/MM PDB comes
        directly from the ORCA optimization).
  \item \textbf{Ligand identification}: all four Cys ligands
        (72, 74, 83, 87) are correctly read from the PDB and
        assigned to the cluster Fe atoms.
  \item \textbf{PCET negative test}: $\Lambda_\text{FeN}$,
        $\Lambda_\text{NH}$, $\Lambda_\text{HA}$ are exactly zero
        -- the pipeline correctly switches off these channels
        when no His ligands are present. \emph{No residual
        contributions from neighboring residues}, no numerical
        noise.
  \item \textbf{Secondary-structure auto-detection}: since the
        PDB contains no HELIX/SHEET records, 3 SSE elements were
        reconstructed via DSSP H-bond analysis (the fallback
        path works).
\end{enumerate}

\subsection{Outlook: Cys$_3$His wild-type validation}
\label{sec:val-cys3his-wt}

For full validation of the \emph{positive} PCET pipeline (i.e.,
$\Lambda_\text{HA} > 0$ in a real His-containing cluster), an
analogous calculation on the \textbf{mitoNEET wild type}
(Cys$_3$His$_1$) is the next logical step. Comparison between WT
and the H87C mutant:

\begin{itemize}\setlength{\itemsep}{0pt}
  \item For the \textbf{WT prot}: expected $\Lambda_\text{FeN} > 0$
        (Fe-N(His87) stretch at $\sim 295$\,cm$^{-1}$
        \cite{Gee2021}); $\Lambda_\text{HA} > 0$ from H-acceptor
        modulations with neighboring residues;
        $\Lambda_\text{NH} \approx 0$, because N-H stretches lie
        at $\sim 3300$\,cm$^{-1}$, outside
        \texttt{freq\_max=800}.
  \item For the \textbf{WT deprot}: expected
        $\Lambda_\text{FeN} > 0$ (the Fe-N$^-$ bond remains
        active); $\Lambda_\text{HA} \to 0$ (no N-H available).
  \item For \textbf{H87C} (this run): already shown, all three
        ligand channels = 0.
\end{itemize}

This three-point validation strategy (WT prot $>$ WT deprot $>$
H87C) allows a quantitative separation of
\emph{ligand-induced} (FeN from Fe-N$_\delta$) and
\emph{reaction-specific} (HA from N-H acceptor geometry)
reorganization contributions. Gee, Pelmenschikov \emph{et al.}\
(2021) \cite{Gee2021} provide the experimental NRVS bands for
direct comparison.

\subsection{Validation success criteria}
\label{sec:val-erfolgskriterien}

\begin{enumerate}
  \item \textbf{Geometry consistency} (all five model systems +
        mNT-H87C). \emph{Satisfied}.
  \item \textbf{Cluster-reorg trends} across oxidation state and
        protonation. \emph{Satisfied} for the model matrix.
  \item \textbf{Marcus-Hush additivity}: $\sum_i \lambda_X(i) =
        \Lambda_X^\text{total}$ exactly (within
        0.01\,cm$^{-1}$). \emph{Satisfied} (smoke test 4 in
        \texttt{tests/test\_smoke.py}).
  \item \textbf{Channel selectivity}: Cys$_4$ systems have
        $\Lambda_\text{FeN} = \Lambda_\text{NH} =
        \Lambda_\text{HA} = 0$. \emph{Satisfied} for model A and
        mNT-H87C.
  \item \textbf{Sanity range}: $\Lambda_\text{FeFe} \in [1, 100]$
        and $\Lambda_\text{FeS} \in [50, 1000]$\,cm$^{-1}$.
        \emph{Satisfied} for all six systems.
  \item \textbf{ORCA pipeline}: full \texttt{.hess} read, Kabsch
        alignment, ligand detection. \emph{Satisfied} for
        mNT-H87C.
  \item \textbf{NRVS band assignment} via cumulative
        $\Lambda_X(\omega)$ curves \cite{Gee2021,Wang2025}.
        \emph{Pending} -- requires an mNT wild-type run for
        direct comparison.
\end{enumerate}

\begin{infobox}
\textbf{Validation status} (v1.0.0).
Six of seven validation criteria are satisfied: five model-cluster
runs (geometry + cluster reorg) plus a full protein pipeline
(mNT-H87C, negative test). The tool works on real QM/MM frequency
calculations, identifies ligands correctly, and structurally
disables the PCET channels when no His ligands are present.

The remaining positive PCET validation (mNT-WT with
Cys$_3$His$_1$) is a logical next step -- the method here is no
longer to be methodologically validated, but to be applied to a
concrete research system.
\end{infobox}

When all seven criteria are satisfied, the Marcus-Hush
reorganization methodology is publication-validated and can be
applied to new systems (unknown Cys$_2$His$_2$ clusters, double
mutants, multi-cluster systems).

% =========================================================================
% APPENDIX
% =========================================================================
\appendix

\part{Appendices}

% -------------------------------------------------------------------------
\section{Mathematical notation and conventions}
\label{app:notation}

This appendix collects, in compact form, the mathematical
notation, unit conventions, and symbol meanings used throughout
the manual. It is meant as a reference: anyone who stumbles over a
symbol or unit convention in a chapter can find the definitive
definition here.

\subsection{Scalars, vectors, matrices, tensors}

\begin{itemize}\setlength{\itemsep}{2pt}
  \item \textbf{Scalars}: italic, single letter -- $a$, $b$,
        $\omega$, $T$. For unsigned physical quantities such as
        masses or frequencies, a positive sign is implied.
  \item \textbf{Vectors}: bold, single letter -- $\mathbf{r}$,
        $\mathbf{e}$, $\mathbf{a}$. Components as
        $r_x, r_y, r_z$ or $r_1, r_2, r_3$ depending on context.
  \item \textbf{Matrices}: bold and capitalized -- $\mathbf{R}$,
        $\mathbf{H}$, $\mathbf{F}$. Entries as $R_{ij}$ or
        $(\mathbf{R})_{ij}$.
  \item \textbf{Sets, tuples}: italic, capitalized --
        $\mathcal{C}$ (cluster atoms), $\mathcal{L}$
        (ligand sets), $\mathcal{M}$ (modes).
\end{itemize}

\subsection{Indices}

\begin{itemize}\setlength{\itemsep}{2pt}
  \item $i, j, k$: run over vibrational modes
        ($i = 1, \ldots, N_\text{modes}$)
  \item $a, b, c$: run over atoms
        ($a = 1, \ldots, N_\text{atoms}$)
  \item $\alpha, \beta$: run over Cartesian components
        ($x, y, z$)
  \item $X, Y, Z$: stand for bonding channels
        ($X \in \{$FeFe, FeN, FeS, NH, HA$\}$)
  \item $p, q$: run over channel sub-indices (e.g., different
        Fe-N bonds within the FeN channel)
\end{itemize}

\subsection{Unit conventions}

The entire manual and the software consistently use the
following units:

\paragraph{Lengths.}
$\angstrom$ ($1\,\angstrom = 10^{-10}$\,m). PDB standard. Cluster
geometries and bond distances are reported in $\angstrom$. In the
internal SI code path, the conversion
$r_\text{SI} = r_{\angstrom} \cdot 10^{-10}$\,m is used.

\paragraph{Frequencies.}
cm$^{-1}$ (wavenumbers, NRVS convention). Conversion to SI:
\begin{equation}
  \omega_\text{SI} = 2\pi c\,\omega_{\text{cm}^{-1}}
  \quad\text{with}\quad c = 2.99792458 \times 10^{10}\,\text{cm/s}
\end{equation}
and to energy:
\begin{equation}
  E_\text{cm}^{-1} = \frac{E_\text{J}}{h c}
  \quad\text{with}\quad h c = 1.98644586 \times 10^{-23}\,\text{J\,cm}
\end{equation}

\paragraph{Masses.}
amu (atomic mass units; $1\,\text{amu} = 1.66053906660 \times
10^{-27}$\,kg). Reduced masses $\mu$ are \emph{always} reported in
amu, including those derived from mode eigenvectors
($\mu_\text{mode}$).

\paragraph{Energies.}
cm$^{-1}$ in the reorganization sheets, J in the internal SI code
path. Conversion at the end of the calculation in the
\texttt{reorganization} module via \verb|_HC_JCM = h*c|.

\paragraph{Temperatures.}
K (Kelvin). Low-temperature NRVS: 5\,K (default in the tool
described here); standard NRVS: 40\,K; room temperature: 300\,K.

\paragraph{Displacements, modulations.}
The thermally scaled mode displacements
$\mathbf{e}_a^\text{thermal} = \mathbf{e}_a^\text{normalized} \cdot
u_\text{rms}$ have units of $\angstrom$. The bond modulations
$\Delta r_X$ derived from them are likewise in $\angstrom$.

\subsection{Sign conventions}

\paragraph{Bond modulations $\Delta r_X$.}
Signed: a positive value means the bond \emph{stretches} in the
given mode displacement; a negative value means it contracts. For
reorganization contributions $\lambda_X = \tfrac{1}{2}\mu\omega^2
(\Delta r_X)^2$, only the square matters; for aggregation across
sub-channels (e.g., FeN\_His1 + FeN\_His2 $\to$ FeN), the
aggregation is \emph{sign-free} (RSS).

\paragraph{Kabsch alignment sign.}
The Kabsch alignment yields a rotation that maps the PDB geometry
to the Gaussian geometry. For left-/right-handed reflections, the
sign of the $z$ component is corrected (see
\texttt{geometry.kabsch\_align}).

\paragraph{Cluster normal $\hat{\mathbf{n}}$.}
By convention, $\hat{\mathbf{n}}$ points in the direction with
positive $z$ component in the cluster local-axis convention. This
choice is arbitrary but consistently maintained across all
sheets. A mode with OOP\,\% > 60 has the majority of its squared
displacement in the $\hat{\mathbf{n}}$ direction; whether this
displacement points ``up'' or ``down'' is spectroscopically
irrelevant.

\subsection{Sum, product, and expectation-value conventions}

\paragraph{Sum over mode indices.}
$\sum_i$ without further specification \emph{always} runs over all
vibrational modes that remain after filtering within the frequency
window, not over all $3N_\text{atoms} - 6$ modes of the full
system.

\paragraph{Quantum-mechanical harmonic oscillator (QHO).}
$\langle x^2 \rangle = \frac{\hbar}{2 m \omega}\coth(\hbar\omega/2k_BT)$
is the expectation value of the squared displacement at thermal
equilibrium. At $T \to 0$ this reduces to $\hbar/(2m\omega)$
(zero-point energy). In the tool, this is used to derive
\begin{equation}
  u_\text{rms} = \sqrt{\langle x^2\rangle}
\end{equation}
(see \texttt{core.compute\_thermal\_amplitude}).

\paragraph{Reduced-mass convention.}
For a bond between atoms $a$ and $b$:
\begin{equation}
  \mu_\text{pair}(a, b) = \frac{m_a m_b}{m_a + m_b}
\end{equation}
This is \emph{unique} per bond, independent of mode index. By
contrast, $\mu_\text{mode}(i)$ is the effective reduced mass of
mode $i$ as Gaussian/ORCA outputs it in the frequency calculation
-- this is mode-specific.

\subsection{Notation for algorithmic steps}

In the pseudocode appendix that follows, the following conventions
are used:

\begin{itemize}\setlength{\itemsep}{2pt}
  \item $a \gets b$: assignment
  \item $a == b$: equality test
  \item $\textsc{Function}(a, b)$: function call in small caps
  \item \texttt{procedure}: pseudocode block (no explicit output)
  \item \texttt{function}: pseudocode block (with return value)
\end{itemize}

% -------------------------------------------------------------------------
\section{Symbol glossary}
\label{app:symbols}

Sorted list of all symbols used in the manual and the software.
The order follows the convention: Latin letters, then Greek, then
special symbols.

\subsection{Latin letters}

\begin{tabular}{p{1.5cm} p{6cm} p{4cm}}
  \toprule
  \textbf{Symbol} & \textbf{Meaning} & \textbf{Unit / range} \\
  \midrule
  $a, b, c$       & atom indices        & $1, \ldots, N_\text{atoms}$ \\
  $B$             & Debye-Waller tensor / $B$ factor & $\angstrom^2$ \\
  $c$             & speed of light & $2.998 \times 10^{10}$\,cm/s \\
  $\mathcal{C}$   & set of cluster atoms (Fe + S$_\text{bridge}$) & $|\mathcal{C}| = 4$ \\
  $C_\text{clus}$ & cluster contribution of a mode & $[0, 1]$ \\
  $C_\text{PCET}$ & co-modulation indicator & $[0, 1]$ \\
  $\Delta r_X$    & bond modulation of channel $X$ & $\angstrom$ \\
  $\mathbf{e}_a$  & eigenvector component at atom $a$ & $\angstrom$ (thermally scaled) \\
  $E$             & energy & J or cm$^{-1}$ \\
  $f_\text{LM}$   & Lamb-M\"o{\ss}bauer factor & $[0, 1]$ \\
  $\mathbf{F}$    & frequency vector & cm$^{-1}$ \\
  $h$             & Planck constant & $6.626 \times 10^{-34}$\,Js \\
  $\hbar$         & reduced Planck constant & $h/(2\pi)$ \\
  $H$             & Hessian of force constants & a.u.\ in Gaussian \\
  $\hat{H}$       & Hamiltonian operator & J \\
  $i, j, k$       & mode indices        & $1, \ldots, N_\text{modes}$ \\
  $k_B$           & Boltzmann constant & $1.381 \times 10^{-23}$\,J/K \\
  $\mathcal{L}$   & set of ligand atoms (Fe-coordinated) & $|\mathcal{L}| \approx 4$--$8$ \\
  $m_a$           & mass of atom $a$ & amu \\
  $\mathbf{n}$    & cluster normal (best-fit plane) & dimensionless unit vector \\
  $N_\text{atoms}$ & number of atoms in the system & $\approx 50$--$3000$ \\
  $N_\text{modes}$ & number of vibrational modes in the filter range & $\approx 100$--$10\,000$ \\
  $\mathbf{r}_a$  & position of atom $a$ in Cartesian coords. & $\angstrom$ \\
  $r_\text{eq}$   & equilibrium bond distance & $\angstrom$ \\
  $S$             & Huang-Rhys factor & dimensionless \\
  $T$             & temperature & K \\
  $u_\text{rms}$  & RMS amplitude of a mode (QHO expectation) & $\angstrom$ \\
  $X, Y, Z$       & bonding channel index & $\in \{\text{FeFe, FeN, FeS, NH, HA}\}$ \\
  \bottomrule
\end{tabular}

\subsection{Greek letters}

\begin{tabular}{p{1.5cm} p{6cm} p{4cm}}
  \toprule
  \textbf{Symbol} & \textbf{Meaning} & \textbf{Unit / range} \\
  \midrule
  $\alpha$        & in SCSD: $D_\text{2h}$ irrep coefficient & $\angstrom$ \\
  $\alpha_X$      & mode contribution to bond $X$ & $[0, 1]$ \\
  $\Delta Q$      & static mode displacement (Marcus-Hush) & amu$^{1/2}\angstrom$ \\
  $\Delta r_X$    & bond modulation & $\angstrom$ \\
  $\hat{\mathbf{n}}$ & cluster normal & dimensionless unit vector \\
  $\lambda$       & Marcus-Hush reorganization energy & cm$^{-1}$ \\
  $\lambda_X(i)$  & per-mode contribution, channel $X$, mode $i$ & cm$^{-1}$ \\
  $\Lambda_X^\text{total}$ & system-total reorg, channel $X$ & cm$^{-1}$ \\
  $\Lambda_X(\omega)$ & cumulative reorg curve & cm$^{-1}$ \\
  $\mu_\text{pair}$ & reduced mass of a bond & amu \\
  $\mu_\text{mode}$ & reduced mass of a mode (from DFT) & amu \\
  $\sigma$        & Gaussian broadening of the spectra & cm$^{-1}$ \\
  $\sigma_X$      & statistical uncertainty of $X$ & like $X$ \\
  $\Phi$          & phase of a mode displacement (internal) & rad \\
  $\omega$        & mode frequency & cm$^{-1}$ or rad/s \\
  $\omega_i$      & frequency of mode $i$ & cm$^{-1}$ \\
  $\mathbf{\Omega}$ & vector of all mode frequencies & cm$^{-1}$ \\
  \bottomrule
\end{tabular}

\subsection{Special symbols and abbreviations}

\begin{tabular}{p{2cm} p{8.5cm}}
  \toprule
  \textbf{Symbol} & \textbf{Meaning} \\
  \midrule
  $\angstrom$    & \AA{}ngstrom unit \\
  $\langle\cdot\rangle$ & quantum-mechanical / thermal expectation value \\
  $|\cdot|$      & magnitude (scalar) or norm (vector) \\
  $\|\cdot\|$    & L2 norm of a vector \\
  $\hat{\cdot}$  & unit vector (e.g.\ $\hat{\mathbf{n}}$) \\
  $\cdot^*$      & complex conjugate \\
  $\sum_i$       & sum over mode indices (within the filter range) \\
  $\sum_a$       & sum over atom indices \\
  $\sum_X$       & sum over channel indices (FeFe, FeN, FeS, NH, HA) \\
  RSS            & Root Sum of Squares: $\sqrt{\sum_i x_i^2}$ \\
  RMS            & Root Mean Square: $\sqrt{\frac{1}{n}\sum_i x_i^2}$ \\
  pDOS           & projected density of states \\
  pCH            & pyrrole C-H (in NRVS literature) \\
  HP             & High-Precision (Gaussian eigenvectors) \\
  Std            & standard precision (Gaussian eigenvectors) \\
  PCET           & Proton-Coupled Electron Transfer \\
  ET             & Electron Transfer \\
  PT             & Proton Transfer \\
  NRVS           & Nuclear Resonance Vibrational Spectroscopy \\
  SCSD           & Symmetry-Coordinate Structural Decomposition \\
  OOP            & Out-of-Plane (perpendicular to the cluster plane) \\
  INP            & In-Plane (within the cluster plane) \\
  QHO            & quantum-mechanical harmonic oscillator \\
  WT             & wild type \\
  Mut            & mutant \\
  prot           & protonated (imidazole NH present) \\
  deprot         & deprotonated (imidazole NH removed) \\
  $D_{2h}$       & point group $D_{2h}$ (rhombic cluster) \\
  $A_g$, $B_{1g}$, $\ldots$ & irreps of $D_{2h}$ symmetry \\
  \bottomrule
\end{tabular}

% -------------------------------------------------------------------------
\section{Glossary}
\label{app:glossary}

Terms and concepts used in the scientific literature and this
manual, with brief definitions. Sorted alphabetically.

\paragraph{Acceptor (H-bond).}
Atom with a lone pair of electrons that receives a hydrogen bond.
In [2Fe-2S] proteins, the typical acceptors are: backbone carbonyl
oxygens (C=O), hydroxyl oxygens of Tyr/Ser/Thr, or water
molecules. The ``main acceptor'' of a His N-H is the one with the
largest Gaussian weight (distance closest to the optimum
$r_0 = 2.8\,\angstrom$).

\paragraph{$B$ factor (Debye-Waller factor).}
Mean squared displacement of an atom in thermal equilibrium:
$B_a = 8\pi^2\langle u_a^2\rangle/3$. Units of $\angstrom^2$.
Computed from the modes as
$B_a = 8\pi^2/3 \cdot \sum_i u_\text{rms}^2(i) |e_a(i)|^2$.

\paragraph{Bonding channel.}
A bonding channel is a class of all chemically equivalent bonds
whose modulation is aggregated together. Five channels are
computed in the tool: FeFe (cluster breathing), FeN
(Fe-N$_\text{His}$), FeS
(Fe-S$_\text{Cys}$ + Fe-S$_\text{bridge}$), NH
(imidazole stretching), HA (H-acceptor;
reaction-coordinate mode).

\paragraph{Cluster breathing.}
Symmetric mode in which the Fe-Fe distance changes
(stretching/compression of the cluster rhombus). $A_g$ symmetry in
$D_{2h}$. Frequency $\sim 200$\,cm$^{-1}$ in typical [2Fe-2S]
proteins.

\paragraph{Cluster normal $\hat{\mathbf{n}}$.}
Unit vector perpendicular to the best-fit plane through the four
cluster atoms (2 Fe + 2 S$_\mu$). Computed via SVD of the centered
atomic coordinates. Defines the OOP/INP separation.

\paragraph{Co-modulation.}
Simultaneous modulation of two bonding channels by the same mode.
Indicator: $C_\text{PCET}(\omega) = (M_\text{NH}(\omega) \cdot
M_\text{HA}(\omega))^{1/2}$. Large values for reaction-relevant
PCET modes.

\paragraph{Cys$_2$His$_2$, Cys$_3$His$_1$, Cys$_4$.}
Ligand schemes of the [2Fe-2S] cluster:
\begin{itemize}\setlength{\itemsep}{0pt}
  \item Cys$_4$: all four Fe-coordinating ligands are cysteinates
        (S$^-$). Standard ferredoxins.
  \item Cys$_2$His$_2$: two Cys + two His. Rieske type. PCET
        possible.
  \item Cys$_3$His$_1$: three Cys + one His. mitoNEET type.
\end{itemize}

\paragraph{$D_{2h}$ irreps.}
Irreducible representations of the point group $D_{2h}$ (rhombic
cluster with three orthogonal $C_2$ axes plus an inversion center).
Eight irreps: $A_g, B_{1g}, B_{2g}, B_{3g}, A_u, B_{1u}, B_{2u},
B_{3u}$. SCSD decomposes every mode displacement into exactly one
linear combination of these eight irreps.

\paragraph{DFT (Density Functional Theory).}
Quantum-chemical method for computing electronic structures and
vibrational frequencies. The tool expects DFT frequencies from
Gaussian-16 or ORCA-6 (typically B3LYP/6-31G* converged to
$10^{-8}$\,Hartree).

\paragraph{Eigenvector.}
Mass-weighted displacement components of a vibrational mode,
relative to the equilibrium geometry. For atom $a$ in mode $i$:
$\mathbf{e}_a^{(i)}$. Normalization in Gaussian: $\sum_a m_a
|\mathbf{e}_a^{(i)}|^2 = 1$ (mass-weighted).

\paragraph{Embedding.}
Low-dimensional representation of high-dimensional mode feature
vectors. In the tool: a 31-dimensional feature matrix is mapped to
2D via UMAP.

\paragraph{Feature matrix.}
$N_\text{modes} \times 31$ matrix with 31 scalars per mode
(frequency, OOP\%, cluster contribution, four ligand OOP, SCSD
irrep contributions, etc.). Input to the UMAP embedding.

\paragraph{Frequency window.}
Frequency range \texttt{(lo, hi)} in cm$^{-1}$ to be analyzed
separately. Multi-window mode: several windows in parallel, each
with their own reorganization aggregates and embeddings.

\paragraph{Gaussian eigenvectors (HP / Std).}
Gaussian writes vibrational-mode eigenvectors in two blocks:
\begin{itemize}\setlength{\itemsep}{0pt}
  \item \textbf{Standard (Std)}: 4 decimal digits, compact, always
        present.
  \item \textbf{High Precision (HP)}: 5 decimal digits, activated
        via \texttt{freq=hpmodes}; requires more storage in the
        log file. Recommended for this tool.
\end{itemize}

\paragraph{HDBSCAN.}
Hierarchical Density-Based Spatial Clustering of Applications
with Noise. A clustering algorithm that groups UMAP points into
clusters and marks noise separately.

\paragraph{HPmodes.}
Gaussian keyword (\texttt{freq=hpmodes}) that writes the HP
eigenvectors additionally to the log file.

\paragraph{Huang-Rhys factor.}
$S = \frac{\mu\omega(\Delta q)^2}{2\hbar}$. The dimensionless
displacement of a harmonic mode between electronic states, in
units of the zero-point amplitude. Related to our $\lambda_X$ but
in a different convention (displacement-squared vs.\ amplitude
squared).

\paragraph{Imidazole.}
Aromatic five-ring with two nitrogens (N$_\delta$ and N$_\epsilon$)
in histidine. One of the nitrogens (typically N$_\epsilon$) binds
the Fe; the other can be protonated or not, depending on the
tautomer.

\paragraph{INP (in-plane).}
Mode displacement predominantly in the cluster plane (perpendicular
to $\hat{\mathbf{n}}$). Classified with OOP\,\% < 40 (binary
threshold at 60).

\paragraph{Kabsch alignment.}
Algorithm for the optimal superposition of two atom-coordinate
sets (PDB $\leftrightarrow$ Gaussian). Yields a rotation matrix
and an RMSD.

\paragraph{Kernel classification.}
Heuristic mode-type assignment: cluster breathing, cluster hinge,
cluster twist, ligand stretch, etc. Implemented in
\texttt{core.py:classify\_kern}. Not orthogonal to the SCSD
irreps; multiple kernel types can correspond to the same irrep.

\paragraph{Lamb-M\"o{\ss}bauer factor.}
$f_\text{LM} = \exp(-2W)$, with $W$ the Debye-Waller factor of the
M\"o{\ss}bauer-active nucleus ($^{57}$Fe). Measures the fraction of
the resonance absorption that occurs without phonon excitation
(elastic component in the NRVS spectrum).

\paragraph{Marcus-Hush theory.}
Classical theory of electron transfer (Marcus 1956, Hush 1958)
that captures the reaction rate via a harmonic approximation of
the reactant and product energy surfaces with a single
reorganization coordinate.

\paragraph{mitoNEET.}
A Cys$_3$His$_1$-coordinated [2Fe-2S] protein. Known for its
particularly well-studied NRVS band structure
(Gee 2021, Bergner 2017).

\paragraph{Modulation spectrum $M_X(\omega)$.}
Gaussian-broadened spectrum of bond modulations per channel:
$M_X(\omega) = \sum_i |\Delta r_X(i)| \cdot G(\omega - \omega_i,
\sigma)$. Shows in which frequency range bond $X$ is strongly
modulated.

\paragraph{NRVS (Nuclear Resonance Vibrational Spectroscopy).}
Synchrotron spectroscopy on M\"o{\ss}bauer-active nuclei
(typically $^{57}$Fe). Measures the $^{57}$Fe-projected
vibrational density of states. Complementary to Raman/IR
(element-selective, not symmetry-selective). The tool does not
itself compute the NRVS spectrum; it provides reorganization
contributions that can be used for band identification.

\paragraph{OOP (out-of-plane).}
Mode displacement predominantly perpendicular to the cluster plane
(parallel to $\hat{\mathbf{n}}$). Classified with OOP\,\% > 60.

\paragraph{OOP\,\%.}
Squared fraction of the mode displacement along
$\hat{\mathbf{n}}$, relative to the total displacement.
$\text{OOP\,\%} = 100 \cdot \sum_a (\mathbf{e}_a \cdot
\hat{\mathbf{n}})^2 / \sum_a |\mathbf{e}_a|^2$.

\paragraph{PCET (Proton-Coupled Electron Transfer).}
The concerted transfer of an electron and a proton. In [2Fe-2S]
proteins, typically: cluster reduction
($\text{Fe}^{3+} + e^- \to \text{Fe}^{2+}$) coupled with
imidazole NH deprotonation (Rieske: the $pK_a$ shifts by 4--5
units between oxidized and reduced state).

\paragraph{Pseudocode.}
Algorithm description in an informal but precise notation, without
a specific programming language. Used in
appendix~\ref{app:pseudocodes} for the central algorithms.

\paragraph{Reaction-coordinate mode (HA).}
Definition of the HA modulation: instead of modulating the
H-acceptor distance, the motion of the H atom is computed relative
to the donor N along the N$\to$acceptor axis. Classical notation:
$\Delta r_\text{HA} = (\mathbf{e}_H - \mathbf{e}_N)
\cdot \hat{\mathbf{n}}_{N\to\text{Acc}}$. Measures genuine proton
motion rather than nonspecific skeletal modes.

\paragraph{Reorganization energy $\lambda_X$.}
In the tool: thermal bond-length-fluctuation energy per mode along
bond $X$, computed in Marcus-Hush form with thermally scaled mode
displacements. Related to but not identical with the classical
Marcus-Hush reorganization energy (see \S\ref{sec:reorg}).

\paragraph{Rieske protein.}
A Cys$_2$His$_2$-coordinated [2Fe-2S] protein. A component of
cytochrome $bc_1$ complexes, involved in the Q-cycle. A classical
PCET system (imidazole NH deprotonation coupled to cluster
reduction).

\paragraph{RMS (Root Mean Square).}
$\sqrt{\frac{1}{n}\sum_i x_i^2}$. Mean of the squares, then square
root.

\paragraph{RSS (Root Sum of Squares).}
$\sqrt{\sum_i x_i^2}$. Sum of the squares, then square root. Not
to be confused with RMS (which is divided by $n$).

\paragraph{SCSD (Symmetry-Coordinate Structural Decomposition).}
Method by Kingsbury \& Senge (2024), which represents structural
distortions of an arbitrary geometry as a linear combination of
symmetry coordinates of a canonical reference geometry. In the
tool: $D_{2h}$ decomposition with a canonical reference [2Fe-2S]
cluster (Fe-Fe = 2.73\,$\angstrom$, Fe-S = 2.20\,$\angstrom$).

\paragraph{Secondary structure (SSE).}
Local 3D structure of a protein (helix, beta sheet, loop). In the
tool, derived from the PDB file (HELIX/SHEET records, DSSP
fallback, or $\varphi/\psi$ heuristic).

\paragraph{TOML.}
Configuration format (\emph{Tom's Obvious, Minimal Language}).
In the tool: \texttt{config.toml} or \texttt{run.toml} for
reproducible runs.

\paragraph{Thermally scaled eigenvectors.}
$\mathbf{e}_a^\text{thermal}(i) = \mathbf{e}_a^\text{normalized}(i)
\cdot u_\text{rms}(i)$. Eigenvectors multiplied by the
mode-specific RMS amplitude (the QHO expectation value at the
given temperature). Directly usable for computing thermal bond
modulations.

\paragraph{UMAP (Uniform Manifold Approximation and Projection).}
Topology-preserving dimensionality-reduction method (McInnes 2018).
Local and global structures are preserved simultaneously;
deterministically reproducible with a seed.

\paragraph{Hydrogen bond.}
Weak but directional interaction between a proton donor (X-H with
$X = \text{N, O}$) and an acceptor ($Y$ with a lone pair, $Y = \text{O,
N}$). Energies typically 5--30\,kJ/mol. Responsible in proteins for
secondary structure and PCET reaction channels.

% -------------------------------------------------------------------------
\section{Algorithm pseudocodes}
\label{app:pseudocodes}

The central algorithms of the tool are documented here as
pseudocode. Language level: between formal mathematics and Python
-- no type details, but all non-trivial control flow explicit.
Variable names follow the module names in the code.

The pseudocodes are \emph{not} 1:1 implementations, but sketches of
the logic. The actual code in \texttt{src/modenanalyse\_2fe2s/} has
additional performance optimizations, caching, logging, and error
handling that are omitted here for readability.

\subsection{Cluster detection}
\label{app:pseudo-cluster}

Central task: find all [2Fe-2S] clusters from a list of atoms. In
multi-cluster systems (glutaredoxin dimers etc.), several clusters
must be unambiguously grouped and sorted.

\paragraph{Algorithm.}

\begin{algorithm}[H]
\caption{\textsc{FindAllClusters}}
\begin{algorithmic}[1]
\Require Atom list $A = \{a_1, \ldots, a_N\}$ with element +
         Cartesian position; cutoffs $c_\text{FeFe}, c_\text{FeS}$
         in $\angstrom$
\Ensure List of clusters $\mathcal{C}_1, \mathcal{C}_2, \ldots$,
        sorted by Fe-Fe distance ascending
\State $\text{Fe} \gets \{a \in A : \text{element}(a) = \text{Fe}\}$
\State $\text{S}  \gets \{a \in A : \text{element}(a) = \text{S}\}$
\If{$|\text{Fe}| < 2$ \textbf{or} $|\text{S}| < 2$}
  \State \Return $\emptyset$  \Comment{no cluster possible}
\EndIf
\State Form all Fe pairs $(f_i, f_j)$ with
       $\|\mathbf{r}(f_i) - \mathbf{r}(f_j)\| \leq c_\text{FeFe}$
\State Sort the Fe pairs by distance ascending
\State $\text{used} \gets \emptyset$  \Comment{set of atoms already assigned}
\State $\text{cluster\_list} \gets \emptyset$
\ForAll{Fe pairs $(f_i, f_j)$ in sorted order}
  \If{$f_i \in \text{used}$ \textbf{or} $f_j \in \text{used}$}
    \State \textbf{continue}  \Comment{atom already assigned}
  \EndIf
  \State $\mathcal{S}_\mu \gets \{s \in \text{S} :$ S bridges $f_i, f_j$
         with $\|\mathbf{r}(s) - \mathbf{r}(f_i)\| \leq c_\text{FeS}$
         \textbf{and} $\|\mathbf{r}(s) - \mathbf{r}(f_j)\| \leq c_\text{FeS}\}$
  \If{$|\mathcal{S}_\mu| \geq 2$}  \Comment{enough bridging sulfurs?}
    \State $(s_1, s_2) \gets$ the two closest S in $\mathcal{S}_\mu$
    \State $\text{cluster\_list} \gets \text{cluster\_list} \cup
           \{\{f_i, f_j, s_1, s_2\}\}$
    \State $\text{used} \gets \text{used} \cup \{f_i, f_j, s_1, s_2\}$
  \EndIf
\EndFor
\State \Return $\text{cluster\_list}$
\end{algorithmic}
\end{algorithm}

\paragraph{Implementation subtleties.}
\begin{itemize}\setlength{\itemsep}{0pt}
  \item Sorting Fe pairs by distance is \emph{deterministic} and
        ensures the tightest cluster is assigned first. Hence
        \texttt{cluster\_index = 0} is always the tightest one,
        regardless of atom order in the log file.
  \item The cutoffs $c_\text{FeFe}$ and $c_\text{FeS}$ are config
        parameters (defaults $3.5$ and $3.0\,\angstrom$). For very
        extended clusters (e.g.\ [4Fe-4S] or Mo-Fe cofactor),
        these may need to be raised -- but then the tool only
        reads the [2Fe-2S] sub-fragment.
  \item Once an atom has been ``consumed'' by a tighter cluster,
        it is not reused for another cluster. Hence the algorithm
        finds the true number of clusters even in densely packed
        multi-cluster systems.
\end{itemize}

\subsection{Cluster normal via SVD}
\label{app:pseudo-svd}

From the four cluster atoms, the best-fit plane and hence the
cluster normal $\hat{\mathbf{n}}$ is determined by singular-value
decomposition.

\begin{algorithm}[H]
\caption{\textsc{ClusterNormal}}
\begin{algorithmic}[1]
\Require Cluster atoms $\mathcal{C} = \{a_1, a_2, a_3, a_4\}$
        with positions
        $\mathbf{r}_1, \ldots, \mathbf{r}_4 \in \mathbb{R}^3$
\Ensure Unit vector $\hat{\mathbf{n}}$ perpendicular to the
        best-fit plane
\State $\bar{\mathbf{r}} \gets \frac{1}{4} \sum_{i=1}^{4} \mathbf{r}_i$
       \Comment{centroid}
\State Construct the centered matrix
       $\tilde{\mathbf{R}} \in \mathbb{R}^{4\times 3}$:
       $\tilde{\mathbf{R}}_{i,:} = \mathbf{r}_i - \bar{\mathbf{r}}$
\State $\mathbf{U}, \boldsymbol{\Sigma}, \mathbf{V}^\top \gets
       \textsc{SVD}(\tilde{\mathbf{R}})$
       \Comment{$\mathbf{V}$ is an orthogonal $3\times 3$ matrix}
\State $\hat{\mathbf{n}} \gets \mathbf{V}_{:,3}$
       \Comment{third column = direction of smallest variance}
\If{$\hat{\mathbf{n}}_z < 0$}  \Comment{sign convention}
  \State $\hat{\mathbf{n}} \gets -\hat{\mathbf{n}}$
\EndIf
\State \textbf{return} $\hat{\mathbf{n}}$
\end{algorithmic}
\end{algorithm}

Remark: the folding residual is $\Sigma_{33}$ (the smallest
singular value). Values $> 0.1\,\angstrom$ indicate a strongly
non-planar cluster -- unusual for native [2Fe-2S] clusters but
possible in strongly distorted mutants.

\subsection{Kabsch alignment (PDB $\leftrightarrow$ Gaussian)}
\label{app:pseudo-kabsch}

When a PDB file is provided, the PDB coordinates must be mapped to
the Gaussian coordinates -- the geometries are usually
rotation-shifted relative to each other. The Kabsch algorithm
\cite{Kabsch1976} provides the optimal rotation.

\begin{algorithm}[H]
\caption{\textsc{KabschAlign}}
\begin{algorithmic}[1]
\Require Point sets $\mathbf{P}, \mathbf{Q} \in \mathbb{R}^{n\times 3}$,
         in correspondence (same atom order)
\Ensure Rotation $\mathbf{R}$, translation $\mathbf{t}$, RMSD
\State $\bar{\mathbf{p}} \gets \frac{1}{n}\sum_i \mathbf{P}_{i,:}$,
       $\bar{\mathbf{q}} \gets \frac{1}{n}\sum_i \mathbf{Q}_{i,:}$
\State $\mathbf{P}' \gets \mathbf{P} - \bar{\mathbf{p}}$,
       $\mathbf{Q}' \gets \mathbf{Q} - \bar{\mathbf{q}}$
       \Comment{shift each centroid to the origin}
\State $\mathbf{H} \gets \mathbf{P}'^\top \mathbf{Q}'$
       \Comment{$3\times 3$ cross-correlation}
\State $\mathbf{U}, \boldsymbol{\Sigma}, \mathbf{V}^\top \gets
       \textsc{SVD}(\mathbf{H})$
\State $d \gets \text{sign}(\det(\mathbf{V} \mathbf{U}^\top))$
       \Comment{reflection sign}
\State $\mathbf{D} \gets \text{diag}(1, 1, d)$
\State $\mathbf{R} \gets \mathbf{V} \mathbf{D} \mathbf{U}^\top$
       \Comment{optimal rotation}
\State $\mathbf{t} \gets \bar{\mathbf{q}} - \mathbf{R} \bar{\mathbf{p}}$
\State $\text{RMSD} \gets \sqrt{\frac{1}{n} \|\mathbf{R}\mathbf{P}' -
       \mathbf{Q}'\|_F^2}$
\State \textbf{return} $(\mathbf{R}, \mathbf{t}, \text{RMSD})$
\end{algorithmic}
\end{algorithm}

\paragraph{What the $\mathbf{D}$ matrix does.}
The ``naive'' solution $\mathbf{R} = \mathbf{V}\mathbf{U}^\top$ may
be a reflection (determinant $-1$) if the two point sets are
chirally different (e.g.\ accidentally mirrored). The $\mathbf{D}$
matrix corrects the sign of the third singular-vector axis, so
that an actual rotation (determinant $+1$) is returned. For
correctly corresponding point sets, $d = +1$, hence
$\mathbf{D} = \mathbf{I}$.

\subsection{Mode loop with thermal scaling}
\label{app:pseudo-mode-loop}

The central loop of the tool: for each vibrational mode the
eigenvector is read, thermally scaled, and fed into the OOP/INP,
reorganization, and SCSD pipelines.

\begin{algorithm}[H]
\caption{\textsc{AnalyzeMode}}
\begin{algorithmic}[1]
\Require Mode index $i$, frequency $\omega_i$, reduced mass $\mu_i$,
         temperature $T$, eigenvector $\mathbf{e}^{(i)}$
         (mass-weighted, normalized), cluster atoms $\mathcal{C}$,
         ligands $\mathcal{L}$
\Ensure Result dict $r$ with OOP\%, $\Delta r_X$, $\lambda_X$,
        SCSD contributions, kernel classification
\State \Comment{1. Thermal RMS amplitude (quantum expectation value)}
\State $u_\text{rms}(i) \gets \sqrt{\frac{\hbar}{2\mu_i\omega_i}
       \coth\!\left(\frac{\hbar\omega_i}{2 k_B T}\right)}$
\State $\mathbf{e}^\text{therm}_a \gets \mathbf{e}_a^{(i)} \cdot u_\text{rms}(i)
       \quad \forall a \in \mathcal{C} \cup \mathcal{L}$
\State \Comment{2. OOP/INP classification}
\State $\mathbf{e}_a^\text{OOP} \gets (\mathbf{e}^\text{therm}_a \cdot
       \hat{\mathbf{n}})\,\hat{\mathbf{n}}$
\State $\alpha_\text{OOP} \gets 100 \cdot
       \frac{\sum_{a\in\mathcal{C}} \|\mathbf{e}_a^\text{OOP}\|^2}
            {\sum_{a\in\mathcal{C}} \|\mathbf{e}^\text{therm}_a\|^2}$
\State $r.\text{oop\_pct} \gets \alpha_\text{OOP}$
\State $r.\text{type} \gets \textsc{ClassifyOOP}(\alpha_\text{OOP},
       \theta_\text{OOP}=60)$
\State \Comment{3. Bond modulations}
\ForAll{Bonding channels $X \in \{\text{FeFe, FeN, FeS, NH, HA}\}$}
  \State $\Delta r_X(i) \gets \textsc{ComputeBondModulation}(X,
         \mathbf{e}^\text{therm}, \text{geometry})$
       \Comment{see \S\ref{sec:reorg}}
  \State $\lambda_X(i) \gets \tfrac{1}{2}\mu_i \omega_i^2 \,(\Delta r_X(i))^2$
\EndFor
\State \Comment{4. SCSD symmetry decomposition (optional)}
\If{\texttt{cfg.analyze\_scsd}}
  \State $\textsc{Displacement} \gets \mathbf{e}^\text{therm}|_\mathcal{C}$
       \Comment{cluster atoms only}
  \State $\textsc{Irrep contributions} \gets \textsc{ProjectOnD2hIrreps}(
         \textsc{Displacement}, \text{reference geometry})$
  \State $r.\text{scsd} \gets \textsc{Irrep contributions}$
\EndIf
\State \Comment{5. Heuristic kernel classification}
\State $r.\text{kern\_type} \gets \textsc{ClassifyKern}(\mathbf{e}^\text{therm},
       \mathcal{C}, \alpha_\text{OOP})$
\State \textbf{return} $r$
\end{algorithmic}
\end{algorithm}

\paragraph{Subtleties:}
\begin{itemize}\setlength{\itemsep}{0pt}
  \item For \texttt{temp\_k = None}, the thermal scaling is
        skipped; $u_\text{rms}$ is replaced by
        \texttt{cfg.amplitude} (default 1.0). This is the
        ``classical'' mode in which all modes are weighted equally.
  \item At low frequency and high temperature the coth argument
        is small and $u_\text{rms}$ becomes large -- for modes
        below $\sim 30$\,cm$^{-1}$ this can yield contributions
        $> 50$\,\% of the total reorganization in a typical filter
        range. Hence the recommendation: always set
        \texttt{freq\_min} to 50 or 100 if low-frequency skeletal
        motions should be suppressed.
  \item In the ``no-His'' path (no histidine in the cluster), the
        NH and HA channels are zeroed \emph{before} the
        reorganization loop -- not via zero $\Delta r$, but
        directly by the upstream call in \texttt{runner.py}. This
        saves geometric computations and makes the path
        deterministically clean.
\end{itemize}

\subsection{Reorganization aggregation per mode}
\label{app:pseudo-reorg}

Per mode, the bond modulations $\Delta r_X$ are computed from the
thermally scaled eigenvectors.

\begin{algorithm}[H]
\caption{\textsc{ComputeBondModulation} (FeFe example)}
\begin{algorithmic}[1]
\Require Bond atoms $a, b$ (with positions + thermally scaled
         displacements $\mathbf{e}_a, \mathbf{e}_b$)
\Ensure Bond modulation $\Delta r$ (signed)
\State $\mathbf{r}_a^\text{new} \gets \mathbf{r}_a + \mathbf{e}_a$
\State $\mathbf{r}_b^\text{new} \gets \mathbf{r}_b + \mathbf{e}_b$
\State $r_\text{new} \gets \|\mathbf{r}_a^\text{new} - \mathbf{r}_b^\text{new}\|$
\State $r_\text{eq} \gets \|\mathbf{r}_a - \mathbf{r}_b\|$
\State $\Delta r \gets r_\text{new} - r_\text{eq}$
\State \textbf{return} $\Delta r$
\end{algorithmic}
\end{algorithm}

For multi-bond channels (e.g.\ FeS with four Fe-S bonds),
\textsc{ComputeBondModulation} is called per bond, and the
contributions are aggregated \emph{sign-free} via RSS
(\S\ref{app:notation}):
\begin{equation}
  \Delta r_X^\text{rss} = \sqrt{\sum_{(a,b) \in X} (\Delta r_{(a,b)})^2}
\end{equation}
The channel contribution $\lambda_X = \tfrac{1}{2}\mu\omega^2
(\Delta r_X^\text{rss})^2$ is then identical with the sum of the
pair contributions:
\begin{equation}
  \lambda_X = \sum_{(a,b) \in X} \tfrac{1}{2}\mu\omega^2 (\Delta r_{(a,b)})^2.
  \label{eq:reorg-additivity}
\end{equation}
This Marcus-Hush additivity is the central consistency condition
and smoke test (\S\ref{sec:workflow-worked-example}).

\subsection{HA reaction-coordinate computation}
\label{app:pseudo-ha}

The HA channel measures the H-acceptor distance modulation along
the reaction coordinate (rather than the naive H-Acc distance).

\begin{algorithm}[H]
\caption{\textsc{ComputeHAReactionCoord}}
\begin{algorithmic}[1]
\Require N$_\text{donor}$, H, acceptor $A$ (positions + displacements)
\Ensure Reaction-coordinate modulation $\Delta r_\text{HA}$
\State $\hat{\mathbf{n}}_{N\to A} \gets
       \frac{\mathbf{r}_A - \mathbf{r}_N}
            {\|\mathbf{r}_A - \mathbf{r}_N\|}$
       \Comment{donor-acceptor axis, unit vector}
\State $\Delta\mathbf{e} \gets \mathbf{e}_H - \mathbf{e}_N$
       \Comment{H motion relative to the donor}
\State $\Delta r_\text{HA} \gets \Delta\mathbf{e} \cdot
       \hat{\mathbf{n}}_{N\to A}$
       \Comment{projection onto the reaction coordinate}
\State \textbf{return} $\Delta r_\text{HA}$
\end{algorithmic}
\end{algorithm}

\paragraph{Why this approach?}
The naive H-Acc distance measures the \emph{total motion} of the
H atom relative to the acceptor atom. Pure skeletal modes (in
which the entire complex moves collectively) would then show a
large $\Delta r$, even though the proton is not moving
\emph{relative} to the donor at all -- such modes are not real
PCET-relevant modes. The reaction-coordinate convention extracts
only the genuine proton motion.

\subsection{Modulation spectrum and cumulative reorganization curve}
\label{app:pseudo-spektren}

From the per-mode values, frequency-resolved spectra are built.

\begin{algorithm}[H]
\caption{\textsc{ModulationSpectrum}}
\begin{algorithmic}[1]
\Require Mode list with $(\omega_i, \Delta r_X(i))$, broadening $\sigma$,
         frequency grid $\boldsymbol{\omega}$
\Ensure Modulation spectrum $M_X(\omega)$
\State $M_X(\omega) \gets 0 \quad \forall \omega \in \boldsymbol{\omega}$
\ForAll{Modes $i$}
  \State $g_i(\omega) \gets \frac{1}{\sigma\sqrt{2\pi}}
         \exp\!\left(-\frac{(\omega - \omega_i)^2}{2\sigma^2}\right)$
  \State $M_X(\omega) \gets M_X(\omega) + |\Delta r_X(i)| \cdot g_i(\omega)$
\EndFor
\State \textbf{return} $M_X(\boldsymbol{\omega})$
\end{algorithmic}
\end{algorithm}

\begin{algorithm}[H]
\caption{\textsc{CumulativeReorg}}
\begin{algorithmic}[1]
\Require Mode list with $(\omega_i, \lambda_X(i))$, sorted ascending in $\omega$
\Ensure Cumulative curve $\Lambda_X(\omega)$
\State $\Lambda \gets 0$
\ForAll{Modes $i$ (sorted ascending)}
  \State Write $(\omega_i^-, \Lambda)$ and $(\omega_i^+, \Lambda + \lambda_X(i))$
         \Comment{step function: jump at $\omega_i$}
  \State $\Lambda \gets \Lambda + \lambda_X(i)$
\EndFor
\State Write the end value $(\omega_\text{max}, \Lambda)$
\State \textbf{return} $\Lambda_X(\omega)$
\end{algorithmic}
\end{algorithm}

Both spectra are written to the sheets
\texttt{Modulation\_spectra} and \texttt{Lambda\_cumulative}
respectively (\S\ref{sec:ausgabe}).

\subsection{HDBSCAN clustering of the embedding}
\label{app:pseudo-hdbscan}

After the UMAP projection to 2D, the points are clustered
density-based with HDBSCAN \cite{Campello2013}.

\begin{algorithm}[H]
\caption{\textsc{HDBSCANCluster}}
\begin{algorithmic}[1]
\Require 2D embedding points $\{x_1, \ldots, x_N\}$,
         min\_cluster\_size, min\_samples
\Ensure Cluster labels $\ell_1, \ldots, \ell_N \in \{-1, 0, 1, \ldots\}$
       ($-1$ = noise)
\State \Comment{1. Weighted graph: each point has distances to its K nearest neighbors}
\State $\text{core\_dist}(x_i) \gets$ distance to the
       \texttt{min\_samples}-th neighbor
\State $\text{mutual\_reachability}(x_i, x_j) \gets$
       $\max(\text{core\_dist}(x_i), \text{core\_dist}(x_j),
       \|x_i - x_j\|)$
\State \Comment{2. Minimum spanning tree of the mutual\_reachability matrix}
\State $\text{MST} \gets \textsc{KruskalMST}(\text{mutual\_reachability})$
\State \Comment{3. Hierarchical cluster tree}
\State $\text{Tree} \gets \textsc{ConvertMSTtoHierarchy}(\text{MST})$
\State \Comment{4. Select stable clusters (excess of mass)}
\State $\text{cluster\_list} \gets \textsc{ExcessOfMassSelection}(\text{Tree},
       \text{min\_cluster\_size})$
\State \Comment{5. Assign points to clusters or label them as noise}
\ForAll{Points $x_i$}
  \If{$x_i$ is a member of a stable cluster $C_k$}
    \State $\ell_i \gets k$
  \Else
    \State $\ell_i \gets -1$  \Comment{noise}
  \EndIf
\EndFor
\State \textbf{return} $\boldsymbol{\ell}$
\end{algorithmic}
\end{algorithm}

\paragraph{Why HDBSCAN instead of DBSCAN?}
\begin{itemize}\setlength{\itemsep}{0pt}
  \item HDBSCAN does not need a global $\epsilon$ threshold (DBSCAN
        does -- which is problematic for varying cluster
        densities).
  \item HDBSCAN finds clusters of varying density automatically
        (via its hierarchical construction).
  \item HDBSCAN is robust against outliers (they automatically end
        up as noise, without a threshold needing to be tuned).
\end{itemize}

In the default configuration, the tool uses
\texttt{min\_cluster\_size = 8} and \texttt{min\_samples = 5}; see
\S\ref{sec:embeddings} for the rationale.


% -------------------------------------------------------------------------
\section{Excel sheet full reference}
\label{app:sheet-referenz}

Per run, up to \emph{four} Excel workbooks are produced. This
appendix documents every column of every sheet, so that table
values are interpretable without recourse to the code. Sheet names
are listed in the order they are created in \texttt{export.py} via
\texttt{wb.create\_sheet(...)}.

\subsection{Main analysis workbook (\texttt{*\_analysis.xlsx})}
\label{app:wb-hauptauswertung}

Contains 16 sheets with all mode-related analyses, cluster
geometry, and SCSD symmetry decomposition. The required workbook
for any publication.

\subsubsection*{\texttt{Mode\_analysis} -- mode-by-mode detail table}

One row per vibrational mode (typically 60--800 modes, depending
on the \texttt{freq\_max} filter). 32 columns:

\begin{center}
\small
\begin{tabular}{p{3.6cm}p{1.4cm}p{8.4cm}}
\toprule
\textbf{Column} & \textbf{Unit} & \textbf{Meaning} \\
\midrule
\texttt{Mode \#}         & --       & mode index from the log file (1-based) \\
\texttt{Frequency}       & cm$^{-1}$ & frequency $\omega_i$ \\
\texttt{u\_rms}          & $\angstrom$ & thermal RMS amplitude
                                       $\sqrt{\hbar/(2\mu\omega)\coth(\hbar\omega/2k_BT)}$ \\
\texttt{Red.Mass}        & AMU      & reduced mass $\mu_i$ \\
\texttt{Frc.Const.}      & mdyn/$\angstrom$ & force constant \\
\texttt{Sym.}            & --       & Mulliken symbol from Gaussian (e.g., A, $B_g$) \\
\texttt{Type}            & --       & OOP/INP classification (binary, threshold 60\,\%) \\
\texttt{Type (fine)}     & --       & 7-level OOP/INP classification \\
\texttt{Prec.}           & --       & HP/Std eigenvector source \\
\texttt{Lig OOP\%}       & \%       & OOP fraction of the mode on the ligand sphere \\
\texttt{$\pm$sigma}      & \%       & bootstrap scatter of the estimate \\
\texttt{Lig INP\%}       & \%       & INP fraction on the ligands \\
\texttt{$\pm$sigma}      & \%       & scatter \\
\texttt{Lig |d|}         & $\angstrom$ & ligand displacement (norm of the displacement) \\
\texttt{$\pm$sigma}      & $\angstrom$ & scatter \\
\texttt{2nd OOP\%}       & \%       & like \texttt{Lig OOP\%} for the 2nd ligand sphere
                                       (e.g., backbone N) \\
\texttt{$\pm$sigma}      & \%       & scatter \\
\texttt{2nd INP\%}       & \%       & 2nd ligand sphere INP \\
\texttt{$\pm$sigma}      & \%       & scatter \\
\texttt{2nd |d|}         & $\angstrom$ & 2nd ligand sphere displacement \\
\bottomrule
\end{tabular}
\end{center}

Columns 21--32 (not listed above) contain kernel contributions
(\texttt{Kern OOP\%}, \texttt{Kern INP\%}, etc.) and OOP/INP
contributions of the individual cluster atoms (Fe$_1$, Fe$_2$,
S$_1$, S$_2$).

\subsubsection*{\texttt{Groups\_OOP, \_INP, \_Angle, \_Tors}}

Four separate sheets, each with one row per mode and one column
per ligand group (Cys$_n$, His$_n$, etc.). Values are the OOP/INP
contributions or bond angles per group. Useful for the
ligand-specific analysis of individual modes.

\subsubsection*{\texttt{Fe\_S\_Cys} and \texttt{Fe\_N\_His}}

One column per Fe-ligand bond (e.g., Fe1-Cys72-S, Fe2-Cys83-S,
etc.). One row per mode, with content:
\begin{itemize}\setlength{\itemsep}{0pt}
  \item \texttt{stretch}: fraction of the mode motion along the
        bond axis.
  \item \texttt{bend\_inp}: in-plane bend contribution
        (within the cluster plane).
  \item \texttt{bend\_oop}: OOP bend contribution
        (perpendicular to the cluster plane).
\end{itemize}

These sheets are the data source for the mode-loop aggregation
(\S\ref{sec:reorg}); they show how a single mode acts on each
Fe-ligand bond.

\subsubsection*{\texttt{Kernel\_Scores} -- heuristic cluster decomposition}

One row per mode, 10 columns with the heuristic cluster score
functions from \S\ref{sec:scsd}:

\begin{itemize}\setlength{\itemsep}{0pt}
  \item \texttt{Translation}, \texttt{Umbrella}, \texttt{Hinge-Fe},
        \texttt{Hinge-S}: cluster motion classes.
  \item \texttt{OOP-Twist}, \texttt{Fe-stretch},
        \texttt{S-stretch}, \texttt{Breathing},
        \texttt{Rhombus-shear}, \texttt{Rotation-ip}:
        six in-plane / OOP motion modes.
\end{itemize}

\textbf{Important note in the sheet header}: the kernel scores
are \emph{heuristic geometric projections}; they do not generally
sum to 1 and are not orthogonal. For a \textbf{rigorous symmetry
decomposition}, the sheet refers to \texttt{SCSD} (see below).

\subsubsection*{\texttt{Equilibrium\_distances}}

Table of the equilibrium bond lengths of all relevant cluster and
ligand bonds. One row per bond. Columns: bond label, distance in
$\angstrom$, atom indices.

\subsubsection*{\texttt{Coordination} -- ligand and group atom mapping
                  \emph{(since v1.0.3)}}
\label{app:wb-main-coordination}

Diagnostic sheet that exposes the PDB-to-Gaussian atom assignments
used by all subsequent analyses. Three blocks, separated by blank
rows:
\begin{itemize}\setlength{\itemsep}{0pt}
  \item \textbf{Ligands}: one row per Fe-coordinating ligand
        (\texttt{LigandInfo} in \texttt{geometry.py}). Columns:
        \texttt{label} (e.g.\ ``Cys 207''), \texttt{element}
        (S/N/O), \texttt{fe\_center} (1-based Gaussian center
        of the bonded Fe), \texttt{lig\_center} (Gaussian
        center of the donor atom), \texttt{bond\_len\_A},
        \texttt{his\_protonated} (boolean), \texttt{h\_center}
        (Gaussian center of the imidazole NH-H if any),
        \texttt{his\_hn\_center} (the donor N that carries
        that H, NE2 or ND1).
  \item \textbf{Groups (\texttt{group\_map})}: one row per atom
        group used in the \texttt{Groups\_OOP/INP/Angle/Tors}
        sheets. Columns: \texttt{group} (label), \texttt{centers}
        (comma-separated list of 1-based Gaussian centers).
  \item \textbf{Cluster atoms}: a 4-row block with Fe$_1$,
        Fe$_2$, S$_1$, S$_2$, their Gaussian centers, and the
        equilibrium $(x, y, z)$ in $\angstrom$.
\end{itemize}

This sheet is the first place to check when an Excel row looks
unexpectedly empty: an empty \texttt{centers} list or a missing
\texttt{h\_center} explains a downstream all-zero row in the
\texttt{Groups\_*}, \texttt{Fe\_*\_*} or \texttt{His\_HN} sheets
(\S\ref{sec:migration-v104}).

\subsubsection*{\texttt{SCSD} -- rigorous $D_{2h}$ symmetry decomposition}

Per mode, 8 columns for the contributions to the $D_{2h}$ irreps:
$A_g, B_{1g}, B_{2g}, B_{3g}, A_u, B_{1u}, B_{2u}, B_{3u}$. Values
in \% of the cluster-mode contribution. In contrast to
\texttt{Kernel\_Scores}, these values are \emph{orthogonal} and
sum exactly to 100\,\%.

Method: Kingsbury \& Senge (Chem.\ Sci.\ 15, 13638, 2024) with the
canonical $D_{2h}$ reference (Fe-Fe = 2.73\,$\angstrom$,
Fe-S = 2.20\,$\angstrom$, axes $x$ = Fe-Fe, $y$ = S-S, $z$ = normal).

\subsubsection*{\texttt{Reorganization\_energy}}
\label{app:sheet-reorg}

The main sheet for the reorganization pipeline. One row per mode,
21 columns:

\begin{center}
\small
\begin{tabular}{p{4.0cm}p{1.4cm}p{7.6cm}}
\toprule
\textbf{Column} & \textbf{Unit} & \textbf{Meaning} \\
\midrule
\texttt{freq\_cm1}                  & cm$^{-1}$ & mode frequency \\
\texttt{dr\_FeFe\_rss\_a}           & $\angstrom$ & RSS aggregate of the Fe-Fe modulation \\
\texttt{dr\_FeFe\_sum\_signed\_a}   & $\angstrom$ & signed-sum aggregate (can be negative) \\
\texttt{lambda\_FeFe\_pair\_cm1}    & cm$^{-1}$ & reorg per bond pair (sub-breakdown) \\
\texttt{lambda\_FeFe\_mode\_cm1}    & cm$^{-1}$ & reorg per mode (aggregate over all
                                                 pairs of the channel) \\
\midrule
\multicolumn{3}{l}{\textit{Repeated for channels: FeN, FeS, NH, HA
       (4 columns each)}} \\
\bottomrule
\end{tabular}
\end{center}

\textbf{Convention}: \texttt{rss} (Root-Sum-Squared) is always
positive and makes Marcus-Hush additivity exact
(\S\ref{app:pseudo-reorg}). \texttt{sum\_signed} is the signed sum
over all bond modulations of the channel; it shows whether the
bond modulations are coherent (Sum $\approx$ rss) or incoherent
(Sum $\ll$ rss).

\textbf{Mode-vs.-pair convention}: \texttt{lambda\_X\_pair} =
individual bond-pair contributions summed; \texttt{lambda\_X\_mode}
= aggregated per mode with $\mu_i\omega_i^2$ weighting. Both are
preserved for different analysis purposes.

\subsubsection*{\texttt{Reorg\_total} -- system aggregates per channel}

Five rows (one per channel: FeFe, FeN, FeS, NH, HA). Columns:

\begin{itemize}\setlength{\itemsep}{0pt}
  \item \texttt{Channel}: name of the bonding channel.
  \item \texttt{Lambda\_total\_pair\_cm1}: $\Lambda_X^\text{total}$
        from pair aggregation.
  \item \texttt{Lambda\_total\_mode\_cm1}: $\Lambda_X^\text{total}$
        from mode aggregation.
  \item \texttt{n\_modes\_contributing}: number of modes with
        non-trivial contribution
        ($\lambda_X(i) > 10^{-6}$\,cm$^{-1}$).
\end{itemize}

These are the central \emph{publication-ready values} of the
reorganization pipeline.

\subsubsection*{\texttt{Reorg\_per\_bond} -- breakdown per individual bond}

One row per individual bond (e.g., Fe1-Cys72-S, Fe1-Fe2,
Fe1-S$_b$1, etc.). Six columns:

\begin{itemize}\setlength{\itemsep}{0pt}
  \item \texttt{Bond}: bond label.
  \item \texttt{Parent\_channel}: which channel
        (FeFe, FeS, FeN, ...) the bond belongs to.
  \item \texttt{Acceptor\_weight}: for the HA channel, the
        Gaussian weight of the H-acceptor pair (default 1.0 for
        the main acceptor).
  \item \texttt{Lambda\_total\_pair\_cm1},
        \texttt{Lambda\_total\_mode\_cm1}: as in
        \texttt{Reorg\_total}, but per bond.
  \item \texttt{n\_modes\_contributing}: modes with contribution
        $> 10^{-6}$\,cm$^{-1}$.
\end{itemize}

Allows the question to be answered: \emph{which individual Fe-S
bond contributes most to $\Lambda_\text{FeS}$?}

\subsubsection*{\texttt{Modulation\_spectra}}

Frequency-resolved modulation spectra $M_X(\omega)$, plus
co-spectra for PCET analysis. One row per frequency grid point
(default: 0.5\,cm$^{-1}$ step). 8 columns:

\begin{itemize}\setlength{\itemsep}{0pt}
  \item \texttt{freq\_cm1}: frequency grid point.
  \item \texttt{M\_FeFe, M\_FeN, M\_FeS, M\_NH, M\_HA}:
        modulation spectra
        $M_X(\omega) = \sum_i |\Delta r_X(i)| \cdot
        \mathcal{N}(\omega - \omega_i, \sigma)$ with default
        $\sigma=5$\,cm$^{-1}$.
  \item \texttt{C\_ET\_FeS}: co-spectrum
        $\sqrt{M_\text{FeFe}(\omega) \cdot M_\text{FeS}(\omega)}$
        -- marks frequency regions where FeFe and FeS are
        modulated simultaneously. High-indicator for
        ET-relevant cluster-breathing modes.
  \item \texttt{C\_PCET}: co-spectrum
        $\sqrt{M_\text{NH}(\omega) \cdot M_\text{HA}(\omega)}$
        -- marks PCET-relevant modes (simultaneous N-H and
        H-acceptor modulation).
\end{itemize}

\subsubsection*{\texttt{Lambda\_cumulative}}

Cumulative reorganization curves $\Lambda_X(\omega) = \sum_{\omega_i \leq \omega}
\lambda_X(i)$. One row per frequency grid point. 6 columns:
\texttt{freq\_cm1}, \texttt{Lambda\_FeFe}, \texttt{Lambda\_FeN},
\texttt{Lambda\_FeS}, \texttt{Lambda\_NH}, \texttt{Lambda\_HA}.

Each column is a \textbf{monotone step function}: jumps mark
modes with $\lambda_X$ contribution. Directly overlayable on an
NRVS spectrum for band identification.

\subsubsection*{\texttt{B\_factors}}

Per atom: thermal $B$ factors $B_a = \frac{8\pi^2}{3} \langle u_a^2
\rangle$ from the QHO averaging over all modes. Columns: atom
index, element, $B$ factor in $\angstrom^2$.

\subsubsection*{\texttt{Info}}

Per-run metadata: log-file path, number of modes, frequency
range, cluster geometry, ligand list, configuration parameters.
Full-text description of the run in a single sheet.

\subsection{Embeddings workbook (\texttt{*\_analysis\_Embeddings.xlsx})}
\label{app:wb-embeddings}

Contains the UMAP projections and HDBSCAN clusterings. The workbook
hosts \emph{three} parallel embeddings of the same set of modes,
each in its own pair of \texttt{*\_clusters} / \texttt{*\_profile}
sheets, plus the per-mode raw C$_\alpha$ amplitude matrix used as
input for the Ca-UMAP.

\subsubsection*{\texttt{Projections}}

Global UMAP. One row per mode. Columns: \texttt{Mode \#},
\texttt{Frequency}, \texttt{UMAP\_x}, \texttt{UMAP\_y}. Computed
on the full feature matrix built by
\texttt{embedding.build\_feature\_matrix} (Marcus-Hush reorganization
features + bond modulation features + kernel scores).
Directly plottable as a 2D embedding.

\subsubsection*{\texttt{Ca\_amplitudes}}

Per mode and per C$_\alpha$ residue: the thermally scaled
$\lVert\evec\rVert$ amplitude of the backbone $\alpha$-carbon for
the residue.  Used as the feature matrix for the Ca-UMAP
(below).  Format: rows = modes, columns = residue numbers.

\subsubsection*{\texttt{Cluster} / \texttt{Cluster\_Profile}}
\label{app:wb-embed-cluster}

Global UMAP cluster assignments. \texttt{Cluster} has one row per
mode with its HDBSCAN cluster ID (\texttt{cluster\_label}: $-1$ =
noise; $0, 1, \ldots$ = cluster index), distance to cluster centre,
and all feature values.
\texttt{Cluster\_Profile} has one row per HDBSCAN cluster with
the mean feature vector, Fisher-F statistic per feature, and a
ranked list of representative modes (closest-to-centroid).

\subsubsection*{\texttt{SSE\_UMAP\_clusters} / \texttt{SSE\_UMAP\_profile}
                  \emph{(since v1.0.3)}}
\label{app:wb-embed-ssumap}

Secondary-structure UMAP, restricted to modes that have at least one
non-zero entry in the per-SSE-element feature row (\texttt{sse}-dict
non-empty after \texttt{analyze\_all\_sse}). Useful for identifying
collective protein-skeleton modes that visit multiple secondary
structures simultaneously. The \texttt{*\_clusters} sheet lists
modes with their SSE-UMAP coordinates and HDBSCAN cluster ID; the
\texttt{*\_profile} sheet gives the Z-score profile of the top-30
discriminating SSE-elements per cluster.

\subsubsection*{\texttt{Ca\_UMAP\_clusters} / \texttt{Ca\_UMAP\_profile}
                  \emph{(since v1.0.3)}}
\label{app:wb-embed-caumap}

C$_\alpha$-amplitude UMAP. Each mode is represented as a point in an
$N_{C_\alpha}$-dimensional space, one dimension per backbone $\alpha$-carbon
(the columns of \texttt{Ca\_amplitudes}). Features are z-scored
per residue before UMAP, so that no single high-amplitude residue
dominates the embedding distance.

Conceptually, the three UMAPs answer complementary questions:
\begin{itemize}\setlength{\itemsep}{0pt}
  \item \emph{Global UMAP} (\texttt{Projections}): which modes
        couple to the same reorganization channels (Fe-X, NH, HA)?
        Useful for identifying ET-active and PCET-active modes
        respectively.
  \item \emph{SSE-UMAP} (\texttt{SSE\_UMAP\_*}): which modes share
        the same secondary-structure footprint? Useful for
        identifying domain-collective modes.
  \item \emph{Ca-UMAP} (\texttt{Ca\_UMAP\_*}): which modes have
        similar spatial distributions of backbone motion? Useful
        for identifying long-range allosteric pathways and modes
        localised to specific regions of the protein.
\end{itemize}

Companion PNGs (only if \texttt{cfg.export\_embedding\_plots = true}):
\texttt{\_embedding\_UMAP.png} (2 panels),
\texttt{\_embedding\_SSE\_UMAP.png} (3 panels: frequency / mode type /
cluster), \texttt{\_embedding\_Ca\_UMAP.png} (3 panels), and
\texttt{\_ca\_amplitudes\_heatmap.png} (log-scale residue $\times$
frequency heatmap of the C$_\alpha$ amplitudes).

\subsection{Interpolated pDOS workbook (\texttt{*\_analysis\_interp0.05.xlsx})}
\label{app:wb-interp}

A single sheet \texttt{Interpolated} with the pDOS analysis on a
fine frequency grid (default 0.05\,cm$^{-1}$). Columns:
\texttt{freq\_cm1}, \texttt{pDOS\_Fe1}, \texttt{pDOS\_Fe2},
\texttt{pDOS\_S1}, \texttt{pDOS\_S2}, \texttt{pDOS\_total}. Used
for direct comparison with NRVS spectra.

\subsection{Secondary-structure workbook
            (\texttt{*\_analysis\_SSE.xlsx}, optional)}
\label{app:wb-sse}

Created only when a PDB file with (or DSSP-extracted)
secondary-structure records is provided. Typically contains
9--11 sheets, depending on SSE complexity:

\subsubsection*{Per-SSE-element sheets (\texttt{SSE\_*})}

One sheet per SSE metric:
\texttt{SSE\_amplitude\_mean},
\texttt{SSE\_amplitude\_max},
\texttt{SSE\_com\_amplitude} (center-of-mass amplitude),
\texttt{SSE\_lateral\_std} (lateral motion scatter),
\texttt{SSE\_lateral\_amplitude},
\texttt{SSE\_stretching} (axial length change),
\texttt{SSE\_axial\_amplitude},
\texttt{SSE\_tilting\_angle} (helix tilt),
\texttt{SSE\_internal\_amplitude} (internal motion).

Each sheet has one row per mode and one column per SSE element.

\subsubsection*{\texttt{SSE\_UMAP\_Cluster} and \texttt{SSE\_UMAP\_Profile}}

Like \texttt{UMAP\_Cluster} and \texttt{UMAP\_Cluster\_Profile},
but on the SSE feature matrix (instead of the cluster-ligand
matrix). Identifies structural clusters of modes based on their
effect on secondary-structure elements.

\begin{infobox}
\textbf{Which sheets are required for a publication?}
For a reorganization manuscript table, \texttt{Reorg\_total} and
\texttt{Reorg\_per\_bond} (main analysis) are typically
sufficient. For band assignment, \texttt{Lambda\_cumulative} is
added (NRVS overlay). For mode characterization, the sheet
\texttt{Mode\_analysis} is needed and (for symmetry-relevant
questions) \texttt{SCSD}. The rest serves internal diagnostic
purposes.
\end{infobox}


% =========================================================================
% Bibliography
% =========================================================================

\begin{thebibliography}{99}

\bibitem{Kingsbury2024}
C.~J. Kingsbury, M.~O. Senge,
\emph{Quantifying near-symmetric molecular distortion using
symmetry-coordinate structural decomposition},
\textbf{Chem.\ Sci.} \textbf{15}, 13638--13649 (2024).
DOI: 10.1039/D4SC01670J. Web implementation:
\url{https://www.kingsbury.id.au/scsd}.

\bibitem{KingsburyPreprint2024}
C.~J. Kingsbury, M.~O. Senge,
\emph{Quantifying near-symmetric molecular distortion using
symmetry-coordinate structural decomposition},
ChemRxiv preprint, doi:10.26434/chemrxiv-2024-6b25s (2024).

\bibitem{Iwata1996}
S.~Iwata, M.~Saynovits, T.~A. Link, H.~Michel,
\emph{Structure of a water soluble fragment of the {`Rieske'} iron-sulfur
protein of the bovine heart mitochondrial cytochrome bc$_1$ complex
determined by {MAD} phasing at 1.5 \AA{} resolution},
\textbf{Structure} \textbf{4}, 567--579 (1996).

\bibitem{SauerHeyden2023}
P.~Sauer, M.~Heyden,
\emph{Frequency-selective anharmonic mode analysis of thermally excited
vibrations in proteins},
\textbf{J.\ Chem.\ Phys.} (2023).

\bibitem{Achterhold2002}
K.~Achterhold, F.~G. Parak,
\emph{Protein dynamics: determination of anisotropic vibrations at the
heme iron of myoglobin},
\textbf{J.\ Phys.: Condens.\ Matter} \textbf{14}, S1399--S1421 (2002).

\bibitem{Paulsen1999}
H.~Paulsen, H.~Winkler, A.~X. Trautwein \emph{et al.},
\emph{Measurement and simulation of nuclear inelastic-scattering spectra
of molecular crystals},
\textbf{Phys.\ Rev.\ B} \textbf{59}, 975 (1999).

\bibitem{Gee2021}
L.~B. Gee, V.~Pelmenschikov, C.~Mons, N.~Mishra, H.~Wang, Y.~Yoda,
K.~Tamasaku, M.-P. Golinelli-Cohen, S.~P. Cramer,
\emph{NRVS and DFT of mitoNEET: Understanding the special vibrational
structure of a [2Fe--2S] cluster with (Cys)$_3$(His)$_1$ ligation},
\textbf{Biochemistry} \textbf{60}, 2419--2424 (2021).
DOI: 10.1021/acs.biochem.1c00252.

\bibitem{Conlan2011}
A.~R. Conlan, M.~L. Paddock, C.~Homer \emph{et al.},
\emph{Mutation of the His ligand in mitoNEET stabilizes the 2Fe--2S
cluster despite conformational heterogeneity in the ligand environment},
\textbf{Acta Cryst.\ F} \textbf{67}, 516--523 (2011).

\bibitem{GoNoguti1983}
N.~Go, T.~Noguti, T.~Nishikawa,
\emph{Dynamics of a small globular protein in terms of low-frequency
vibrational modes},
\textbf{Proc.\ Natl.\ Acad.\ Sci.\ USA} \textbf{80}, 3696 (1983).

\bibitem{ParakKnapp1984}
F.~Parak, E.~W. Knapp,
\emph{A consistent picture of protein dynamics},
\textbf{Proc.\ Natl.\ Acad.\ Sci.\ USA} \textbf{81}, 7088 (1984).

\bibitem{WilsonDeciusCross1955}
E.~B. Wilson, J.~C. Decius, P.~C. Cross,
\emph{Molecular Vibrations: The Theory of Infrared and Raman
Vibrational Spectra},
McGraw-Hill, New York (1955); Dover Reprint (1980).

\bibitem{Sturhahn2004}
W.~Sturhahn,
\emph{Nuclear resonant spectroscopy},
\textbf{J.\ Phys.: Condens.\ Matter} \textbf{16}, S497--S530 (2004).

\bibitem{Sturhahn2000}
W.~Sturhahn, A.~Chumakov,
\emph{Lamb-M\"ossbauer factor and second-order Doppler shift from
inelastic nuclear resonant absorption},
\textbf{Hyperfine Interact.} \textbf{123/124}, 809--824 (2000).

\bibitem{Lipkin1995}
H.~J. Lipkin,
\emph{M\"ossbauer sum rules for use with synchrotron sources},
\textbf{Phys.\ Rev.\ B} \textbf{52}, 10073--10079 (1995).

\bibitem{SingwiSjolander1960}
K.~S. Singwi, A.~Sj\"olander,
\emph{Resonance absorption of nuclear gamma rays and the dynamics of
atomic motions},
\textbf{Phys.\ Rev.} \textbf{120}, 1093--1102 (1960).

\bibitem{KohnRuocco1998}
V.~G. Kohn, A.~I. Chumakov, R.~R\"uffer,
\emph{Nuclear resonant inelastic absorption of synchrotron radiation
in an anisotropic single crystal},
\textbf{Phys.\ Rev.\ B} \textbf{58}, 8437--8444 (1998).

\bibitem{Kabsch1976}
W.~Kabsch,
\emph{A solution for the best rotation to relate two sets of vectors},
\textbf{Acta Cryst.\ A} \textbf{32}, 922--923 (1976).

\bibitem{HammesSchifferStuchebrukhov2010}
S.~Hammes-Schiffer, A.~A. Stuchebrukhov,
\emph{Theory of coupled electron and proton transfer reactions},
\textbf{Chem.\ Rev.} \textbf{110}, 6939--6960 (2010).

\bibitem{Marcus1956}
R.~A. Marcus,
\emph{On the theory of oxidation-reduction reactions involving
electron transfer. I},
\textbf{J.\ Chem.\ Phys.} \textbf{24}, 966--978 (1956).
DOI: 10.1063/1.1742723.

\bibitem{Marcus1993}
R.~A. Marcus,
\emph{Electron transfer reactions in chemistry. Theory and experiment}
(Nobel Lecture),
\textbf{Angew.\ Chem.\ Int.\ Ed.\ Engl.} \textbf{32}, 1111--1121 (1993).
DOI: 10.1002/anie.199311113.

\bibitem{Marcus1985}
R.~A. Marcus, N.~Sutin,
\emph{Electron transfers in chemistry and biology},
\textbf{Biochim.\ Biophys.\ Acta} \textbf{811}, 265--322 (1985).

\bibitem{Sturhahn1999}
W.~Sturhahn, V.~G.~Kohn,
\emph{Theoretical aspects of incoherent nuclear resonant scattering},
\textbf{Hyperfine Interact.} \textbf{123/124}, 367--377 (1999).

\bibitem{Chumakov1999}
A.~I.~Chumakov, W.~Sturhahn,
\emph{Experimental aspects of inelastic nuclear resonant scattering},
\textbf{Hyperfine Interact.} \textbf{123/124}, 781--808 (1999).

\bibitem{MaradudinFlinn1963}
A.~A.~Maradudin, P.~A.~Flinn,
\emph{Anharmonic contributions to the Debye-Waller factor},
\textbf{Phys.\ Rev.} \textbf{129}, 2529 (1963).

\bibitem{Lipkin1960}
H.~J.~Lipkin,
\emph{Some simple features of the M\"o\ss bauer effect},
\textbf{Ann.\ Phys.} (N.Y.) \textbf{9}, 332 (1960).

\bibitem{Hu2003}
M.~Y.~Hu, W.~Sturhahn, T.~S.~Toellner \emph{et al.},
\emph{Measuring velocity of sound with nuclear resonant inelastic
X-ray scattering},
\textbf{Phys.\ Rev.\ B} \textbf{67}, 094304 (2003).

\bibitem{Leu2008}
B.~M.~Leu, M.~Z.~Hu, J.~Zhao, S.~P.~Cramer,
\emph{Quantitative vibrational dynamics of iron in carbonyl compounds},
\textbf{J.\ Am.\ Soc.\ Mass Spectrom.} \textbf{17}, 235 (2008).

\bibitem{ThompsonCoxHastings1987}
P.~Thompson, D.~E.~Cox, J.~B.~Hastings,
\emph{Rietveld refinement of Debye-Scherrer synchrotron X-ray data
from $\mathrm{Al}_2\mathrm{O}_3$},
\textbf{J.\ Appl.\ Crystallogr.} \textbf{20}, 79 (1987).

\bibitem{ina32}
K.~Achterhold, F.~G.~Parak,
\emph{ina32.f -- Fortran source code for analysis of nuclear inelastic
scattering spectra}; internal publication, TU Munich, Department of
Biophysics (2005--).

\bibitem{Codata2018}
E.~Tiesinga, P.~J.~Mohr, D.~B.~Newell, B.~N.~Taylor,
\emph{CODATA recommended values of the fundamental physical constants:
2018},
\textbf{Rev.\ Mod.\ Phys.} \textbf{93}, 025010 (2021).

% --- Embedding theory ----------------------------------------------------

\bibitem{Pearson1901}
K.~Pearson,
\emph{On lines and planes of closest fit to systems of points in space},
\textbf{Philos.\ Mag.} \textbf{2}, 559--572 (1901).

\bibitem{Hotelling1933}
H.~Hotelling,
\emph{Analysis of a complex of statistical variables into principal
components},
\textbf{J.\ Educ.\ Psychol.} \textbf{24}, 417--441 (1933).

\bibitem{Jolliffe2002}
I.~T. Jolliffe,
\emph{Principal Component Analysis},
2nd ed., Springer (2002).

\bibitem{vanDerMaaten2008}
L.~van der Maaten, G.~Hinton,
\emph{Visualizing data using t-SNE},
\textbf{J.\ Mach.\ Learn.\ Res.} \textbf{9}, 2579--2605 (2008).

\bibitem{vanDerMaaten2014}
L.~van der Maaten,
\emph{Accelerating t-SNE using tree-based algorithms},
\textbf{J.\ Mach.\ Learn.\ Res.} \textbf{15}, 3221--3245 (2014).

\bibitem{McInnes2018}
L.~McInnes, J.~Healy, J.~Melville,
\emph{UMAP: Uniform Manifold Approximation and Projection for
Dimension Reduction},
\textbf{arXiv} 1802.03426 (2018).

\bibitem{Becht2019}
E.~Becht, L.~McInnes, J.~Healy \emph{et al.},
\emph{Dimensionality reduction for visualizing single-cell data using
UMAP},
\textbf{Nat.\ Biotechnol.} \textbf{37}, 38--44 (2019).

\bibitem{Campello2013}
R.~J.~G.~B. Campello, D.~Moulavi, J.~Sander,
\emph{Density-based clustering based on hierarchical density estimates},
in \textbf{Adv.\ Knowl.\ Discov.\ Data Min.} (PAKDD 2013), pp.~160--172.

\bibitem{McInnes2017}
L.~McInnes, J.~Healy, S.~Astels,
\emph{hdbscan: Hierarchical density based clustering},
\textbf{J.\ Open Source Softw.} \textbf{2}, 205 (2017).

% --- Extended PCET/ET theory --------------------------------------------

\bibitem{HammesSchiffer2010}
S.~Hammes-Schiffer,
\emph{Theory of proton-coupled electron transfer in energy conversion
processes},
\textbf{Acc.\ Chem.\ Res.} \textbf{42}, 1881--1889 (2009).
DOI: 10.1021/ar9001284.

\bibitem{EdwardsSoudackov2010}
S.~J. Edwards, A.~V. Soudackov, S.~Hammes-Schiffer,
\emph{Driving force dependence of rates for nonadiabatic proton and
proton-coupled electron transfer: conditions for inverted region behavior},
\textbf{J.\ Phys.\ Chem.\ B} \textbf{113}, 14545--14548 (2009).

\bibitem{LayfieldHammesSchiffer2014}
J.~P. Layfield, S.~Hammes-Schiffer,
\emph{Hydrogen tunneling in enzymes and biomimetic models},
\textbf{Chem.\ Rev.} \textbf{114}, 3466--3494 (2014).

\bibitem{CuiKarplus2008}
Q.~Cui, M.~Karplus,
\emph{Allostery and cooperativity revisited},
\textbf{Protein Sci.} \textbf{17}, 1295--1307 (2008).

\bibitem{Bergner2017}
M.~Bergner, S.~Dechert, S.~Demeshko, C.~Kupper, J.~M.~Mayer, F.~Meyer,
\emph{Model of the {MitoNEET} {[}2Fe-2S{]} cluster shows proton coupled
electron transfer},
\textbf{J.\ Am.\ Chem.\ Soc.} \textbf{139}, 701--707 (2017).
DOI: 10.1021/jacs.6b09180.

\bibitem{Kuila1992}
D.~Kuila, J.~R. Schoonover, R.~B. Dyer, C.~J. Batie, D.~P. Ballou,
J.~A. Fee, W.~H. Woodruff,
\emph{Resonance {R}aman studies of {R}ieske-type proteins},
\textbf{Biochim.\ Biophys.\ Acta} \textbf{1140}, 175--183 (1992).

\bibitem{Wang2025}
H.~Wang, V.~Pelmenschikov, Y.~Yoda, S.~P. Cramer,
\emph{NRVS of {Fe-S} cluster proteins and models -- A bestiary of
nifty normal modes},
\textbf{J.\ Inorg.\ Biochem.} \textbf{270}, 112935 (2025).
DOI: 10.1016/j.jinorgbio.2025.112935.

\bibitem{Zuris2011}
J.~A. Zuris, D.~A. Halim, A.~R. Conlan \emph{et al.},
\emph{Engineering the redox potential over a wide range within a
new class of FeS proteins},
\textbf{J.\ Am.\ Chem.\ Soc.} \textbf{132}, 13120--13122 (2010).

\bibitem{KurodaHammesSchiffer2020}
T.~Kuroda, S.~Hammes-Schiffer,
\emph{Vibronic coupling for proton-coupled electron transfer reactions:
recent insights from theory},
\textbf{Curr.\ Opin.\ Electrochem.} \textbf{29}, 100789 (2021).

% --- PCET/ET literature, iron-sulfur spectroscopy -----------------------

\bibitem{HammesSchifferSoudackov2008}
S.~Hammes-Schiffer, A.~V.~Soudackov,
\emph{Proton-coupled electron transfer in solution, proteins, and
electrochemistry},
\textbf{J.\ Phys.\ Chem.\ B} \textbf{112}, 14108--14123 (2008).

\bibitem{HuangRhys1950}
K.~Huang, A.~Rhys,
\emph{Theory of light absorption and non-radiative transitions in
$F$-centres},
\textbf{Proc.\ R.\ Soc.\ Lond.\ A} \textbf{204}, 406--423 (1950).

\bibitem{Beinert1997}
H.~Beinert, R.~H.~Holm, E.~M\"unck,
\emph{Iron-sulfur clusters: nature's modular, multipurpose structures},
\textbf{Science} \textbf{277}, 653--659 (1997).

\bibitem{Han1989}
S.~Han, R.~S.~Czernuszewicz, T.~G.~Spiro,
\emph{Vibrational spectra and normal mode analysis for {[}2Fe-2S{]}
protein analogs using sulfur-34, iron-54 and deuterium substitution:
coupling of iron-sulfur stretching and sulfur-carbon-carbon bending modes},
\textbf{J.\ Am.\ Chem.\ Soc.} \textbf{111}, 3496--3504 (1989).

\bibitem{Czernuszewicz1986}
R.~S.~Czernuszewicz, J.~LeGall, I.~Moura, T.~G.~Spiro,
\emph{Resonance Raman spectra of rubredoxin: new assignments and
vibrational coupling mechanism from iron-54/iron-56 isotope shifts
and variable-wavelength excitation},
\textbf{Inorg.\ Chem.} \textbf{25}, 696--700 (1986).

\bibitem{Bergmann2003}
U.~Bergmann, W.~Sturhahn, D.~E.~Linn, F.~E.~Jenney, M.~W.~W.~Adams,
K.~Rupnik, B.~J.~Hales, E.~E.~Alp, A.~Mayse, S.~P.~Cramer,
\emph{Observation of {Fe-H/D} modes by nuclear resonant vibrational
spectroscopy},
\textbf{J.\ Am.\ Chem.\ Soc.} \textbf{125}, 4016--4017 (2003).

\bibitem{Saouma2014}
C.~T.~Saouma, M.~M.~Pinney, J.~M.~Mayer,
\emph{Electron transfer and proton-coupled electron transfer reactivity
and self-exchange of synthetic {[}2Fe-2S{]} complexes:
Models for Rieske and mitoNEET clusters},
\textbf{Inorg.\ Chem.} \textbf{53}, 3153--3161 (2014).

\bibitem{Albers2014Rieske}
A.~Albers, S.~Demeshko, S.~Dechert, C.~T.~Saouma, J.~M.~Mayer, F.~Meyer,
\emph{Fast proton-coupled electron transfer observed for a high-fidelity
structural and functional {[}2Fe-2S{]} {Rieske} model},
\textbf{J.\ Am.\ Chem.\ Soc.} \textbf{136}, 3946--3954 (2014).

% --- Methodological sources for vibrational analysis --------------------

\bibitem{Cyvin1968}
S.~J.~Cyvin,
\emph{Molecular Vibrations and Mean Square Amplitudes},
Universitetsforlaget, Oslo / Elsevier, Amsterdam (1968).

\bibitem{Ochterski1999}
J.~W.~Ochterski,
\emph{Vibrational Analysis in Gaussian},
Gaussian, Inc., White Paper (1999).
\url{https://gaussian.com/vib/}

\bibitem{Frisch2016Gaussian16}
M.~J.~Frisch \emph{et al.},
\emph{Gaussian 16, Revision C.01},
Gaussian, Inc., Wallingford CT (2016).

\bibitem{NeeseORCAmanual2023}
F.~Neese,
\emph{The ORCA program system: user's manual},
Version 6.0, Max-Planck-Institut f\"ur Kohlenforschung,
M\"ulheim a.\,d.\,Ruhr (2023).
\url{https://www.faccts.de/docs/orca/6.0/manual/}

\bibitem{Neese2022WIRES}
F.~Neese,
\emph{Software update: The ORCA program system---Version 5.0},
\textbf{WIREs Comput.\ Mol.\ Sci.} \textbf{12}, e1606 (2022).

\end{thebibliography}

\end{document}
